\chapter{A Taste of Modular Forms}

\begin{center}\itshape
There are five elementary arithmetical operations: addition, subtraction, multiplication, division, and modular forms.\\
\upshape --- attributed to Martin Eichler
\end{center}

In 1916 Srinivasa Ramanujan, two years after arriving in Cambridge, published a paper on sums of divisors. In it he expanded the infinite product
\[
\begin{aligned}q\prod_{n\ge1}(1-q^n)^{24}&=\sum_{n\ge1}\tau(n)q^n\\&=q-24q^2+252q^3-1472q^4\eqbr\quad+4830q^5-\cdots\end{aligned}
\]
and looked at the coefficients. They seemed to have no reason to be anything in particular. Ramanujan noticed three things. First, $\tau(6)=-6048=\tau(2)\tau(3)$, and more generally $\tau(mn)=\tau(m)\tau(n)$ whenever
$\gcd(m,n)=1$. Second, $\tau(n)\equiv\sigma_{11}(n)\pmod{691}$, where $\sigma_{11}(n)$ is the sum of the eleventh powers of the divisors of $n$. Third, $|\tau(p)|\le2p^{11/2}$ for every prime $p$. He could prove none of them.

Mordell proved the first a year later. The second follows from a single identity between two functions, and we shall prove it in a few lines. The third resisted for fifty-eight years: Deligne proved it in 1974,
as a consequence of his proof of the Weil conjectures, one of the deepest results of twentieth-century mathematics. All three are statements about \emph{modular forms}, functions on the upper half-plane with an
enormous group of symmetries. That group forces the space of such functions to be finite-dimensional, and finite-dimensional spaces are where coincidences become theorems.

\begin{think}
Two power series built from divisor sums, $1+240\sum\sigma_3(n)q^n$ and $1+480\sum\sigma_7(n)q^n$, turn out to satisfy $(1+240\sum\sigma_3(n)q^n)^2=1+480\sum\sigma_7(n)q^n$. Compare the coefficients of $q^2$ by hand.
What identity between divisor sums does this give, and why would anyone expect it to hold for every $n$?
\end{think}

This last chapter is a taste, not a meal. We prove what is needed to explain Ramanujan's first two observations, with full proofs, and use only complex analysis from Volume~4. Volume~12 develops the theory properly.

\section{The modular group}

Let $\mathbb H=\{\tau\in\C:\operatorname{Im}\tau>0\}$ be the upper half-plane. The group $\mathrm{SL}_2(\Z)$ of integer matrices of determinant $1$ acts on $\mathbb H$ by Möbius transformations:
\[
\gamma=\begin{pmatrix}a&b\\c&d\end{pmatrix},\qq\gamma\tau=\frac{a\tau+b}{c\tau+d},\qq\operatorname{Im}(\gamma\tau)=\frac{\operatorname{Im}\tau}{|c\tau+d|^2}.
\]
The last formula follows from $(a\tau+b)(c\overline\tau+d)=ac|\tau|^2+bd+ad\tau+bc\overline\tau$, whose imaginary part is $(ad-bc)\operatorname{Im}\tau$. So $\gamma\tau\in\mathbb H$, and one checks that $(\gamma\gamma')\tau=\gamma(\gamma'\tau)$
by multiplying out. The matrices $\gamma$ and $-\gamma$ act in the same way. Two elements play a special role:
\[
\begin{gathered}T=\begin{pmatrix}1&1\\0&1\end{pmatrix},\qquad T\tau=\tau+1,\\S=\begin{pmatrix}0&-1\\1&0\end{pmatrix},\qquad S\tau=-\frac1\tau.\end{gathered}
\]

\begin{proposition}\label{prop:generators}
The matrices $S$ and $T$ generate $\mathrm{SL}_2(\Z)$.
\end{proposition}

\begin{proof}
Let $G$ be the subgroup generated by $S$ and $T$; it contains $S^2=-I$. Let $\gamma=\bigl(\begin{smallmatrix}a&b\\c&d\end{smallmatrix}\bigr)\in\mathrm{SL}_2(\Z)$; we show $\gamma\in G$ by induction on $|c|$. If $c=0$, then $ad=1$, so $\gamma=\pm\bigl(\begin{smallmatrix}1&b'\\0&1\end{smallmatrix}\bigr)=\pm T^{b'}$.
If $c\ne0$, Euclid's division gives $d=kc+r$ with $0\le r<|c|$, and
\[
\gamma T^{-k}=\begin{pmatrix}a&b-ka\\c&r\end{pmatrix},\qq\gamma T^{-k}S=\begin{pmatrix}b-ka&-a\\r&-c\end{pmatrix}.
\]
The lower-left entry $r$ satisfies $|r|<|c|$, so $\gamma T^{-k}S\in G$ by induction, and so is $\gamma$.
\end{proof}

\begin{theorem}\label{thm:fundamental}
Let $F=\{\tau\in\mathbb H:|\operatorname{Re}\tau|\le\frac12,\ |\tau|\ge1\}$. Every $\tau\in\mathbb H$ is of the form $\gamma\tau'$ with $\gamma\in\mathrm{SL}_2(\Z)$ and $\tau'\in F$.
\end{theorem}

\begin{proof}
Fix $\tau$. For $c,d$ integers, $|c\tau+d|\le1$ forces $|c|\operatorname{Im}\tau\le1$ and then $|d|\le1+|c||\tau|$, so only finitely many pairs $(c,d)$ satisfy it. Hence $\operatorname{Im}(\gamma\tau)=\frac{\operatorname{Im}\tau}{|c\tau+d|^2}$
takes only finitely many values larger than $\operatorname{Im}\tau$, and there is $\gamma$ for which $\operatorname{Im}(\gamma\tau)$ is maximal. Replacing $\gamma$ by $T^m\gamma$, which does not change the imaginary part, we may assume $|\operatorname{Re}\gamma\tau|\le\frac12$.
If $|\gamma\tau|<1$, then $\operatorname{Im}(S\gamma\tau)=\frac{\operatorname{Im}\gamma\tau}{|\gamma\tau|^2}>\operatorname{Im}\gamma\tau$, contradicting maximality. So $\gamma\tau\in F$, and $\tau=\gamma^{-1}(\gamma\tau)$.
\end{proof}

So $F$ meets every orbit, and Volume~12 shows that it meets each orbit only once, apart from identifications on its boundary. The images of $F$ tile $\mathbb H$ (Figure~\ref{fig:tessellation}). Every point of $F$
has imaginary part at least $\frac{\sqrt3}2$, the height of the corner $\rho=e^{2\pi i/3}$.

\begin{figure}[ht]
\centering
\fig{ch12-tessellation}
\caption{The fundamental domain $F$ (shaded) and its images under $\mathrm{SL}_2(\Z)$, which tile the upper half-plane. Each region is bounded by arcs of circles centred on the real axis, or vertical lines. The tiles
accumulate at every rational point of the real axis.}
\label{fig:tessellation}
\end{figure}

\section{Modular forms and Eisenstein series}

For $\gamma=\bigl(\begin{smallmatrix}a&b\\c&d\end{smallmatrix}\bigr)$ put $j(\gamma,\tau)=c\tau+d$. Multiplying out shows the \emph{cocycle relation}
\[
j(\gamma\gamma',\tau)=j(\gamma,\gamma'\tau)\,j(\gamma',\tau).
\]
For an integer $k$ and a function $f$ on $\mathbb H$, define $(f|_k\gamma)(\tau)=j(\gamma,\tau)^{-k}f(\gamma\tau)$. The cocycle relation says exactly that $f|_k(\gamma\gamma')=(f|_k\gamma)|_k\gamma'$.

\begin{definition}\label{def:modform}
A \emph{modular form of weight $k$} is a holomorphic function $f$ on $\mathbb H$ such that $f|_k\gamma=f$ for all $\gamma\in\mathrm{SL}_2(\Z)$, that is
\[
f\Bigl(\frac{a\tau+b}{c\tau+d}\Bigr)=(c\tau+d)^kf(\tau),
\]
and which is \emph{holomorphic at infinity}: there is a power series $h(q)=\sum_{n\ge0}a_nq^n$ converging for $|q|<1$ with $f(\tau)=h(e^{2\pi i\tau})$. The $a_n$ are the \emph{Fourier coefficients} of $f$; if $a_0=0$, $f$ is a
\emph{cusp form}. We write $M_k$ and $S_k$ for the complex vector spaces of modular forms and of cusp forms of weight $k$.
\end{definition}

From now on $q=e^{2\pi i\tau}$; note $|q|=e^{-2\pi\operatorname{Im}\tau}<1$. By Proposition~\ref{prop:generators} and the composition rule, $f|_k\gamma=f$ holds for all $\gamma$ as soon as it holds for $S$ and $T$, that is,
as soon as $f(\tau+1)=f(\tau)$ and $f(-\frac1\tau)=\tau^kf(\tau)$. Taking $\gamma=-I$ gives $f=(-1)^kf$, so $M_k=0$ for odd $k$. Products of modular forms are modular, and weights add: $M_kM_l\subseteq M_{k+l}$.

The first examples are lattice sums. For an even $k\ge4$ and $\tau\in\mathbb H$, the \emph{Eisenstein series} is
\[
G_k(\tau)=\sum_{(m,n)\ne(0,0)}\frac1{(m\tau+n)^k}.
\]

\begin{lemma}\label{lem:latticesum}
For $\sigma>2$ the series $\sum_{(m,n)\ne(0,0)}|m\tau+n|^{-\sigma}$ converges uniformly for $\tau$ in every compact subset of $\mathbb H$.
\end{lemma}

\begin{proof}
Let $K\subseteq\mathbb H$ be compact. The function $(\tau,x,y)\mapsto|x\tau+y|$ is continuous and positive on the compact set of $(\tau,x,y)$ with $\tau\in K$ and $\max(|x|,|y|)=1$ (it vanishes only if $x=y=0$, since $\tau\notin\R$),
so it has a minimum $c>0$ there. By homogeneity, $|m\tau+n|\ge c\max(|m|,|n|)$ for all $(m,n)$. There are $8r$ pairs with $\max(|m|,|n|)=r\ge1$, so the series is at most $\sum_{r\ge1}8r(cr)^{-\sigma}<\infty$, uniformly on $K$.
\end{proof}

\begin{proposition}\label{prop:eisenstein}
For even $k\ge4$, $G_k$ is holomorphic on $\mathbb H$ and $G_k|_k\gamma=G_k$ for all $\gamma\in\mathrm{SL}_2(\Z)$.
\end{proposition}

\begin{proof}
Holomorphy follows from Lemma~\ref{lem:latticesum} and Weierstrass's theorem (Volume~4, Theorem~9.15). For $\gamma=\bigl(\begin{smallmatrix}a&b\\c&d\end{smallmatrix}\bigr)$,
\[
\frac1{(m\gamma\tau+n)^k}=\frac{(c\tau+d)^k}{\bigl((ma+nc)\tau+(mb+nd)\bigr)^k},
\]
and $(m,n)\mapsto(ma+nc,mb+nd)$ is a bijection of $\Z^2\setminus\{0\}$, with inverse given by the inverse matrix. Summing over $(m,n)$, which is allowed by absolute convergence, $G_k(\gamma\tau)=(c\tau+d)^kG_k(\tau)$.
\end{proof}

\begin{theorem}[Fourier expansion of $G_k$]\label{thm:eisensteinq}
For even $k\ge4$,
\[
G_k(\tau)=2\zeta(k)+\frac{2(2\pi i)^k}{(k-1)!}\sum_{n\ge1}\sigma_{k-1}(n)q^n,
\]
where $\sigma_{k-1}(n)=\sum_{d\mid n}d^{k-1}$. Consequently $E_k=\frac{G_k}{2\zeta(k)}$ is a modular form of weight $k$, with
\[
E_k=1-\frac{2k}{B_k}\sum_{n\ge1}\sigma_{k-1}(n)q^n.
\]
\end{theorem}

\begin{proof}
By Volume~4, Corollary~9.35, for $\tau\in\mathbb H$,
\[
\pi\cot\pi\tau=\frac1\tau+\sum_{n\ge1}\Bigl(\frac1{\tau+n}+\frac1{\tau-n}\Bigr),
\]
and on the other hand $\pi\cot\pi\tau=\pi i\frac{q+1}{q-1}=-\pi i\bigl(1+2\sum_{r\ge1}q^r\bigr)$ since $|q|<1$. Both series converge uniformly on compact subsets of $\mathbb H$, so we may differentiate $k-1$ times term by term
(Volume~4, Theorem~9.15). As $\frac{\mathrm d^{k-1}}{\mathrm d\tau^{k-1}}\frac1{\tau+n}=\frac{(-1)^{k-1}(k-1)!}{(\tau+n)^k}$ and $\frac{\mathrm d}{\mathrm d\tau}q^r=2\pi irq^r$, we get, for even $k$, the \emph{Lipschitz formula}
\[
\sum_{n\in\Z}\frac1{(\tau+n)^k}=\frac{(2\pi i)^k}{(k-1)!}\sum_{r\ge1}r^{k-1}q^r.
\]
In $G_k$, the terms with $m=0$ give $2\zeta(k)$. The terms with $m\ne0$ come in pairs $(m,n)$, $(-m,-n)$ with equal values since $k$ is even, so they give $2\sum_{m\ge1}\sum_n(m\tau+n)^{-k}$. Apply the Lipschitz formula
at $m\tau$, whose $q$ is $q^m$:
\[
\begin{aligned}G_k(\tau)-2\zeta(k)&=\frac{2(2\pi i)^k}{(k-1)!}\sum_{m\ge1}\sum_{r\ge1}r^{k-1}q^{mr}\eqbr=\frac{2(2\pi i)^k}{(k-1)!}\sum_{n\ge1}\sigma_{k-1}(n)q^n,\end{aligned}
\]
grouping the terms with $mr=n$ (the double series converges absolutely). Finally, Euler's formula (Volume~4, Theorem~3.7) gives $2\zeta(k)=\frac{(-1)^{k/2+1}B_k(2\pi)^k}{k!}$, and $(2\pi i)^k=(-1)^{k/2}(2\pi)^k$, so
$\frac{2(2\pi i)^k}{(k-1)!\,2\zeta(k)}=-\frac{2k}{B_k}$.
\end{proof}

With $B_4=-\frac1{30}$, $B_6=\frac1{42}$, $B_8=-\frac1{30}$, $B_{10}=\frac5{66}$, $B_{12}=-\frac{691}{2730}$:
\[
\begin{gathered}E_4=1+240\sum\sigma_3(n)q^n,\qq E_6=1-504\sum\sigma_5(n)q^n,\\E_8=1+480\sum\sigma_7(n)q^n,\qq E_{10}=1-264\sum\sigma_9(n)q^n,\\E_{12}=1+\frac{65520}{691}\sum\sigma_{11}(n)q^n.\end{gathered}
\]
There it is: the prime $691$, hiding in the numerator of a Bernoulli number.

\section{Dedekind's eta function}

Ramanujan's product is the $24$-th power of a function that Dedekind studied in 1877:
\[
\begin{gathered}\eta(\tau)=e^{\pi i\tau/12}\prod_{n\ge1}(1-q^n),\\\Delta(\tau)=\eta(\tau)^{24}=q\prod_{n\ge1}(1-q^n)^{24}.\end{gathered}
\]
Since $\sum|q|^n$ converges, uniformly on $\operatorname{Im}\tau\ge y_0>0$, the function $L(\tau)=\frac{\pi i\tau}{12}+\sum_{n\ge1}\log(1-q^n)$ (principal logarithms, as $|q^n|<1$) is holomorphic on $\mathbb H$, and $\eta=e^L$. In particular $\eta$ and
$\Delta$ are holomorphic and \emph{never zero} on $\mathbb H$. Clearly $\Delta(\tau+1)=\Delta(\tau)$. The heart of the matter is the behaviour under $S$.

\begin{theorem}[Dedekind]\label{thm:eta}
For $\tau\in\mathbb H$, $\eta(-\frac1\tau)=\sqrt{-i\tau}\,\eta(\tau)$, with the principal square root. Hence $\Delta(-\frac1\tau)=\tau^{12}\Delta(\tau)$, and $\Delta$ is a cusp form of weight $12$ with $\tau(1)=1$.
\end{theorem}

The proof below is Siegel's (1954): a single contour integral of cotangents.

\begin{proof}
Both $L(-\frac1\tau)$ and $L(\tau)+\frac12\log(-i\tau)$ are holomorphic on $\mathbb H$ ($-i\tau$ lies in the right half-plane), so by the identity theorem (Volume~4, Theorem~9.18) it suffices to prove that they agree for $\tau=iy$, $y>0$.
Expanding $\log(1-x)=-\sum_m\frac{x^m}m$ and summing the geometric series in $n$,
\[
L(iy)=-\frac{\pi y}{12}-\Sigma(y),\qq\Sigma(y)=\sum_{m\ge1}\frac1m\cdot\frac1{e^{2\pi my}-1}.
\]
Since $-\frac1{iy}=\frac iy$, we must show
\[
\Sigma(y)-\Sigma\Bigl(\frac1y\Bigr)+\frac\pi{12}\Bigl(y-\frac1y\Bigr)=\frac12\log y.\tag{$\ast$}
\]
Let $\tau=iy$, $N=n+\frac12$ with $n\ge1$ an integer, and
\[
f_N(z)=\frac1z\cot(\pi Nz)\cot\Bigl(\frac{\pi Nz}\tau\Bigr).
\]
Let $C$ be the parallelogram (here a rhombus) with vertices $1,\tau,-1,-\tau$, traversed positively. Inside $C$, $f_N$ has poles at $z=0$, at $z=\frac kN$ and at $z=\frac{k\tau}N$ for $1\le|k|\le n$; none lies on $C$, since
$N$ is not an integer.

\emph{Residues.} At $z=\frac kN$, $\cot(\pi Nz)$ has residue $\frac1{\pi N}$, so $f_N$ has residue $\frac1{\pi k}\cot\frac{\pi k}\tau$. At $z=\frac{k\tau}N$, similarly, the residue is $\frac1{\pi k}\cot(\pi k\tau)$. At $z=0$, from
$\cot u=\frac1u-\frac u3+O(u^3)$,
\[
\begin{aligned}\cot(\pi Nz)\cot\Bigl(\frac{\pi Nz}\tau\Bigr)&=\frac\tau{\pi^2N^2z^2}\eqbr-\frac13\Bigl(\tau+\frac1\tau\Bigr)+O(z^2),\end{aligned}
\]
so the residue is $-\frac13(\tau+\frac1\tau)$. The terms $k$ and $-k$ give equal residues. Now $\cot(ix)=-i\coth x=-i\bigl(1+\frac2{e^{2x}-1}\bigr)$, so $\cot(\pi k\tau)=-i\bigl(1+\frac2{e^{2\pi ky}-1}\bigr)$ and
$\cot\frac{\pi k}\tau=\cot\frac{-i\pi k}y=i\bigl(1+\frac2{e^{2\pi k/y}-1}\bigr)$; in the sum the constants cancel. With $\tau+\frac1\tau=i(y-\frac1y)$, the residue theorem (Volume~4, Theorem~9.26) gives
\[
\begin{aligned}\frac1{2\pi i}\int_Cf_N&=-\frac i3\Bigl(y-\frac1y\Bigr)\eqbr+\frac{4i}\pi\sum_{k=1}^n\frac1k\Bigl(\frac1{e^{2\pi k/y}-1}-\frac1{e^{2\pi ky}-1}\Bigr).\end{aligned}
\]
Multiplying by $\frac\pi{4i}$: $-\frac18\int_Cf_N=-\frac\pi{12}(y-\frac1y)+\Sigma_n(\frac1y)-\Sigma_n(y)$, where $\Sigma_n$ is the partial sum up to $n$.

\emph{The limit $N\to\infty$.} For $w=u+iv$, $|\cot w|$ is bounded on the set of $w$ at distance at least $\delta$ from $\pi\Z$, for each $\delta>0$, and $\cot w\to\mp i$ as $v\to\pm\infty$. On the edge from $1$ to $\tau$,
$z=1-t+ity$ with $0\le t\le1$, and $\pi Nz=\pi(n+\frac12)-\pi Nt+i\pi Nty$; its distance to the poles $\pi m$ is at least $\min_{s\ge0}|\frac\pi2-s+isy|=\frac\pi2\frac y{\sqrt{1+y^2}}$, independently of $N$. The same holds for $\frac{\pi Nz}\tau$ and on the other
edges, by symmetry. So $z f_N(z)$ is bounded on $C$ uniformly in $N$. At each point of the open edges, $\operatorname{Im}(\pi Nz)$ and $\operatorname{Im}(\frac{\pi Nz}\tau)=-\frac{\pi N\operatorname{Re}z}y$ tend to $\pm\infty$, and $zf_N(z)$ tends to
$+1$ on the edges $[1,\tau]$ and $[-1,-\tau]$, where $\operatorname{Re}z$ and $\operatorname{Im}z$ have the same sign, and to $-1$ on the two others. By bounded convergence,
\[
\begin{aligned}\int_Cf_N\to&\int_{[1,\tau]}\frac{\mathrm dz}z-\int_{[\tau,-1]}\frac{\mathrm dz}z\eqbr+\int_{[-1,-\tau]}\frac{\mathrm dz}z-\int_{[-\tau,1]}\frac{\mathrm dz}z.\end{aligned}
\]
Along each edge, $\int\frac{\mathrm dz}z$ is the change of $\log|z|+i\arg z$, with the argument varying continuously: the four integrals are $\log y+\frac{i\pi}2$, $-\log y+\frac{i\pi}2$, $\log y+\frac{i\pi}2$, $-\log y+\frac{i\pi}2$. So the limit is $4\log y$, and
letting $n\to\infty$ in the identity above gives $-\frac12\log y=-\frac\pi{12}(y-\frac1y)+\Sigma(\frac1y)-\Sigma(y)$, which is $(\ast)$.

Exponentiating, $\eta(-\frac1\tau)=e^{\frac12\log(-i\tau)}\eta(\tau)$, and $\Delta(-\frac1\tau)=(-i\tau)^{12}\Delta(\tau)=\tau^{12}\Delta(\tau)$. With $\Delta(\tau+1)=\Delta(\tau)$, Proposition~\ref{prop:generators} gives $\Delta|_{12}\gamma=\Delta$ for all $\gamma$.
Expanding the product gives a power series $\sum_{n\ge1}\tau(n)q^n$ with integer coefficients and $\tau(1)=1$.
\end{proof}

\section{Finite dimensions}

Now the symmetry does its work. A modular form is determined by its values on $F$, and $F$ is ``compact up to the point at infinity''. This is enough to make weight-$0$ forms constant.

\begin{lemma}\label{lem:weight0}
Every modular form of weight $0$ is constant.
\end{lemma}

\begin{proof}
Let $f\in M_0$, $f(\tau)=h(q)$ with $h$ holomorphic on the disc $|q|<1$. Since $f$ is invariant, Theorem~\ref{thm:fundamental} shows that $f$ takes on $\mathbb H$ the same values as on $F$. Points of $F$ have
$\operatorname{Im}\tau\ge\frac{\sqrt3}2$, hence $|q|\le r_0=e^{-\pi\sqrt3}$. So
\[
\sup_{|q|<1}|h|=\sup_F|f|\le\max_{|q|\le r_0}|h|,
\]
and $|h|$ attains its supremum over the open disc at some point $q_0$ with $|q_0|\le r_0$. By the
maximum principle, $h$ is constant. (Proof of the principle: by Cauchy's formula $h(q_0)$ is the mean of $h$ on each small circle around $q_0$, so if $|h(q_0)|$ is maximal, $|h|$ is constant near $q_0$; then $\overline h=\frac{|h|^2}h$ is
holomorphic there too (or $h=0$), and the Cauchy--Riemann equations, Volume~4, Proposition~9.2, force $h'=0$ near $q_0$; by the identity theorem $h$ is constant.)
\end{proof}

\begin{theorem}\label{thm:dimension}
\begin{enumerate}[label=(\alph*)]
\item $M_k=0$ for $k<0$ and for $k=2$; $M_0=\C$.
\item For even $k\ge4$, $M_k=\C E_k\oplus S_k$, and multiplication by $\Delta$ is an isomorphism $M_{k-12}\to S_k$.
\item Hence $\dim M_k=\lfloor\frac k{12}\rfloor$ if $k\equiv2\pmod{12}$, and $\lfloor\frac k{12}\rfloor+1$ otherwise, for even $k\ge0$.
\end{enumerate}
\end{theorem}

\begin{proof}
\emph{(b) first.} For $f\in M_k$, $f-a_0(f)E_k$ has constant term $0$, so $M_k=\C E_k+S_k$, and the sum is direct since $E_k\notin S_k$. If $f\in S_k$, then $f/\Delta$ is holomorphic on $\mathbb H$ ($\Delta$ has no zeros), satisfies
$(f/\Delta)|_{k-12}\gamma=f/\Delta$, and is holomorphic at infinity: $h_f(q)=q\,g(q)$ and $h_\Delta(q)=q\,u(q)$ with $g,u$ holomorphic on the disc and $u$ without zeros (as $\Delta\ne0$, and $u(0)=1$). So $f/\Delta\in M_{k-12}$, and
$g\mapsto g\Delta$ is the inverse isomorphism.

\emph{(a).} $M_0=\C$ by Lemma~\ref{lem:weight0}. Let $f\in M_k$ with $k<0$. Then $f^{12}\Delta^{-k}$ has weight $12k-12k=0$, so it is constant; and it vanishes at infinity, since $\Delta$ does and $-k>0$. Hence it is $0$ and
$f=0$. So $M_k=0$ for $k<0$, and by (b), $S_k=0$ for $4\le k\le10$; thus $M_4=\C E_4$ and $M_6=\C E_6$. Let $f\in M_2$. A cusp form of weight $2$ would give $f/\Delta\in M_{-10}=0$, so if $f\ne0$ then $c=a_0(f)\ne0$. Now
$f^2\in M_4$ and $f^3\in M_6$, with constant terms $c^2$ and $c^3$, so $f^2=c^2E_4$ and $f^3=c^3E_6$. Then $c^6E_4^3=f^6=c^6E_6^2$, so $E_4^3=E_6^2$. But $E_4^3-E_6^2=(1+720q+\cdots)-(1-1008q+\cdots)=1728q+\cdots\ne0$.
So $M_2=0$.

\emph{(c)} follows from (a) and (b) by induction on $k$, since $\dim M_k=1+\dim M_{k-12}$ for $k\ge4$, with the values $1,0,1,1,1,1$ for $k=0,2,4,6,8,10$.
\end{proof}

The payoff is immediate. $M_8$, $M_{10}$ and $M_{14}$ are one-dimensional, spanned by $E_8$, $E_{10}$ and $E_{14}$, and they also contain $E_4^2$, $E_4E_6$ and $E_4^2E_6$, with constant term $1$. So $E_4^2=E_8$, $E_4E_6=E_{10}$, and
so on. And $S_{12}$ is the line spanned by $\Delta$.

\begin{corollary}\label{cor:sigma7}
For every $n\ge1$,
\[
\sigma_7(n)=\sigma_3(n)+120\sum_{m=1}^{n-1}\sigma_3(m)\sigma_3(n-m).
\]
\end{corollary}

\begin{proof}
Compare the coefficients of $q^n$ in $E_4^2=E_8$:
\[
\begin{aligned}480\sigma_7(n)&=2\cdot240\sigma_3(n)\eqbr+240^2\sum_{m=1}^{n-1}\sigma_3(m)\sigma_3(n-m),\end{aligned}
\]
and divide by $480$.
\end{proof}

Elementary proofs exist, but they are long; here the identity falls out of a dimension count. For $n=2$ it says $129=9+120$.

\begin{corollary}\label{cor:delta}
$1728\Delta=E_4^3-E_6^2$, and Ramanujan's congruence holds: $\tau(n)\equiv\sigma_{11}(n)\pmod{691}$ for every $n\ge1$.
\end{corollary}

\begin{proof}
$E_4^3-E_6^2$ is a cusp form of weight $12$ with $q$-coefficient $720+1008=1728$, so it is $1728\Delta$, as $S_{12}=\C\Delta$ and $\tau(1)=1$. Next, $E_4^3-E_{12}\in S_{12}$ has $q$-coefficient $720-\frac{65520}{691}=\frac{432000}{691}$, so
\[
691E_4^3=691E_{12}+432000\,\Delta.
\]
The coefficient of $q^n$ gives $691a_n=65520\,\sigma_{11}(n)+432000\,\tau(n)$, where $a_n$, the coefficient of $E_4^3$, is an integer. Modulo $691$, $65520\equiv566\equiv-125$ and $432000\equiv125$, so
$125(\tau(n)-\sigma_{11}(n))\equiv0$, and $125$ is invertible modulo the prime $691$.
\end{proof}

\section{Hecke operators}

Why is $\tau$ multiplicative? In 1937 Hecke found the reason: a family of linear operators on each $M_k$ that commute with each other and, on the one-dimensional space $S_{12}$, can only act by scalars. Let $p$ be a prime.
For $f=\sum a_nq^n\in M_k$, define
\[
\begin{aligned}T_pf(\tau)&=p^{k-1}f(p\tau)+\frac1p\sum_{j=0}^{p-1}f\Bigl(\frac{\tau+j}p\Bigr)\eqbr=\sum_{n\ge0}\bigl(a_{pn}+p^{k-1}a_{n/p}\bigr)q^n,\end{aligned}
\]
with $a_{n/p}=0$ if $p\nmid n$. The second expression follows from $\frac1p\sum_je^{2\pi in(\tau+j)/p}=q^{n/p}$ if $p\mid n$ and $0$ otherwise (a sum of $p$-th roots of unity).

To show that $T_pf$ is modular, extend the slash operator to integer matrices $\alpha=\bigl(\begin{smallmatrix}a&b\\c&d\end{smallmatrix}\bigr)$ with $\det\alpha>0$: $(f|_k\alpha)(\tau)=(\det\alpha)^{k-1}(c\tau+d)^{-k}f(\alpha\tau)$. The cocycle relation holds for these
matrices too (it is an identity between polynomials), and $\det$ is multiplicative, so $f|_k(\alpha\beta)=(f|_k\alpha)|_k\beta$. Let $R_p$ be the set of $p+1$ matrices
\[
\begin{pmatrix}p&0\\0&1\end{pmatrix},\qq\begin{pmatrix}1&j\\0&p\end{pmatrix}\ (0\le j<p).
\]
Then $T_pf=\sum_{\rho\in R_p}f|_k\rho$, as one checks directly.

\begin{lemma}\label{lem:hermite}
Every integer matrix $\alpha$ with $\det\alpha=p$ can be written as $\alpha=\gamma\rho$ with $\gamma\in\mathrm{SL}_2(\Z)$ and $\rho\in R_p$, and $\rho$ is unique.
\end{lemma}

\begin{proof}
\emph{Existence.} Let $\alpha=\bigl(\begin{smallmatrix}a&b\\c&d\end{smallmatrix}\bigr)$ and $g=\gcd(a,c)>0$. By Bézout there are $x,y$ with $xa+yc=g$, and $\gamma_1=\bigl(\begin{smallmatrix}x&y\\-c/g&a/g\end{smallmatrix}\bigr)\in\mathrm{SL}_2(\Z)$ satisfies
$\gamma_1\alpha=\bigl(\begin{smallmatrix}g&b'\\0&d'\end{smallmatrix}\bigr)$, with $gd'=p$, so $d'>0$. Multiplying on the left by $T^{-m}$ replaces $b'$ by $b'-md'$, so we may take $0\le b'<d'$. As $p$ is prime, either $g=p$, $d'=1$, $b'=0$,
or $g=1$, $d'=p$, $0\le b'<p$: the result is in $R_p$.

\emph{Uniqueness.} If $\gamma\rho=\gamma'\rho'$ with $\rho,\rho'\in R_p$, then $\rho'\rho^{-1}=\gamma'^{-1}\gamma\in\mathrm{SL}_2(\Z)$. Both are upper triangular, so $\rho'\rho^{-1}$ is upper triangular with positive rational diagonal entries $\frac{g'}g$, $\frac{d'}d$, which
are integers of product $1$; so $g=g'$ and $d=d'$. Its upper right entry is $\frac{b'-b}d$, an integer, so $b\equiv b'\pmod d$, and $b=b'$.
\end{proof}

\begin{theorem}[Hecke]\label{thm:hecke}
If $f\in M_k$, then $T_pf\in M_k$, and $T_p$ maps $S_k$ to $S_k$.
\end{theorem}

\begin{proof}
$T_pf$ is holomorphic on $\mathbb H$, and holomorphic at infinity by its $q$-expansion, whose constant term is $(1+p^{k-1})a_0$; so cusp forms go to cusp forms. Let $\gamma\in\mathrm{SL}_2(\Z)$. For $\rho\in R_p$, $\det(\rho\gamma)=p$, so by
Lemma~\ref{lem:hermite} $\rho\gamma=\gamma_\rho\rho'$ for some $\gamma_\rho\in\mathrm{SL}_2(\Z)$ and a unique $\rho'=\pi(\rho)\in R_p$. The map $\pi$ is injective: if $\pi(\rho_1)=\pi(\rho_2)=\rho'$, then $\rho_1=\gamma_{\rho_1}\rho'\gamma^{-1}$ and
$\rho_2=\gamma_{\rho_2}\rho'\gamma^{-1}$ give $\rho_1=(\gamma_{\rho_1}\gamma_{\rho_2}^{-1})\rho_2$, and uniqueness in the lemma (applied to $\alpha=\rho_1$) forces $\rho_1=\rho_2$. So $\pi$ is a permutation of $R_p$, and
\[
\begin{aligned}(T_pf)|_k\gamma&=\sum_\rho f|_k\rho\gamma=\sum_\rho(f|_k\gamma_\rho)|_k\pi(\rho)\eqbr=\sum_\rho f|_k\pi(\rho)=T_pf.\end{aligned}
\]
\end{proof}

\begin{theorem}[Mordell, 1917]\label{thm:mordell}
For every prime $p$ and $n\ge1$, $\tau(p)\tau(n)=\tau(pn)+p^{11}\tau(n/p)$, with $\tau(n/p)=0$ if $p\nmid n$. Consequently $\tau(mn)=\tau(m)\tau(n)$ whenever $\gcd(m,n)=1$, and
$\tau(p^{r+1})=\tau(p)\tau(p^r)-p^{11}\tau(p^{r-1})$ for $r\ge1$.
\end{theorem}

\begin{proof}
By Theorem~\ref{thm:hecke}, $T_p\Delta\in S_{12}=\C\Delta$, so $T_p\Delta=\lambda\Delta$, and comparing the coefficients of $q$, $\lambda=\tau(p)+p^{11}\tau(1/p)=\tau(p)$. Comparing the coefficients of $q^n$ gives the first identity.
With $n=p^r$ it is the recursion. With $p\nmid n$ it gives $\tau(pn)=\tau(p)\tau(n)$; more generally, for $p\nmid n$ the sequences $r\mapsto\tau(p^rn)$ and $r\mapsto\tau(p^r)\tau(n)$ satisfy the same recursion (apply the first identity to $p^rn$)
and agree for $r=0,1$, so $\tau(p^rn)=\tau(p^r)\tau(n)$. Induction on the number of prime factors of $m$ gives $\tau(mn)=\tau(m)\tau(n)$ for coprime $m,n$.
\end{proof}

So $\tau$ is multiplicative because $S_{12}$ is a line. The recursion says that the Dirichlet series of $\tau$ has an Euler product, like $\zeta$ and the $L(s,\chi)$ of this book:
\[
\sum_{n\ge1}\frac{\tau(n)}{n^s}=\prod_p\frac1{1-\tau(p)p^{-s}+p^{11-2s}},
\]
and Hecke proved that it extends to an entire function with a functional equation, just as Riemann did for $\zeta$.

\begin{remark}[Ramanujan's third observation, and beyond]
Write $1-\tau(p)X+p^{11}X^2=(1-\alpha_pX)(1-\beta_pX)$. Ramanujan's inequality $|\tau(p)|\le2p^{11/2}$ says exactly that $\alpha_p$ and $\beta_p$ are complex conjugates of modulus $p^{11/2}$: a Riemann hypothesis
for the local factors. Deligne deduced it in 1974 from his proof of the Riemann hypothesis for varieties over finite fields. Write then $\tau(p)=2p^{11/2}\cos\theta_p$. The Sato--Tate conjecture, proved in 2011 by Barnet-Lamb, Geraghty, Harris and Taylor,
says that the angles $\theta_p$ are distributed with density $\frac2\pi\sin^2\theta$ on $[0,\pi]$, so that $x_p=\cos\theta_p$ has density $\frac2\pi\sqrt{1-x^2}$ (Figure~\ref{fig:satotate}). Lehmer asked in 1947 whether $\tau(n)$ is ever $0$;
nobody knows. The same circle of ideas, for modular forms of weight $2$ attached to elliptic curves, is how Wiles proved Fermat's Last Theorem in 1995.
\end{remark}

\begin{figure}[ht]
\centering
\fig{ch12-satotate}
\caption{The normalised values $x_p=\frac{\tau(p)}{2p^{11/2}}$ for the $5133$ primes $p\le50\,000$ (histogram), against the Sato--Tate density $\frac2\pi\sqrt{1-x^2}$ (orange). All the values lie in $[-1,1]$, as Deligne's
theorem says; the largest $|x_p|$ is $0.990$.}
\label{fig:satotate}
\end{figure}

\begin{beauty}
Ramanujan's three observations about the coefficients of one power series turned out to be shadows of three structures: a finite-dimensional space (the $691$), a commuting family of operators acting on it
(multiplicativity), and the geometry of varieties over finite fields (Deligne's bound). The symmetry group of Figure~\ref{fig:tessellation} is the reason for all three. This book began with Gauss counting primes; it ends at the door
of the modern theory, where the primes, the $L$-functions of Chapters~3 to~7, and modular forms are parts of a single landscape. Volume~12 opens that door.
\end{beauty}

\section{Exercises}

\begin{exercise}[$\star$]\label{ex:sigma9}
Prove that $11\sigma_9(n)=21\sigma_5(n)-10\sigma_3(n)+5040\sum_{m=1}^{n-1}\sigma_3(m)\sigma_5(n-m)$ for every $n\ge1$.
\end{exercise}

\begin{exercise}[$\star$]\label{ex:e2}
Let $E_2=1-24\sum_{n\ge1}\sigma_1(n)q^n$. Show that $E_2$ is holomorphic on $\mathbb H$ and $E_2(\tau+1)=E_2(\tau)$, but that $E_2$ is not a modular form of weight $2$.
\end{exercise}

\begin{exercise}[$\star$]\label{ex:zeros}
Show that $E_6(i)=0$ and $E_4(\rho)=0$, where $\rho=e^{2\pi i/3}$. Deduce that $\Delta(i)=\frac{E_4(i)^3}{1728}$.
\end{exercise}

\begin{exercise}[$\star$]\label{ex:tausmall}
Using Theorem~\ref{thm:mordell} and $\tau(2)=-24$, $\tau(3)=252$, compute $\tau(4)$, $\tau(6)$, $\tau(8)$ and $\tau(12)$, and check them against the product. Check Corollary~\ref{cor:delta} for $n=2$ and $n=3$.
\end{exercise}

\begin{exercise}[$\star\star$]\label{ex:eisensteinhecke}
Show that $T_pE_k=(1+p^{k-1})E_k$, and that $T_pT_\ell=T_\ell T_p$ on $M_k$ for distinct primes $p,\ell$.
\end{exercise}

\begin{exercise}[$\star\star$, the signature exercise]\label{ex:heckebound}
(a) Show that $\phi(\tau)=(\operatorname{Im}\tau)^6|\Delta(\tau)|$ is invariant under $\mathrm{SL}_2(\Z)$ and bounded on $\mathbb H$.
(b) Deduce Hecke's bound $\tau(n)=O(n^6)$, by writing $\tau(n)=\int_0^1\Delta(x+iy)e^{-2\pi in(x+iy)}\dd x$ and choosing $y=\frac1n$. (Deligne's theorem gives $O(n^{11/2+\varepsilon})$.)
\end{exercise}

\begin{exercise}[$\star\star\star$]\label{pb:singer}
\emph{Singer's difference sets (1938).} Let $q$ be a prime power, $N=q^2+q+1$, and $g$ a generator of $\mathbb F_{q^3}^*$. View
$\mathbb F_{q^3}$ as a $3$-dimensional vector space over $\mathbb F_q$, fix a $2$-dimensional subspace $U$, and let
$D=\{i\in\Z/N\Z:g^i\in U\}$.
\begin{enumerate}[label=(\alph*)]
\item Show that $D$ is well defined, that $|D|=q+1$, and that every non-zero $d\in\Z/N\Z$ is a difference $i-j$ of elements of
$D$ in exactly one way.
\item Show that the $N$ translates $D+t$ are the lines of a projective plane on the points $\Z/N\Z$: two points lie on exactly one
translate, and two translates meet in exactly one point.
\item Check that $\{0,1,3\}\bmod7$, $\{0,1,3,9\}\bmod13$ and $\{0,1,4,14,16\}\bmod21$ are perfect difference sets.
\end{enumerate}
\end{exercise}

\begin{exercise}[$\star\star$]\label{pb:consecsqfree}
\emph{Squarefree neighbours.}
\begin{enumerate}[label=(\alph*)]
\item Show that the density of the integers $n$ such that $n$ and $n+1$ are both squarefree is
$\prod_p\bigl(1-\frac2{p^2}\bigr)=0.3226\ldots$
\item Show that $n,n+1,n+2$ are all squarefree with density $\prod_p\bigl(1-\frac3{p^2}\bigr)=0.1254\ldots$, and that no four
consecutive integers are all squarefree.
\item Show that there are arbitrarily long runs of consecutive integers none of which is squarefree.
\end{enumerate}
\end{exercise}

\begin{exercise}[$\star\star$]\label{pb:pairwise}
\emph{Pairwise coprime.} Let $k\ge2$.
\begin{enumerate}[label=(\alph*)]
\item Show that $k$ random integers are pairwise coprime with probability
\[
\prod_p\Bigl(1-\frac1p\Bigr)^{k-1}\Bigl(1+\frac{k-1}p\Bigr).
\]
\item For $k=3$ this is $0.2867\ldots$, larger than $(6/\pi^2)^3=0.2247\ldots$, the value if the three events
``$\gcd(a,b)=1$'', ``$\gcd(b,c)=1$'', ``$\gcd(a,c)=1$'' were independent. Explain the sign of the difference prime by prime.
\end{enumerate}
\end{exercise}

\begin{exercise}[$\star\star\star$]\label{pb:twosqfree}
\emph{Two squarefree summands.} Show that every integer $n\ge2$ is the sum of two squarefree positive integers. Prove it for
$n\ge14000$ by counting, and check the smaller $n$ by computer.
\end{exercise}

\begin{exercise}[$\star\star$]\label{pb:dsquared}
\emph{Ramanujan's $\sum d(n)^2$.} Let $d(n)$ be the number of divisors of $n$.
\begin{enumerate}[label=(\alph*)]
\item Show that $\sum_{n\ge1}d(n)^2n^{-s}=\frac{\zeta(s)^4}{\zeta(2s)}$ for $s>1$.
\item Deduce that $\sum_{n\le x}d(n)^2=\frac1{\pi^2}x\log^3x+O(x\log^2x)$.
\end{enumerate}
\end{exercise}

\begin{exercise}[$\star\star$]\label{pb:energy}
\emph{Energy and the multiplication table.} For a finite set $A$ of numbers, its \emph{additive energy} $E_+(A)$ is the number of
quadruples $(a,b,c,d)\in A^4$ with $a+b=c+d$, and its multiplicative energy $E_\times(A)$ the number with $ab=cd$.
\begin{enumerate}[label=(\alph*)]
\item Show that $|A+A|\ge\frac{|A|^4}{E_+(A)}$ and $|A\cdot A|\ge\frac{|A|^4}{E_\times(A)}$.
\item Let $A=\{1,\dots,N\}$. Show that $E_\times(A)=\sum\bigl\lfloor\frac N{\max(u,v)}\bigr\rfloor^2$, the sum running over the pairs of coprime
integers $1\le u,v\le N$, and deduce $E_\times(A)=\frac{12}{\pi^2}N^2\log N+O(N^2)$.
\item Deduce that the $N\times N$ multiplication table contains at least $\bigl(\frac{\pi^2}{12}+o(1)\bigr)\frac{N^2}{\log N}$ distinct numbers.
\end{enumerate}
\end{exercise}

\begin{exercise}[$\star\star$]\label{pb:kleitman}
\emph{Harris and Kleitman.} A family $\mathcal A$ of subsets of $[n]$ is \emph{increasing} if $S\in\mathcal A$ and $S\subseteq T$ imply $T\in\mathcal A$,
and \emph{decreasing} if its complement in $2^{[n]}$ is increasing.
\begin{enumerate}[label=(\alph*)]
\item (Harris, 1960; Kleitman, 1966.) Show that $2^n|\mathcal A\cap\mathcal B|\ge|\mathcal A||\mathcal B|$ if $\mathcal A,\mathcal B$ are both increasing or both decreasing.
\item A family is \emph{intersecting} if any two members meet. Show that the union of $k$ intersecting families of subsets of $[n]$ has
at most $2^n-2^{n-k}$ members, and that this is attained.
\end{enumerate}
\end{exercise}

\begin{exercise}[$\star$]\label{pb:chebyshevsum}
\emph{Chebyshev's sum inequality.}
\begin{enumerate}[label=(\alph*)]
\item Let $a_1\le\dots\le a_n$ and $b_1\le\dots\le b_n$. Show that $\frac1n\sum_ia_ib_i\ge\bigl(\frac1n\sum_ia_i\bigr)\bigl(\frac1n\sum_ib_i\bigr)$.
\item Let $f,g$ be non-decreasing functions and $X$ a real random variable with $f(X),g(X)$ square-integrable. Show that
$\mathbb E[f(X)g(X)]\ge\mathbb E[f(X)]\,\mathbb E[g(X)]$.
\item Deduce that the moments $M_k=\lim\frac1N\sum_{n\le N}\bigl(\frac{\varphi(n)}n\bigr)^k$ of Problem~\ref{pb:phiovern} satisfy
$M_{j+k}\ge M_jM_k$, and check it on $M_1=0.6079\ldots$, $M_2=0.4282\ldots$
\end{enumerate}
\end{exercise}

\begin{exercise}[$\star\star$]\label{pb:primitive}
\emph{Primitive sets (Erdős, 1935).} A set $A$ of integers $\ge2$ is \emph{primitive} if no element divides another. For $a\ge2$ let $P(a)$ be
its largest prime factor, and $M_a$ the set of integers $am$ whose every prime factor $q$ of $m$ satisfies $q\ge P(a)$.
\begin{enumerate}[label=(\alph*)]
\item Show that the sets $M_a$, $a\in A$, are pairwise disjoint when $A$ is primitive.
\item Show that $M_a$ has density $\frac1a\prod_{p<P(a)}\bigl(1-\frac1p\bigr)$.
\item Deduce that $\sum_{a\in A}\frac1{a\log a}\le C$ for every primitive set $A$, with an absolute constant $C$.
\end{enumerate}
\end{exercise}

\begin{exercise}[$\star\star$]\label{pb:divisorbound}
\emph{Champions of divisors.} Let $d(n)$ be the number of divisors of $n$ and $\alpha>0$. (Volume~5, Exercise~2.6, treats $\alpha=\frac12$.)
\begin{enumerate}[label=(\alph*)]
\item Show that $\max_n\frac{d(n)}{n^\alpha}$ exists and equals $\prod_p\max_{e\ge0}\frac{e+1}{p^{e\alpha}}$, a product with finitely many factors
different from $1$.
\item Let $n_\alpha$ be the smallest $n$ at which this maximum is reached, and write $n_\alpha=\prod_pp^{e_p}$. Show that $e_p$ is the least
$e\ge0$ with $\frac{e+2}{e+1}\le p^\alpha$. Deduce that $e_2\ge e_3\ge e_5\ge\cdots$, and that $p$ divides $n_\alpha$ if and only if $p^\alpha<2$.
\item Show that $d(n)\le3.53\,n^{1/3}$ for all $n$, with the maximum of $d(n)/n^{1/3}$ reached at $n=2520$, and find $n_{1/4}$.
\end{enumerate}
\end{exercise}

\begin{exercise}[$\star\star$]\label{pb:logtwo}
\emph{The irrationality of $\log2$.}
\begin{enumerate}[label=(\alph*)]
\item Show that $P_n$ has integer coefficients, and that
\[
I_n=\int_0^1\frac{P_n(x)}{1+x}\dd x=\int_0^1\frac{x^n(1-x)^n}{(1+x)^{n+1}}\dd x.
\]
\item Show that $d_nI_n=a_n+b_n\log2$ with $a_n,b_n\in\Z$.
\item Show that $0<I_n\le(3-2\sqrt2)^n$, and deduce that $\log2$ is irrational.
\end{enumerate}
\end{exercise}

\begin{exercise}[$\star$]\label{pb:doubleint}
\emph{Double integrals that know $\zeta(2)$ and $\zeta(3)$.} Let $r,s\ge0$ be integers.
\begin{enumerate}[label=(\alph*)]
\item Show that $\displaystyle\int_0^1\!\!\int_0^1\frac{x^ry^s}{1-xy}\dd x\dd y$ equals $\frac1{r-s}\bigl(\frac1{s+1}+\cdots+\frac1r\bigr)$ if $r>s$, and
$\zeta(2)-\sum_{k=1}^r\frac1{k^2}$ if $r=s$.
\item Show that $\displaystyle\int_0^1\!\!\int_0^1\frac{-\log(xy)}{1-xy}\,x^ry^s\dd x\dd y$ equals $\frac1{r-s}\bigl(\frac1{(s+1)^2}+\cdots+\frac1{r^2}\bigr)$ if $r>s$, and
$2\bigl(\zeta(3)-\sum_{k=1}^r\frac1{k^3}\bigr)$ if $r=s$.
\item Show that $\displaystyle\int_0^1\!\!\int_0^1\frac{\mathrm dx\dd y}{1+xy}=\frac{\pi^2}{12}$.
\end{enumerate}
\end{exercise}

\begin{exercise}[$\star\star\star$]\label{pb:beukerstwo}
\emph{Beukers's proof for $\zeta(2)$.} Let
\[
J_n=\int_0^1\!\!\int_0^1\frac{(1-y)^nP_n(x)}{1-xy}\dd x\dd y.
\]
\begin{enumerate}[label=(\alph*)]
\item Show that $d_n^2J_n=a_n+b_n\zeta(2)$ with $a_n,b_n\in\Z$.
\item Show that
\[
J_n=(-1)^n\int_0^1\!\!\int_0^1\frac{x^n(1-x)^ny^n(1-y)^n}{(1-xy)^{n+1}}\dd x\dd y.
\]
\item Show that $\frac{x(1-x)y(1-y)}{1-xy}\le\bigl(\frac{\sqrt5-1}2\bigr)^5$ on $[0,1]^2\setminus\{(1,1)\}$, and deduce that $0<|J_n|\le\zeta(2)\bigl(\frac{\sqrt5-1}2\bigr)^{5n}$.
\item Deduce that $\zeta(2)$ is irrational. (This gives a second proof that $\pi^2$ is irrational.)
\end{enumerate}
\end{exercise}

\begin{exercise}[Journey]\label{pb:beukersthree}
\emph{Beukers's proof that $\zeta(3)$ is irrational.} Let
\[
K_n=\int_0^1\!\!\int_0^1\frac{-\log(xy)}{1-xy}\,P_n(x)P_n(y)\dd x\dd y.
\]
\begin{enumerate}[label=(\alph*)]
\item Show that $d_n^3K_n=a_n+b_n\zeta(3)$ with $a_n,b_n\in\Z$.
\item Show that $\frac{-\log(xy)}{1-xy}=\int_0^1\frac{\mathrm dz}{1-(1-xy)z}$ for $0<xy<1$.
\item Integrating by parts $n$ times in $x$, then substituting $w=\frac{1-z}{1-(1-xy)z}$, then integrating by parts $n$ times in $y$, show that
\[
K_n=\int_0^1\!\!\int_0^1\!\!\int_0^1\frac{x^n(1-x)^ny^n(1-y)^nw^n(1-w)^n}{\bigl(1-(1-xy)w\bigr)^{n+1}}\dd x\dd y\dd w.
\]
\item Show that $\frac{x(1-x)y(1-y)w(1-w)}{1-(1-xy)w}\le(\sqrt2-1)^4$ on $[0,1]^3$ (where defined), and that $0<K_n\le2\zeta(3)(\sqrt2-1)^{4n}$.
\item Using the prime number theorem in the form $\log d_n\sim n$, conclude that $\zeta(3)$ is irrational.
\end{enumerate}
\end{exercise}

\begin{exercise}[$\star\star$]\label{pb:fracint}
\emph{Integrals of $\{1/x\}$.}
\begin{enumerate}[label=(\alph*)]
\item Show that $\displaystyle\int_0^1\Bigl\{\frac1x\Bigr\}\dd x=1-\gamma$.
\item Show that $\displaystyle\int_0^1\Bigl\{\frac1x\Bigr\}^2\dd x=\log(2\pi)-\gamma-1$.
\item Show that $\displaystyle\int_0^1\frac{\mathrm dx}{\lfloor1/x\rfloor}=\zeta(2)-1$.
\end{enumerate}
\end{exercise}

\begin{exercise}[$\star\star\star$]\label{pb:primefrac}
\emph{Remainders modulo primes (de la Vallée Poussin, 1898).} Show that
\[
\frac1{\pi(x)}\sum_{p\le x}\Bigl\{\frac xp\Bigr\}\longrightarrow1-\gamma\qquad(x\to\infty).
\]
\end{exercise}

\begin{exercise}[$\star$]\label{pb:leastpf}
\emph{Sums of primes, and of least prime factors.}
\begin{enumerate}[label=(\alph*)]
\item Show that $\sum_{p\le x}p\sim\frac{x^2}{2\log x}$.
\item Show that $\sum_{2\le n\le x}P^-(n)\sim\frac{x^2}{2\log x}$: the least prime factors of all the integers up to $x$ add up to essentially the sum of the primes alone.
\item Show that the sum of the first $n$ primes is asymptotic to $\frac12n^2\log n$.
\item Numerically, one finds $\sum_{p\le10^7}p=3\,203\,324\,994\,356$, which is $1.0326\,\frac{x^2}{2\log x}$ but $0.99950\operatorname{Li}(x^2)$ for $x=10^7$. Explain why
$\operatorname{Li}(x^2)$ is the natural approximation, and show that $\sum_{p\le x}p=\operatorname{Li}(x^2)+O(x^2e^{-c\sqrt{\log x}})$ for some $c>0$.
\end{enumerate}
\end{exercise}

\begin{exercise}[$\star\star\star$]\label{pb:largestpf}
\emph{The largest prime factor on average.} For $n\ge2$ let $A(n)$ be the sum of the prime factors of $n$ counted with multiplicity, so $A(12)=2+2+3=7$. For $y\ge1$ and
$z\ge2$, let $\Psi(y,z)$ be the number of $m\le y$ with no prime factor larger than $z$ (with $m=1$ counted).
\begin{enumerate}[label=(\alph*)]
\item Show that $\sum_{2\le n\le x}P^+(n)=\sum_{p\le x}p\,\Psi(\frac xp,p)$.
\item Show that the primes $p\le\sqrt x$ contribute $O(x^{3/2})$, and that $\Psi(\frac xp,p)=\lfloor\frac xp\rfloor$ when $p>\sqrt x$.
\item Show that $\sum_{\sqrt x<p\le x}p\lfloor\frac xp\rfloor=\sum_{1\le k<\sqrt x}\Bigl(\sum_{p\le x/k}p-\sum_{p\le\sqrt x}p\Bigr)$, and deduce that
\[
\sum_{2\le n\le x}P^+(n)\sim\frac{\pi^2}{12}\,\frac{x^2}{\log x}.
\]
\item Show that $A(n)-P^+(n)\le\sqrt n\log_2n$, and deduce the same asymptotic formula for $\sum_{n\le x}A(n)$.
\item Numerically, one finds $\sum_{n\le10^7}P^+(n)=1.0767\cdot\frac{\pi^2}{12}\frac{x^2}{\log x}$. Why is the ratio larger than $1$, and why does it approach $1$ so slowly?
\end{enumerate}
\end{exercise}

\begin{exercise}[$\star\star$]\label{pb:factorialdiv}
\emph{How many divisors has $n!$?} Let $\tau(m)$ be the number of divisors of $m$, and $C=\sum_{k\ge1}\frac1k\log\bigl(1+\frac1k\bigr)=1.25774\ldots$
\begin{enumerate}[label=(\alph*)]
\item Show that $\tau(n!)=\prod_{p\le n}\bigl(1+v_p(n!)\bigr)$, where $v_p(n!)=\sum_{j\ge1}\lfloor\frac n{p^j}\rfloor$.
\item Show that the primes $p\le\sqrt n$ contribute $O(\sqrt n)$ to $\log\tau(n!)$, and that $v_p(n!)=\lfloor\frac np\rfloor$ for $p>\sqrt n$.
\item Show that $\log\tau(n!)\sim C\frac n{\log n}$. In other words, $\tau(n!)=e^{C(1+o(1))n/\log n}$, where $e^C=2\cdot\sqrt{\frac32}\cdot\sqrt[3]{\frac43}\cdots=3.5174\ldots$
\item Numerically, one finds $\log\tau(n!)\big/\frac n{\log n}=1.58$, $1.51$, $1.45$, $1.42$ for $n=10^4,10^5,10^6,10^7$. Comment.
\end{enumerate}
\end{exercise}

\begin{exercise}[$\star\star$]\label{pb:centralbinom}
\emph{The prime factors of $\binom{2n}n$.} Every prime in $]n,2n]$ divides $\binom{2n}n$, and no prime in $]\frac{2n}3,n]$ does once $n\ge5$ (this is the heart of Erd\H os's proof of Bertrand's postulate).
\begin{enumerate}[label=(\alph*)]
\item Show that for $p>\sqrt{2n}$, $v_p\bigl(\binom{2n}n\bigr)=\lfloor\frac{2n}p\rfloor-2\lfloor\frac np\rfloor$, which is $1$ if the fractional part of $\frac np$ is at least $\frac12$, and $0$ otherwise.
\item Deduce that the primes $p>\sqrt{2n}$ dividing $\binom{2n}n$ are those of $\bigcup_{k\ge0}\bigl]\frac n{k+1},\frac n{k+1/2}\bigr]$.
\item Show that $\omega\bigl(\binom{2n}n\bigr)\sim\log4\cdot\frac n{\log n}$: the central binomial coefficient has about $39\%$ more prime factors than the primes of $]n,2n]$ alone.
\item Numerically, one finds $\omega\bigl(\binom{2n}n\bigr)=453\,587$ for $n=5\cdot10^6$. Compare.
\end{enumerate}
\end{exercise}

\begin{exercise}[$\star$]\label{pb:superprimes}
\emph{Primes with a prime index.} The primes $p_{p_n}$ are $3,5,11,17,31,41,59,\dots$
\begin{enumerate}[label=(\alph*)]
\item Show that $p_{p_n}\sim n\log^2n$ and that the number of them up to $x$, $\pi(\pi(x))$, is asymptotic to $\frac x{\log^2x}$.
\item Show that $\sum_n\frac1{p_{p_n}}$ converges. Up to $10^8$ the partial sum is $0.9833$; how large, roughly, is the remaining tail?
\item Numerically, one finds $\pi(\pi(10^8))=397\,557$, which is $1.349\,\frac x{\log^2x}$. Find a better approximation using $\pi(y)\approx\frac y{\log y}(1+\frac1{\log y})$.
\end{enumerate}
\end{exercise}

\begin{exercise}[$\star\star$]\label{pb:ramprimes}
\emph{Ramanujan primes.} In 1919 Ramanujan gave a short proof of Bertrand's postulate which shows more: $\pi(x)-\pi(\frac x2)\to\infty$. Let $f(x)=\pi(x)-\pi(\frac x2)$ for real $x>0$.
\begin{enumerate}[label=(\alph*)]
\item For $n\ge1$ let $R_n$ be the least real number such that $f(x)\ge n$ for all $x\ge R_n$. Show that $R_n$ exists and is a prime, and that $f(R_n)=n$.
(So $R_1=2$, $R_2=11$, $R_3=17$, $R_4=29$, $R_5=41$; Sondow named them Ramanujan primes in 2009.)
\item Show that $R_n\sim p_{2n}$.
\item Numerically, one finds $R_{10\,000}=242\,057$ while $p_{20\,000}=224\,737$. Using $\pi(y)\approx\frac y{\log y}(1+\frac1{\log y})$, show that one should expect
$\frac{R_n}{p_{2n}}\approx1+\frac{\log2}{\log p_{2n}}$, and compare.
\end{enumerate}
\end{exercise}

\begin{exercise}[$\star\star\star$]\label{pb:pminusone}
\emph{Artin's constant twice.} Let $A=\prod_q\bigl(1-\frac1{q(q-1)}\bigr)=0.3739558\ldots$, the product over all primes $q$.
\begin{enumerate}[label=(\alph*)]
\item (\emph{The large squares are rare.}) Show that for $y\ge2$ and $x\ge x_0$, the number of primes $p\le x$ such that $q^2\mid p-1$ for some prime $q>y$ is at most $\frac{10x}{y\log x}+2x^{3/4}$.
\item Show that the proportion of primes $p\le x$ for which $p-1$ is squarefree tends to $A$.
\item Show that $\frac1{\pi(x)}\sum_{p\le x}\frac{\varphi(p-1)}{p-1}\to A$ as well.
\item Explain why (c) says that a random residue modulo a random prime is a primitive root with probability $A$. Numerically, one finds that $2$ is a primitive root modulo
$248\,491$ of the $664\,578$ odd primes below $10^7$, a proportion of $0.37391$. (That this proportion tends to $A$ is Artin's conjecture, still unproved.)
\end{enumerate}
\end{exercise}

\begin{exercise}[$\star\star$]\label{pb:titchmarsh}
\emph{The divisors of $p-1$.} Let $Z=\frac{\zeta(2)\zeta(3)}{\zeta(6)}=1.94359\ldots$ and $B=\prod_q\bigl(1+\frac1{(q-1)^2}\bigr)=2.82642\ldots$
\begin{enumerate}[label=(\alph*)]
\item Show that $\sum_{d\le y}\frac1{\varphi(d)}\le c(1+\log y)$ for some constant $c$, writing $\frac1{\varphi(d)}=\frac1d\sum_{e\mid d}\frac{\mu(e)^2}{\varphi(e)}$.
\item Show that $\sum_{p\le x}\tau(p-1)=O(x)$.
\item Show that $\frac1{\pi(x)}\sum_{p\le x}\frac{p-1}{\varphi(p-1)}\to B$.
\item Assuming that $\sum_{d\le y}\frac1{\varphi(d)}\sim Z\log y$ (a theorem of Landau), give a heuristic argument that $\sum_{p\le x}\tau(p-1)\sim Zx$. Numerically, one finds
$\frac1x\sum_{p\le x}\tau(p-1)=1.916$, $1.927$, $1.932$ for $x=10^5,10^6,10^7$.
\end{enumerate}
\end{exercise}

\begin{exercise}[$\star\star$]\label{pb:menon}
\emph{The roots of $x^n-1$ modulo $p$, on average.} Fix $n\ge1$.
\begin{enumerate}[label=(\alph*)]
\item Show that for every prime $p$, the polynomial $x^n-1$ has exactly $\gcd(n,p-1)$ roots in $\mathbb F_p$.
\item (\emph{Menon's identity, 1965.}) Let the group $(\Z/n\Z)^\times$ act on $\Z/n\Z$ by multiplication. Count the orbits in two ways and deduce that
\[
\sum_{a\in(\Z/n\Z)^\times}\gcd(a-1,n)=\varphi(n)\tau(n).
\]
\item Deduce that the average number of roots of $x^n-1$ modulo $p$, over the primes $p\le x$, tends to $\tau(n)$.
\item Show directly that for $d\mid n$, the cyclotomic polynomial $\Phi_d$ has on average one root modulo $p$, and explain (b) and (c) through the factorisation
$x^n-1=\prod_{d\mid n}\Phi_d(x)$. Numerically, one finds the averages $5.994$, $7.983$ and $11.979$ for $n=12,24,60$, over the primes below $10^6$.
\end{enumerate}
\end{exercise}

\begin{exercise}[$\star\star\star$]\label{pb:omegapoly}
\emph{How many prime factors has $n^2+1$?} Let $\rho(p)$ be the number of solutions of $n^2\equiv-1\pmod p$, and $\chi$ the non-principal character modulo $4$.
\begin{enumerate}[label=(\alph*)]
\item Show that $\rho(2)=1$ and $\rho(p)=1+\chi(p)$ for odd $p$.
\item Show that $n^2+1$ has at most one prime factor larger than $x$ when $n\le x$, and deduce that $\sum_{n\le x}\omega(n^2+1)=x\sum_{p\le x}\frac{\rho(p)}p+O(x)$.
\item Deduce that $\sum_{n\le x}\omega(n^2+1)=x\log\log x+O(x)$.
\item Show that $\sum_{n\le x}\bigl(\omega(n^2+1)-\log\log x\bigr)^2=O(x\log\log x)$, and deduce that $\omega(n^2+1)$ has normal order $\log\log n$.
\item Numerically, one finds $\frac1x\sum_{n\le x}\omega(n^2+1)-\log\log x=0.598$, $0.603$, $0.606$ for $x=10^4,10^5,10^6$. Is this consistent with (c)?
\end{enumerate}
\end{exercise}

\begin{exercise}[$\star\star$]\label{pb:rankin}
\emph{Rankin's trick.}
\begin{enumerate}[label=(\alph*)]
\item Show that $\Psi(x,y)\le x^\sigma\prod_{p\le y}(1-p^{-\sigma})^{-1}$ for every $\sigma\in\left]0,1\right]$.
\item Take $\sigma=1-\frac\eta{\log y}$ with $0<\eta\le\frac13\log y$. Using $e^{\eta t}\le1+(e^\eta-1)t$ for $0\le t\le1$ and the first two theorems of Mertens, show that
there is an absolute constant $C$ with
\[
\Psi(x,y)\le Cx\log y\,\exp\bigl(-\eta u+Ce^\eta\bigr).
\]
\item Choose $\eta=\log u$ and deduce that $\Psi(x,y)\le Cx\log y\,e^{-u\log u+Cu}$ when $u\ge1$ and $\log u\le\frac13\log y$. In particular, for $y=e^{\sqrt{\log x}}$,
the proportion of $y$-smooth integers up to $x$ is at most $e^{-(\frac12+o(1))\sqrt{\log x}\log\log x}$. (This is the scale on which the quadratic sieve lives.)
\item Direct computation gives $\Psi(10^{10},100)=8\,800\,084$ and $\Psi(10^{12},1000)=6\,471\,274\,933$, and minimises the bound of (a) over $\sigma$: it gets
$9.95\cdot10^7$ and $1.26\cdot10^{11}$. By what factor does the best Rankin bound miss? Where is the loss?
\end{enumerate}
\end{exercise}

\begin{exercise}[$\star\star\star$]\label{pb:renyi}
\emph{The law of $\Omega(n)-\omega(n)$.} Volume~5, Exercise~9.1, shows that $\Omega(n)-\omega(n)$ is bounded on average. Here is its exact law (R\'enyi, 1955).
\begin{enumerate}[label=(\alph*)]
\item For complex $z$ with $|z|<2$, show that $f_z(n)=z^{\Omega(n)-\omega(n)}$ is multiplicative and $f_z(n)=\sum_{d\mid n}h_z(d)$, where $h_z(p)=0$ and $h_z(p^k)=z^{k-2}(z-1)$ for $k\ge2$.
Show that $\sum_d\frac{|h_z(d)|}d<\infty$.
\item Deduce from Wintner's theorem (Volume~5, Theorem~2.2) that
\[
\frac1x\sum_{n\le x}z^{\Omega(n)-\omega(n)}\to R(z)=\prod_p\Bigl(1-\frac1p\Bigr)\Bigl(1+\frac1{p-z}\Bigr).
\]
\item Show that for every $k\ge0$ the set $\{n:\Omega(n)-\omega(n)=k\}$ has a density $d_k$, and that $\sum_kd_kz^k=R(z)$. (Write $d_k$ for $n\le x$ as an integral over the
unit circle.) Check that $d_0=\frac6{\pi^2}$ and that the mean of the law is $\sum_p\frac1{p(p-1)}=0.77315\ldots$
\item Show that $2^kd_k\to\frac1{4C_2}=0.37869\ldots$, where $C_2=\prod_{p\ge3}\bigl(1-\frac1{(p-1)^2}\bigr)$ is the twin prime constant.
\item Numerically, one finds $d_1=0.200756$, $d_2=0.096302$, $d_3=0.047615$, and, among the integers up to $10^7$, the proportions $0.200758$, $0.096306$, $0.047616$.
\end{enumerate}
\end{exercise}

\begin{exercise}[Journey]\label{pb:twosquares}
\emph{How many integers are sums of two squares?} Let $b(n)=1$ if $n\ge1$ is a sum of two squares and $b(n)=0$ otherwise, $B(x)=\sum_{n\le x}b(n)$, and
$F(s)=\sum_{n\ge1}b(n)n^{-s}$ for $s>1$. Let $\chi_4$ be the non-principal character modulo $4$.
\begin{enumerate}[label=(\alph*)]
\item Show that for $s>1$,
\[
F(s)^2=\frac{\zeta(s)L(s,\chi_4)}{1-2^{-s}}\prod_{p\equiv3\,(4)}\frac1{1-p^{-2s}},
\]
and deduce that $F(\sigma)\sim C(\sigma-1)^{-1/2}$ as $\sigma\to1^+$, with $C^2=\frac\pi2\prod_{p\equiv3\,(4)}(1-p^{-2})^{-1}$.
\item Let $S(y)=\sum_{n\le y}\frac{b(n)}n$ and $\delta=\frac1{\log y}$. Show that $S(y)\le e\,F(1+\delta)$ and $S(y)\ge F(1+\delta)-e^{-1/2}F(1+\frac\delta2)$, and
deduce that $S(y)$ lies between two positive multiples of $\sqrt{\log y}$ for large $y$.
\item Show that $\sum_{p^k\le y}k\log p=O(y)$ and $\sum_{n\le x}b(n)\log n\le c\,x\,S(x)$ for an absolute constant $c$, and deduce that $B(x)=O\bigl(\frac x{\sqrt{\log x}}\bigr)$.
\item Show that $\sum_{n\le x}b(n)\log n\ge\sum_{m\le x/\log^3x}b(m)\,\theta(\frac xm;4,1)$, where $\theta(y;4,1)=\sum_{p\le y,\,p\equiv1\,(4)}\log p$, and
deduce that $\frac{B(x)\sqrt{\log x}}x$ is bounded below by a positive constant.
\item Show that \emph{if} $B(x)\sim\frac{Kx}{\sqrt{\log x}}$ for some constant $K$, then $K=\frac C{\sqrt\pi}=0.7642\ldots$ Numerically, one finds $B(10^8)=18\,457\,847$, for which
$\frac{B(x)\sqrt{\log x}}x=0.7922$.
\end{enumerate}
\end{exercise}

\begin{exercise}[Journey]\label{pb:ramanujanE2}
\emph{Ramanujan's differential equations.} With $q=e^{2\pi i\tau}$ for $\tau$ in the upper half-plane, let $E_2=1-24\sum_{n\ge1}\sigma(n)q^n$, $E_4=1+240\sum_{n\ge1}\sigma_3(n)q^n$,
$E_6=1-504\sum_{n\ge1}\sigma_5(n)q^n$, $\eta(\tau)=q^{1/24}\prod_{n\ge1}(1-q^n)$, and $D=\frac1{2\pi i}\frac{\mathrm d}{\mathrm d\tau}=q\frac{\mathrm d}{\mathrm dq}$.
\begin{enumerate}[label=(\alph*)]
\item Show that $\frac{D\eta}\eta=\frac{E_2}{24}$.
\item Using Dedekind's transformation $\eta(-\frac1\tau)=\sqrt{-i\tau}\,\eta(\tau)$, show that $E_2(-\frac1\tau)=\tau^2E_2(\tau)+\frac{6\tau}{\pi i}$.
\item Show that $f=DE_2-\frac{E_2^2}{12}$ satisfies $f(-\frac1\tau)=\tau^4f(\tau)$, and deduce \emph{Ramanujan's equations}
\[
DE_2=\frac{E_2^2-E_4}{12},\qq DE_4=\frac{E_2E_4-E_6}3,\qq DE_6=\frac{E_2E_6-E_4^2}2.
\]
\item Deduce that for every $n\ge1$,
\[
\sum_{m=1}^{n-1}\sigma(m)\sigma(n-m)=\frac{5\sigma_3(n)+(1-6n)\sigma(n)}{12}.
\]
\item Deduce similarly that $\sum_{m=1}^{n-1}\sigma(m)\sigma_3(n-m)=\frac{21\sigma_5(n)+(10-30n)\sigma_3(n)-\sigma(n)}{240}$.
\item For a prime $p$, show that $\sum_{m=1}^{p-1}\sigma(m)\sigma(p-m)=\frac{(p-1)(5p^2-p-6)}{12}$, and check $p=7$ by hand.
\end{enumerate}
\end{exercise}

\begin{exercise}[Journey]\label{pb:landauzeros}
\emph{The zeros hear the primes.} Fix a real number $x>1$, and let $\Lambda(x)=\Lambda(n)$ if $x$ is an integer $n$, and $\Lambda(x)=0$ otherwise. The zeros
$\rho=\beta+i\gamma$ of $\zeta$ in the critical strip are counted with multiplicity.
\begin{enumerate}[label=(\alph*)]
\item Show that for $T\ge2$,
\[
\frac1{2\pi}\int_1^Tx^{2+it}\,\frac{\zeta'}\zeta(2+it)\dd t=-\frac T{2\pi}\Lambda(x)+O_x(1).
\]
\item Using $\xi(s)=\xi(1-s)$ (Volume~5, Proposition~6.2), show that
$\frac{\zeta'}\zeta(s)=-\frac{\zeta'}\zeta(1-s)+\log\pi-\frac12\frac{\Gamma'}\Gamma\bigl(\frac s2\bigr)-\frac12\frac{\Gamma'}\Gamma\bigl(\frac{1-s}2\bigr)$, and deduce that
$\int_1^Tx^{-1+it}\frac{\zeta'}\zeta(-1+it)\dd t=O_x(\log T)$.
\item Deduce \emph{Landau's formula}: for $T\ge3$,
\[
\sum_{0<\gamma\le T}x^\rho=-\frac T{2\pi}\Lambda(x)+O_x(\log^2T).
\]
\item Assume the Riemann hypothesis. Show that $\sum_{0<\gamma\le T}\cos(\gamma\log x)=-\frac T{2\pi}\frac{\Lambda(x)}{\sqrt x}+O_x(\log^2T)$: the frequencies $\gamma$
resonate exactly at the prime powers. (Figure~\ref{fig:landau} shows it with the first $2000$ zeros.)
\item Still under the Riemann hypothesis, let $0<\gamma_1\le\gamma_2\le\cdots$ be the positive ordinates. Show that for every real $\alpha\ne0$ the sequence
$(\alpha\gamma_n)$ is equidistributed modulo $1$.
\end{enumerate}
\end{exercise}

\begin{exercise}[Journey]\label{pb:hearzeros}
\emph{The primes hear the zeros.} Fix $s>0$ and $T\ge0$, let $h(t)=e^{-(t-T)^2/2s^2}+e^{-(t+T)^2/2s^2}$, and
\[
g(u)=\frac1{2\pi}\int_{-\infty}^\infty h(t)e^{-iut}\dd t=\frac{2s}{\sqrt{2\pi}}\cos(Tu)\,e^{-s^2u^2/2}.
\]
\begin{enumerate}[label=(\alph*)]
\item Let $H(z)=h(-i(z-\frac12))$. Show that $H$ is entire, that $H(1-z)=H(z)$, that $H(\frac12+it)=h(t)$, and that $H(\sigma+it)$ tends to $0$ faster than any
power of $|t|$, uniformly for $\sigma$ in a bounded interval.
\item Integrate $H(z)\frac{\zeta'}\zeta(z)$ around rectangles with vertices $-1\pm iT'$, $2\pm iT'$ to prove the \emph{Guinand--Weil explicit formula}
\[
\begin{aligned}\sum_\rho H(\rho)=H(0)+H(1)&-2\sum_{n\ge1}\frac{\Lambda(n)}{\sqrt n}g(\log n)-g(0)\log\pi\\&+\frac1{2\pi}\int_{-\infty}^\infty h(t)\operatorname{Re}\frac{\Gamma'}\Gamma\Bigl(\frac14+\frac{it}2\Bigr)\dd t,\end{aligned}
\]
the sum running over all non-trivial zeros, and converging absolutely.
\item Assume the Riemann hypothesis, and let $Z(T)=\sum_{\gamma>0}e^{-(\gamma-T)^2/2s^2}$, a sum of bumps of width $s$, one at each zero. Show that for $T\ge1$,
\[
Z(T)=\frac s{\sqrt{2\pi}}\log\frac T{2\pi}-\frac{2s}{\sqrt{2\pi}}\sum_{n\ge2}\frac{\Lambda(n)}{\sqrt n}\cos(T\log n)\,e^{-s^2\log^2n/2}+O_s\Bigl(\frac1T\Bigr).
\]
\item The right side uses only the primes. With $s=0.6$ each term with $n>e^{10}$ is below $10^{-9}$, and all of them together are below $10^{-6}$. Figure~\ref{fig:hearzeros} plots it for $5\le T\le50$. Read off
the first zeros. Why does a smaller $s$ separate the zeros better, and what does it cost?
\end{enumerate}
\end{exercise}

\begin{exercise}[Journey]\label{pb:cramersign}
\emph{The coin and the sign of $\pi(x)-\operatorname{li}(x)$.} Let $\pi_{\mathcal C}(x)=\sum_{n\le x}X_n$ be the counting function of Cramér's random primes
(Volume~5, Definition~10.1), and let $\Pi(x)=\sum_{k\ge1}\frac1k\pi(x^{1/k})$ be Riemann's prime-power counting function.
\begin{enumerate}[label=(\alph*)]
\item Show that $\mathbb E\,\pi_{\mathcal C}(x)=\operatorname{li}(x)+O(1)$ and $\operatorname{Var}\pi_{\mathcal C}(x)\sim\frac x{\log x}$, and that
$P\bigl(\pi_{\mathcal C}(x)>\operatorname{li}(x)\bigr)\to\frac12$.
\item Show that, almost surely, $\pi_{\mathcal C}(x)-\operatorname{li}(x)$ changes sign infinitely often. (Use the independent blocks $]n_{j-1},n_j]$ with $n_j=2^{2^j}$, and
both Borel--Cantelli lemmas.)
\item Show that $\pi(x)=\Pi(x)-\frac12\pi(\sqrt x)+O(x^{1/3})$, and that
\[
\Pi(x)-\operatorname{li}(x)=\frac{\psi(x)-x}{\log x}+\int_2^x\frac{\psi(t)-t}{t\log^2t}\dd t+\frac2{\log2}-\operatorname{li}(2).
\]
\item Numerically, one finds, for $10^3\le x\le10^8$: $\pi(x)<\operatorname{li}(x)$ throughout, with $\frac{\pi(x)-\operatorname{li}(x)}{\sqrt x/\log x}$ between $-2.55$ and
$-0.53$; while $\Pi(x)-\operatorname{li}(x)$ changes sign $66\,112$ times below $10^8$, with $\frac{\Pi(x)-\operatorname{li}(x)}{\sqrt x/\log x}$ between $-0.75$ and $0.71$. In the
model, $\pi_{\mathcal C}(x)-\operatorname{li}(x)$ is typically of size $\sqrt{x/\log x}$, which is $2330$ at $x=10^8$. Explain the sign of $\pi(x)-\operatorname{li}(x)$
with (c). Then explain, with (c) and the explicit formula $\psi(x)=x-\sum_\rho\frac{x^\rho}\rho+\cdots$ of Volume~5, Theorem~7.5, why the true fluctuations are of
size $\frac{\sqrt x}{\log x}$ rather than $\sqrt{x/\log x}$ if the Riemann hypothesis holds.
\end{enumerate}
\end{exercise}

\begin{exercise}[Journey]\label{pb:selberg}
\emph{Selberg's formula.} Let $\Lambda_2=\mu*\log^2$, that is, $\Lambda_2(n)=\sum_{d\mid n}\mu(d)\log^2\frac nd$.
\begin{enumerate}[label=(\alph*)]
\item Show that $\Lambda_2(n)=\Lambda(n)\log n+\sum_{dm=n}\Lambda(d)\Lambda(m)$. In particular $\Lambda_2\ge0$, and $\Lambda_2(n)=0$ when $n$ has three distinct prime
factors.
\item Show that $\sum_{d\le x}\frac{\mu(d)}d\log\frac xd=O(1)$ and $\sum_{d\le x}\frac{\mu(d)}d\log^2\frac xd=2\log x+O(1)$.
\item Deduce that $\sum_{n\le x}\Lambda_2(n)=2x\log x+O(x)$.
\item Deduce \emph{Selberg's formula}
\[
\theta(x)\log x+\sum_{pq\le x}\log p\,\log q=2x\log x+O(x),
\]
the sum running over ordered pairs of primes.
\item Let $a=\liminf\frac{\theta(x)}x$ and $A=\limsup\frac{\theta(x)}x$. Show that $a+A=2$.
\item Direct computation gives $\frac1x\bigl(\theta(x)\log x+\sum_{pq\le x}\log p\log q-2x\log x\bigr)$: it is $-3.668$, $-3.667$, $-3.668$, $-3.666$ at $x=10^4$, $10^5$,
$10^6$, $10^7$, although the left side of the formula is only $0.886$ times $2x\log x$ at $10^7$. Explain why so slow an approach to $2x\log x$ is consistent with an
error term $O(x)$.
\end{enumerate}
\end{exercise}

\begin{exercise}[Journey]\label{pb:schur}
\emph{Schur's determinant and the sign of the Gauss sum.} Let $p$ be an odd prime, $\zeta=e(\frac1p)$, $\chi$ the Legendre symbol modulo $p$, and
$G=\sum_{n=0}^{p-1}\zeta^{n^2}$. Let $F$ be the $p\times p$ matrix $(\zeta^{jk})_{0\le j,k<p}$.
\begin{enumerate}[label=(\alph*)]
\item Show that $G=\sum_{n=1}^{p-1}\chi(n)\zeta^n$, that $G^2=\chi(-1)p$, and that $\operatorname{tr}F=G$.
\item Show that $F^2=pR$, where $R$ is the permutation matrix of $j\mapsto-j$ modulo $p$. Deduce that the eigenvalues of $F$ lie in $\{\sqrt p,-\sqrt p,i\sqrt p,-i\sqrt p\}$,
and that their multiplicities $a,b,c,d$ satisfy $a+b+c+d=p$, $a+b-c-d=1$ and $G=\sqrt p\,\bigl(a-b+i(c-d)\bigr)$.
\item Writing $F$ as a Vandermonde matrix, show that $\det F=i^{p(p-1)/2}\,\Delta$ with $\Delta>0$.
\item Comparing with $\det F=\sqrt p^{\,p}\,(-1)^b\,i^{\,c-d}$, show that $G=\sqrt p$ if $p\equiv1\pmod4$ and $G=i\sqrt p$ if $p\equiv3\pmod4$.
\item Let $p\equiv3\pmod4$, $S=\sum_{0<a<p/2}\chi(a)$, $U=\sum_{0<a<p/2}a\chi(a)$ and $T=\sum_{a=1}^{p-1}a\chi(a)$. Show that $T=2U-pS$ and $T=\chi(2)(4U-pS)$, and
using Dirichlet's formula for $L(1,\chi)$ (Volume~5, Theorem~3.10) together with (d), that
\[
S=\bigl(2-\chi(2)\bigr)\frac{\sqrt p}\pi\,L(1,\chi).
\]
Deduce that for $p\equiv3\pmod4$ there are more quadratic residues than non-residues in $]0,\frac p2[$, and that $L(1,\chi)\ge\frac\pi{3\sqrt p}$.
\item Numerically, one finds $(S,L(1,\chi))=(1,\,1.18741)$, $(3,\,1.96520)$, $(3,\,0.246069)$, $(39,\,1.27940)$, $(77,\,2.41818)$ for $p=7$, $23$, $163$, $1019$, $10007$.
Check the formula of (e) on these values, and explain why $p=163$ gives exactly $L(1,\chi)=\frac\pi{\sqrt{163}}$.
\end{enumerate}
\end{exercise}

\begin{exercise}[Journey]\label{pb:euler41}
\emph{Euler's polynomial $n^2+n+41$.} A positive definite binary quadratic form $[a,b,c]=ax^2+bxy+cy^2$ has discriminant $D=b^2-4ac<0$ and $a>0$. Two forms are
\emph{equivalent} if one is obtained from the other by a substitution $(x,y)\mapsto(\alpha x+\beta y,\gamma x+\delta y)$ with $\alpha\delta-\beta\gamma=1$ and integer
coefficients. A form is \emph{reduced} if $|b|\le a\le c$, with $b\ge0$ when $|b|=a$ or $a=c$.
\begin{enumerate}[label=(\alph*)]
\item Show that every positive definite form is equivalent to a reduced form of the same discriminant, that a reduced form has $a\le\sqrt{|D|/3}$, and that equivalent
forms represent the same integers with coprime $(x,y)$.
\item Let $q$ be a prime, $f(n)=n^2+n+q$ and $D=1-4q$. Prove \emph{Rabinowitsch's theorem}: $f(n)$ is prime for $n=0,1,\dots,q-2$ if and only if $[1,1,q]$ is the only
reduced form of discriminant $D$. (If $p<q$ divides $f(n)$, look at $[p,2n+1,\frac{f(n)}p]$.)
\item Find the reduced forms of discriminant $-163$, and deduce Euler's observation. Numerically, one finds that the primes $q\le5000$ with the property of (b) are $2$, $3$,
$5$, $11$, $17$, $41$.
\item Let $\chi(p)=\bigl(\frac{-163}p\bigr)$. Show that for an odd prime $p$ the number of solutions of $f(n)\equiv0\pmod p$, for $q=41$, is $1+\chi(p)$, and that $f(n)$ is
never even. The \emph{Hardy--Littlewood conjecture} predicts that $\#\{n<N:f(n)\text{ prime}\}\sim C\sum_{n<N}\frac1{\log f(n)}$, with
\[
C=2\prod_{p>2}\Bigl(1-\frac{\chi(p)}{p-1}\Bigr).
\]
Explain where $C$ comes from, and show that $C=\frac43\cdot\frac{\sqrt{163}}\pi\prod_{p>2}\frac{1-\chi(p)/(p-1)}{1-\chi(p)/p}$, a product that converges much faster.
Numerically, one finds $C=6.6395$, and $261\,081$ primes among $f(n)$, $n<10^6$, against the prediction $261\,020$.
\item For $n^2+1$ the same heuristic gives the constant $C'=\prod_{p>2}\bigl(1-\frac{\chi_4(p)}{p-1}\bigr)=1.3728$, and a numerical calculation counts $54\,110$ primes among
$n^2+1$, $1\le n<10^6$, against the prediction $53\,972$. Explain in words why $n^2+n+41$ is almost five times richer in primes than $n^2+1$.
\end{enumerate}
\end{exercise}

\begin{exercise}[Journey]\label{pb:hecke}
\emph{The angles of the Gaussian primes.} Every prime $p\equiv1\pmod4$ is $p=a^2+b^2$ with unique integers $a>b>0$ (Volume~1, Theorems~7.3 and~10.4). Let
$\varphi_p=\arctan\frac ba\in\,]0,\frac\pi4[$. \emph{Hecke's theorem} says that the angles $\varphi_p$ are equidistributed in $]0,\frac\pi4[$.
\begin{enumerate}[label=(\alph*)]
\item For $k\in\Z$ and $\alpha\in\Z[i]\setminus\{0\}$ let $\lambda_k(\alpha)=\bigl(\frac\alpha{|\alpha|}\bigr)^{4k}$. Show that $\lambda_k(\alpha)$ depends only on the ideal
$(\alpha)$ and that $\lambda_k(\alpha\beta)=\lambda_k(\alpha)\lambda_k(\beta)$. Let $S_k(x)=\sum_{0<|\alpha|^2\le x}\lambda_k(\alpha)$. Show that $S_k(x)=O\bigl((1+|k|)\sqrt x\bigr)$
for $k\ne0$, by comparing each term with the average of $\lambda_k$ over a unit square.
\item Let $L_k(s)=\frac14\sum_{\alpha\ne0}\lambda_k(\alpha)|\alpha|^{-2s}$. Show that $L_k$ is holomorphic on $\operatorname{Re}s>\frac12$ for $k\ne0$, that
$L_0(s)=\zeta(s)L(s,\chi_4)$, and that for $\operatorname{Re}s>1$
\[
L_k(s)=\frac1{1-(-1)^k2^{-s}}\prod_{p\equiv1\,(4)}\frac1{(1-e^{4ik\varphi_p}p^{-s})(1-e^{-4ik\varphi_p}p^{-s})}\prod_{p\equiv3\,(4)}\frac1{1-p^{-2s}}.
\]
\item Using $3+4\cos\theta+\cos2\theta\ge0$, show that $L_k(1+it)\ne0$ for all real $t$ when $k\ne0$.
\item Let $\vartheta_k(x)=\sum_{p\le x,\ p\equiv1\,(4)}2\cos(4k\varphi_p)\log p$. Using Newman's theorem (Volume~5, Theorem~5.3) and Volume~5, Lemma~5.4, applied to the
non-decreasing functions $\vartheta_0+\vartheta_k$ and $\vartheta_0-\vartheta_k$, show that $\vartheta_k(x)=o(x)$ for $k\ne0$.
\item Deduce Hecke's theorem. (Only cosines appear: reflect the angles $4\varphi_p\in\,]0,\pi[$ to $]-\pi,0[$ and use Weyl's criterion, Volume~4, Theorem~6.13.)
\item Among the $332\,180$ primes $p\equiv1\pmod4$ below $10^7$, one finds numerically a proportion $0.499795$ with $\varphi_p\le\frac\pi8$, and a proportion $0.12658$
with $b<\frac a{10}$. What does Hecke's theorem predict for the second number? Figure~\ref{fig:angles} shows the histogram.
\end{enumerate}
\end{exercise}

\begin{exercise}[Journey]\label{pb:shifted}
\emph{How many prime factors has $p-1$?} Let $\omega(n)$ be the number of distinct prime factors of $n$, and for $y\ge2$ let $\omega_y(n)$ count only those $\le y$.
\begin{enumerate}[label=(\alph*)]
\item Fix $y$. Show that the law of $\omega_y(p-1)$, for $p$ chosen uniformly among the primes $\le x$, tends as $x\to\infty$ to the law of $\sum_{q\le y}B_q$, where the
$B_q$ are independent and $P(B_q=1)=\frac1{q-1}=1-P(B_q=0)$.
\item Show that $\sum_{q\le y}\frac1{q-1}=\log\log y+M+\sum_q\frac1{q(q-1)}+o(1)$, where $M=0.26149\ldots$ is the constant in Mertens's theorem (Volume~5, Theorem~1.6),
so that the constant is $1.03465\ldots$, and that $\frac{\sum_{q\le y}B_q-\log\log y}{\sqrt{\log\log y}}$ tends in law to $\mathcal N(0,1)$ as $y\to\infty$.
\item Using the Brun--Titchmarsh inequality (Volume~5, Exercise~8.4), show that $\sum_{p\le x}\omega(p-1)\le(2+o(1))\,\pi(x)\log\log x$.
\item Assuming the Bombieri--Vinogradov theorem (stated after Volume~5, Proposition~8.11), show that $\frac1{\pi(x)}\sum_{p\le x}\omega(p-1)=\log\log x+O(1)$.
\item Numerically, one finds, over the primes $p\le x$ for $x=10^4,\dots,10^8$: mean of $\omega(p-1)$ equal to $3.027$, $3.298$, $3.518$, $3.696$, $3.847$; mean minus
$\log\log x$ equal to $0.807$, $0.855$, $0.892$, $0.916$, $0.933$; variance $0.540$, $0.656$, $0.781$, $0.906$, $1.021$. Below $10^8$ the proportions of
$\omega(p-1)=2,3,4,5,6,7$ are $0.0835$, $0.2922$, $0.3709$, $0.2041$, $0.0459$, $0.0033$. For comparison, $\omega(n)$ for $n\le10^8$ has mean $3.150$ and variance $1.214$,
while the model of (a) with all $q\le10^8$ has mean $3.948$ and variance $2.573$. Why might one expect the mean to approach $\log\log x+1.035$, and why is the variance
so much smaller than in the model?
\end{enumerate}
\end{exercise}

\begin{exercise}[Journey]\label{pb:romanoff}
\emph{Romanoff's theorem.} Let $r(n)$ be the number of pairs $(p,k)$, $p$ prime and $k\ge1$, with $n=p+2^k$, and let $\mathcal R=\{n:r(n)>0\}$.
\begin{enumerate}[label=(\alph*)]
\item Show that $\sum_{n\le x}r(n)\ge c\,x$ for some $c>0$ and all large $x$, and that $\#(\mathcal R\cap[1,x])\ge\frac{(\sum_{n\le x}r(n))^2}{\sum_{n\le x}r(n)^2}$.
\item Using the upper bound $\#\{p\le x:p+d\text{ prime}\}\le c\,\bigl(\frac d{\varphi(d)}\bigr)^2\frac x{\log^2x}$ for even $d\ge2$, which follows from Selberg's sieve as in
Volume~5, Exercise~8.5, show that $\sum_{n\le x}r(n)^2=O(x)$ provided $\sum_{m\le L}f(m)=O(L)$, where $f(m)=\bigl(\frac{2^m-1}{\varphi(2^m-1)}\bigr)^2$.
\item Show that $\bigl(\frac n{\varphi(n)}\bigr)^2=\sum_{e\mid n}\mu(e)^2g(e)$, where $g$ is multiplicative with $g(q)=\frac{2q-1}{(q-1)^2}\le\frac5q$ for odd primes $q$, and
deduce that $\sum_{m\le L}f(m)\le L\sum_{e\text{ odd}}\frac{\mu(e)^2g(e)}{\operatorname{ord}_e(2)}$.
\item Show that every odd squarefree $e$ with $\operatorname{ord}_e(2)\le J$ divides $P_J=\prod_{j\le J}(2^j-1)<2^{J^2}$, deduce that
$\sum_{\operatorname{ord}_e(2)\le J}\mu(e)^2g(e)=O\bigl((\log J)^5\bigr)$, and conclude \emph{Romanoff's theorem}: $\mathcal R$ has positive lower density.
\item Let $n\equiv7\,629\,217\pmod{11\,184\,810}$. Check that $11\,184\,810=2\cdot3\cdot5\cdot7\cdot13\cdot17\cdot241$ and that for every $k\ge1$ the number $n-2^k$ is
divisible by one of $3,5,7,13,17,241$. (Look at $k$ modulo $2,4,3,12,8,24$.) Deduce that infinitely many odd numbers are not in $\mathcal R$.
\item Numerically, one finds that the proportion of integers $n\le x$ in $\mathcal R$ is $0.475$, $0.4662$, $0.4605$, $0.4578$ for $x=10^4,\dots,10^7$, and among odd $n$
it is $0.947$, $0.932$, $0.921$, $0.916$, while $\frac1L\sum_{m\le L}f(m)=2.155$, $2.372$, $2.487$, $2.535$ for $L=10$, $20$, $40$, $80$. Comment.
\end{enumerate}
\end{exercise}

\begin{figure}[ht]
\centering
\fig{ch14-landau}
\caption{The zeros hear the primes: $-\frac{2\pi\sqrt x}T\sum_{0<\gamma\le T}\cos(\gamma\log x)$ for $1.5\le x\le30$, with the first $2000$ zeros
($T=\gamma_{2000}$). By Problem~\ref{pb:landauzeros}(d) it tends to $\Lambda(x)$: spikes of height $\log p$ (red dots) at the prime powers $p^k$, and nothing in
between.}\label{fig:landau}
\end{figure}

\begin{figure}[ht]
\centering
\fig{ch14-hearzeros}
\caption{The primes hear the zeros: the right side of Problem~\ref{pb:hearzeros}(c) with $s=0.6$, computed from the prime powers $n\le e^{10}$ alone (blue). Its peaks
sit at the ordinates of the zeros of $\zeta$ (red ticks: $14.13$, $21.02$, $25.01$, $30.42$, $32.94$, $37.59$, $40.92$, $43.33$, $48.01$, $49.77$).}\label{fig:hearzeros}
\end{figure}

\begin{figure}[ht]
\centering
\fig{ch14-angles}
\caption{Angles $\varphi_p=\arctan\frac ba$ of the $332\,180$ primes $p=a^2+b^2\equiv1\pmod4$ below $10^7$, in $20$ bins, each bar divided by the count expected for the
uniform distribution. All bars lie between $0.991$ and $1.007$.}\label{fig:angles}
\end{figure}

There is an $x_0$ such that
\begin{equation}\label{eq:bt13}
\pi(x;d,1)\le\frac{5x}{\varphi(d)\log x}\quad(x\ge x_0,\ 1\le d\le\sqrt x).
\end{equation}

\section{Hints}

\begin{hint}{ex:sigma9}
$E_4E_6=E_{10}$. Divide by $-24$.
\end{hint}
\begin{hint}{ex:e2}
If it were modular, Theorem~\ref{thm:dimension}(a) would apply.
\end{hint}
\begin{hint}{ex:zeros}
$S$ fixes $i$; and $\rho$ is fixed by $ST$, or note $-\frac1\rho=\rho+1$.
\end{hint}
\begin{hint}{ex:tausmall}
$\tau(4)=\tau(2)^2-2^{11}$; $\tau(8)=\tau(2)\tau(4)-2^{11}\tau(2)$.
\end{hint}
\begin{hint}{ex:eisensteinhecke}
Use the $q$-expansion of $T_p$ and the multiplicativity of $\sigma_{k-1}$; for commutation, compute the coefficient of $q^n$ in $T_pT_\ell f$, which is symmetric in $p$ and $\ell$.
\end{hint}
\begin{hint}{ex:heckebound}
(a) $\operatorname{Im}(\gamma\tau)=\frac{\operatorname{Im}\tau}{|c\tau+d|^2}$. On $F$, $\Delta=q\,u(q)$ with $u$ bounded, so $\phi\le(\operatorname{Im}\tau)^6e^{-2\pi\operatorname{Im}\tau}\sup|u|$. (b) $|\Delta(x+iy)|\le Cy^{-6}$.
\end{hint}

\begin{hint}{pb:singer}
(a) $i-j=d$ with $i,j\in D$ means that the point $g^j$ lies in $U\cap g^{-d}U$. Why are these two planes different?
\end{hint}

\begin{hint}{pb:consecsqfree}
(a) For the primes $p\le y$ the condition is periodic modulo $\prod_{p\le y}p^2$; the primes $p>y$ exclude at most $2N/y+O(\sqrt N)$
integers. (c) Chinese remainders.
\end{hint}

\begin{hint}{pb:pairwise}
(a) Locally at $p$: at most one of the $k$ integers is divisible by $p$. (b) Compare the two local factors.
\end{hint}

\begin{hint}{pb:twosqfree}
Inclusion--exclusion over the primes up to $13$, which leaves an error of at most $4^6$ terms, and a tail estimate for the
primes from $17$ on.
\end{hint}

\begin{hint}{pb:dsquared}
(a) Compare Euler factors. (b) $d^2=\tau_4*g$ with $g$ supported on squares, and Volume~5, Exercise~2.7.
\end{hint}

\begin{hint}{pb:energy}
(a) $E_+(A)=\sum_sr(s)^2$ with $\sum_sr(s)=|A|^2$. (b) Write $a/c$ in lowest terms.
\end{hint}

\begin{hint}{pb:kleitman}
(a) Induction on $n$, splitting according to whether $n\in S$, and Chebyshev's inequality for two terms. (b) Replace each family by
the increasing family it generates.
\end{hint}

\begin{hint}{pb:chebyshevsum}
$\sum_{i,j}(a_i-a_j)(b_i-b_j)\ge0$; for (b), take an independent copy of $X$.
\end{hint}

\begin{hint}{pb:primitive}
(a) $a$ is made of the smallest primes of $am$. (c) Mertens: $\prod_{p<y}(1-\frac1p)\ge\frac c{\log y}$.
\end{hint}

\begin{hint}{pb:divisorbound}
$d$ is multiplicative and $d(p^e)=e+1$. Compare $\frac{e+1}{p^{e\alpha}}$ with its value at $e+1$.
\end{hint}

\begin{hint}{pb:logtwo}
(a) $\frac{\mathrm d^n}{\mathrm dx^n}x^{n+k}=\frac{(n+k)!}{k!}x^k$, and integrate by parts $n$ times. (b) $\int_0^1\frac{x^k}{1+x}\dd x$ is $\pm\log2$ plus a rational with
denominator dividing $d_k$. (c) Maximise $\frac{x(1-x)}{1+x}$; bound $d_n$ with Chebyshev's $\prod_{p\le x}p<4^x$.
\end{hint}

\begin{hint}{pb:doubleint}
Expand $\frac1{1-xy}=\sum_k(xy)^k$ and integrate term by term; for (b) differentiate $\int\!\!\int x^{r+t}y^{s+t}/(1-xy)$ in $t$, or use
$\int_0^1x^k(-\log x)\dd x=\frac1{(k+1)^2}$.
\end{hint}

\begin{hint}{pb:beukerstwo}
(a) Problem~\ref{pb:doubleint}(a). (b) For fixed $y$, $\frac{\partial^n}{\partial x^n}\frac1{1-xy}=\frac{n!\,y^n}{(1-xy)^{n+1}}$. (c) $(1-x)(1-y)\le(1-\sqrt{xy})^2$.
(d) $d_n\le3^n$ for large $n$.
\end{hint}

\begin{hint}{pb:beukersthree}
(c) After the substitution, $\frac{(xyz)^n\dd z}{(1-(1-xy)z)^{n+1}}$ becomes $\frac{(1-w)^n\dd w}{1-(1-xy)w}$. (d) First maximise in $w$ for fixed $xy$; the
answer is $(1+\sqrt{xy})^{-2}$. (e) $27(\sqrt2-1)^4<1$.
\end{hint}

\begin{hint}{pb:fracint}
Substitute $x=\frac1t$ and cut $[1,\infty[$ into the intervals $[k,k+1[$. For (b) you will need Stirling's formula.
\end{hint}

\begin{hint}{pb:primefrac}
The function $\{1/t\}$ is Riemann integrable on $[\frac1K,1]$; sandwich it between step functions and count primes in each interval
$]xt_{j-1},xt_j]$ by the prime number theorem. The primes below $x/K$ are few.
\end{hint}

\begin{hint}{pb:leastpf}
(a) Abel summation (Volume~5, Lemma~1.1): $\sum_{p\le x}p=x\pi(x)-\int_2^x\pi(t)\dd t$. (b) A composite $n$ has $P^-(n)\le\sqrt n$. (c) $p_n\sim n\log n$. (d) Substitute $s=t^2$ in
$\int\frac t{\log t}\dd t$, and use $\pi=\operatorname{Li}+E$ with Volume~5, Corollary~7.9.
\end{hint}

\begin{hint}{pb:largestpf}
(a) Write $n=pm$ with $p=P^+(n)$. (c) $\lfloor\frac xp\rfloor$ is the number of $k\ge1$ with $p\le\frac xk$. Then let $x\to\infty$ in each term $\frac{\log x}{x^2}\sum_{p\le x/k}p$ and dominate them by $\frac C{k^2}$.
(d) $A(n)-P^+(n)$ is a sum of at most $\log_2n$ primes, each at most the second largest prime factor of $n$.
\end{hint}

\begin{hint}{pb:factorialdiv}
(c) Telescope: $\log(1+m)=\sum_{k=1}^m\log(1+\frac1k)$ with $m=\lfloor\frac np\rfloor$. Exchange the sums, and dominate as in Problem~\ref{pb:largestpf}.
\end{hint}

\begin{hint}{pb:centralbinom}
(a) $\lfloor2t\rfloor-2\lfloor t\rfloor\in\{0,1\}$. (c) For $k\le n^{1/4}$ the intervals are long; use Brun--Titchmarsh (Volume~5, Theorem~8.4) to dominate. The primes below
$n^{3/4}$ are negligible.
\end{hint}

\begin{hint}{pb:superprimes}
(a) $p_m\sim m\log m$ with $m=p_n$. (b) Compare with $\sum\frac1{n\log^2n}$.
\end{hint}

\begin{hint}{pb:ramprimes}
(a) $f$ jumps up by $1$ at each prime and down by $1$ at each $2p$. (b) Compare $n=f(R_n)\sim\frac{R_n}{2\log R_n}$ with $2n=\pi(p_{2n})\sim\frac{p_{2n}}{\log p_{2n}}$.
\end{hint}

\begin{hint}{pb:pminusone}
(a) Split $q\le x^{1/4}$, where~\eqref{eq:bt13} applies to $d=q^2$, and $q>x^{1/4}$. (b) For fixed $y$, the condition ``$q^2\nmid p-1$ for all $q\le y$'' depends only on $p$
modulo $\prod_{q\le y}q^2$. (c) $\frac{\varphi(m)}m=\sum_{d\mid m}\frac{\mu(d)}d$; treat $d\le y$, $y<d\le\sqrt x$ and $d>\sqrt x$ separately. For the last range put $e=\frac{p-1}d$.
\end{hint}

\begin{hint}{pb:titchmarsh}
(b) $\tau(m)\le2\#\{d\mid m:d\le\sqrt m\}$. (c) $\frac m{\varphi(m)}=\sum_{d\mid m}\frac{\mu(d)^2}{\varphi(d)}$, and as in Problem~\ref{pb:pminusone}; you need $\varphi(d)\ge c_\varepsilon d^{1-\varepsilon}$,
which follows from $\frac d{\varphi(d)}\le\tau(d)$ and Volume~5, Exercise~2.6.
\end{hint}

\begin{hint}{pb:menon}
(a) $\mathbb F_p^\times$ is cyclic of order $p-1$. (b) Burnside's lemma: the orbit of $x$ is determined by $\gcd(x,n)$; the fixed points of $a$ are the $x$ with $(a-1)x\equiv0$.
(c) $\gcd(n,p-1)$ depends only on $p\bmod n$. (d) $\Phi_d$ has $\varphi(d)$ roots modulo $p\nmid n$ if $d\mid p-1$, and none otherwise.
\end{hint}

\begin{hint}{pb:omegapoly}
(b) Exchange: $\sum_n\omega(n^2+1)=\sum_p\#\{n\le x:p\mid n^2+1\}$. (c) $\sum_p\frac{\chi(p)}p$ converges (Volume~5, Corollary~4.5). (d) Work with $\omega_y(m)$, the number of prime
factors $p\le y=x^{1/4}$ of $m$, and compute $\sum_n\omega_y(n^2+1)^2$ by the Chinese remainder theorem.
\end{hint}

\begin{hint}{pb:rankin}
(a) For a $y$-smooth $n\le x$, $1\le(\frac xn)^\sigma$. (b) Split $\log\prod(1-p^{-\sigma})^{-1}=\sum_pp^{-\sigma}+\sum_p\sum_{k\ge2}\frac{p^{-k\sigma}}k$, and note $\sigma\ge\frac23$.
(d) Compare $\prod_{p\le y}(1-\frac1p)^{-1}\sim e^\gamma\log y$ with the true size.
\end{hint}

\begin{hint}{pb:renyi}
(b) Compute the Euler factor $1+\sum_{k\ge2}\frac{h_z(p^k)}{p^k}$. (c) Cauchy's formula on $|z|=1$ and dominated convergence. (d) The factor $p=2$ of $R$ has a pole at $z=2$; the rest
is analytic in $|z|<3$.
\end{hint}

\begin{hint}{pb:twosquares}
(a) Compare Euler factors, using Volume~1, Theorem~7.3. (b) For $n\le y$, $\frac1n\le e\,n^{-1-\delta}$; for the lower bound, the tail $n>y$ is at most
$y^{-\delta/2}F(1+\frac\delta2)$. (c) Write $\log n=\sum_{p^k\parallel n}k\log p$; removing $p^k$ from $n$ keeps $b=1$. For $B(x)$, compare $\sum b(n)\log n$ with
$\log\frac x{\log x}\,\bigl(B(x)-B(\frac x{\log x})\bigr)$. (d) $m\cdot p$ with $p\equiv1\pmod4$ is a sum of two squares when $m$ is. (e) Abel's theorem for Dirichlet series: if $B(x)\sim\frac{Kx}{\sqrt{\log x}}$,
compute $F(\sigma)$ by partial summation.
\end{hint}

\begin{hint}{pb:ramanujanE2}
(a) Logarithmic derivative of the product. (b) Differentiate $\log\eta(-\frac1\tau)=\log\eta(\tau)+\frac12\log(-i\tau)$. (c) $f$ is holomorphic at $i\infty$ with $f=-\frac1{12}+\dots$;
use the dimensions $\dim M_4=\dim M_6=\dim M_8=1$ (Volume~5, Theorem~12.9) for $f$, for $DE_4-\frac{E_2E_4}3$ and for $DE_6-\frac{E_2E_6}2$. (d) Compare coefficients of $q^n$ in
$12DE_2=E_2^2-E_4$.
\end{hint}

\begin{hint}{pb:landauzeros}
(a) Expand $\frac{\zeta'}\zeta(2+it)=-\sum\Lambda(n)n^{-2-it}$ and integrate term by term. (b) Take the logarithmic derivative of the functional equation, and integrate by
parts against $x^{it}$. (c) Integrate $x^s\frac{\zeta'}\zeta(s)$ around the rectangle $-1\le\sigma\le2$, $1\le t\le T'$ with a good height $T'$ (Volume~5, Lemma~7.3). (e) Weyl's criterion
with $\alpha\gamma_n$; $\sum_{\gamma\le T}e^{2\pi ik\alpha\gamma}$ is the sum of (d) with $x=e^{2\pi k\alpha}$, and $N(T)\sim\frac T{2\pi}\log T$.
\end{hint}

\begin{hint}{pb:hearzeros}
(a) $h$ extends to an entire even function of $t$. (b) The rectangle encloses the zeros with $|\gamma|<T'$ and the pole at $1$; on $\operatorname{Re}s=-1$ use the
logarithmic derivative of the functional equation, and shift the $\Gamma$ term to the line $\frac12$, picking up the pole of $\frac{\Gamma'}\Gamma(\frac s2)$ at $s=0$. (c) Under the
Riemann hypothesis $H(\rho)=h(\gamma)$. Use $\operatorname{Re}\frac{\Gamma'}\Gamma(\frac14+\frac{it}2)=\log\frac{|t|}2+O(t^{-2})$.
\end{hint}

\begin{hint}{pb:cramersign}
(a) $\sum_{3\le n\le x}\frac1{\log n}-\operatorname{li}(x)$ is bounded; the Lindeberg condition is trivial for bounded variables. (b) With $V_j=\operatorname{Var}\pi_{\mathcal C}(n_j)$, the
block increment exceeds $\operatorname{li}$'s increment by $2\sqrt{V_j}$ with probability bounded below, while $|\pi_{\mathcal C}(n_{j-1})-\operatorname{li}(n_{j-1})|\le\sqrt{V_j}$
eventually. (c) $\psi(x)=\sum_k\theta(x^{1/k})$; for the identity, integrate $\frac{d\psi}{\log t}$ by parts. (d) $\psi(x)-x$ is about $\sqrt x$ times a bounded
oscillating sum, and $\pi(\sqrt x)\approx\frac{2\sqrt x}{\log x}$.
\end{hint}

\begin{hint}{pb:selberg}
(a) Differentiate $\Lambda=\mu*\log$ with the derivation $f\mapsto f\log$, which satisfies $(f*g)\log=(f\log)*g+f*(g\log)$. (b) Write $\log\frac xd$ and $\log^2\frac xd$ through
$\sum_{m\le x/d}$ of simple functions, and use $\sum_{d\mid n}\mu(d)=[n=1]$ and $\sum_{m\le y}\frac{\tau(m)}m=\frac12\log^2y+2\gamma\log y+c+O(y^{-1/2}\log y)$. (c) $\sum_{n\le x}\Lambda_2(n)=\sum_d\mu(d)\sum_{m\le x/d}\log^2m$.
(e) Apply Selberg's formula at $x$ and at $\frac xp$, and average.
\end{hint}

\begin{hint}{pb:schur}
(b) $(F^2)_{jk}=\sum_l\zeta^{l(j+k)}$. $\operatorname{tr}F^2=p$. (c) $\zeta^k-\zeta^j=2i\,e\bigl(\frac{j+k}{2p}\bigr)\sin\frac{\pi(k-j)}p$. (d) $(a-b)^2+(c-d)^2=1$ by (a); go through the
four cases. (e) Pair $a$ with $p-a$, and split $T$ by the parity of $a$.
\end{hint}

\begin{hint}{pb:euler41}
(a) Repeatedly replace $[a,b,c]$ by $[a,b-2ma,\dots]$ (translation) or by $[c,-b,a]$ (rotation); $a$ decreases. (b) For $\Leftarrow$, a reduced form $[a,2n+1,c]$ with $a>1$ gives
$ac=f(n)$. For $\Rightarrow$, the principal form does not represent $p<q$. (d) Count roots of $(2n+1)^2\equiv-163\pmod p$. Dirichlet's formula for $L(1,\chi)$ gives the
product $\prod(1-\chi(p)/p)^{-1}$.
\end{hint}

\begin{hint}{pb:hecke}
(a) A unit $i^j$ has $(i^j)^{4k}=1$. Across a unit square $\lambda_k$ varies by $O(\frac{|k|}{|\alpha|})$, and $\int_0^{2\pi}e^{4ik\theta}\dd\theta=0$. (c) Look at
$|L_0(\sigma)^3L_k(\sigma+it)^4L_{2k}(\sigma+2it)|\ge1$. (d) The Laplace transform of $\vartheta_k$ is a combination of $\frac{L_k'}{L_k}$ and $\frac{L_{-k}'}{L_{-k}}$.
\end{hint}

\begin{hint}{pb:shifted}
(a) $p\equiv1\pmod q$ has density $\frac1{q-1}$ (Volume~5, Theorem~5.8), and conditions for distinct $q$ combine by the Chinese remainder theorem. (c) At most one prime
$q>\sqrt x$ divides $p-1$; split at $x^\varepsilon$. (e) At most one large prime factor.
\end{hint}

\begin{hint}{pb:romanoff}
(a) $\pi(\frac x2)\cdot\lfloor\log_2\frac x2\rfloor$ pairs give $n\le x$. Cauchy--Schwarz. (b) Off the diagonal, $p'=p+2^l(2^m-1)$. (c) $1+g(q)=\bigl(\frac q{q-1}\bigr)^2$.
$e\mid2^m-1$ exactly when $\operatorname{ord}_e(2)\mid m$. (d) $P_J$ has fewer than $J^2$ prime factors; Mertens. (e) The orders of $2$ modulo $3,5,7,13,17,241$ are
$2,4,3,12,8,24$.
\end{hint}

\section{Solutions}

\begin{solution}{ex:sigma9}
\idea $M_{10}$ is a line, so the two forms $E_4E_6$ and $E_{10}$ with constant term $1$ must be equal.
\writeup Both have constant term $1$, and $\dim M_{10}=1$, so $E_4E_6=E_{10}$. The coefficient of $q^n$ gives
$-264\sigma_9(n)=240\sigma_3(n)-504\sigma_5(n)-240\cdot504\sum_{m=1}^{n-1}\sigma_3(m)\sigma_5(n-m)$. Divide by $-24$: $11\sigma_9(n)=-10\sigma_3(n)+21\sigma_5(n)+5040\sum\sigma_3(m)\sigma_5(n-m)$. For $n=2$: $11\cdot513=5643$ and
$21\cdot33-10\cdot9+5040=693-90+5040=5643$.
\end{solution}

\begin{solution}{ex:e2}
\idea The proof of Theorem~\ref{thm:dimension} kills every modular form of weight $2$, and $E_2$ has constant term $1$, so it cannot be one. (It fails only by a little: $E_2(-\frac1\tau)=\tau^2E_2(\tau)+\frac{6\tau}{\pi i}$.)
\writeup Since $\sigma_1(n)\le n^2$, the series converges uniformly on $|q|\le r<1$, so $E_2$ is holomorphic in $q$, hence in $\tau$, and is $1$-periodic. If it were a modular form of weight $2$, Theorem~\ref{thm:dimension}(a) would give $E_2=0$,
but its constant term is $1$.
\end{solution}

\begin{solution}{ex:zeros}
\idea A fixed point of a symmetry is a place where the automorphy factor must equal $1$, unless the form vanishes.
\writeup $S(i)=i$, so $E_6(i)=E_6(-\frac1i)=i^6E_6(i)=-E_6(i)$, hence $E_6(i)=0$. For $\rho=e^{2\pi i/3}=-\frac12+\frac{\sqrt3}2i$, $\rho^2+\rho+1=0$ gives $-\frac1\rho=\rho+1$. So
$E_4(\rho)=E_4(\rho+1)=E_4(-\frac1\rho)=\rho^4E_4(\rho)=\rho E_4(\rho)$, and $\rho\ne1$ forces $E_4(\rho)=0$. Then $1728\Delta(i)=E_4(i)^3-E_6(i)^2=E_4(i)^3$; as $\Delta(i)\ne0$, $E_4(i)\ne0$.
\end{solution}

\begin{solution}{ex:tausmall}
\idea Mordell's relations generate all values from the values at primes.
\writeup $\tau(4)=576-2048=-1472$; $\tau(6)=-24\cdot252=-6048$; $\tau(8)=\tau(2)\tau(4)-2^{11}\tau(2)=35328+49152=84480$; $\tau(12)=\tau(4)\tau(3)=-370944$. These match the product. For the congruence:
$\sigma_{11}(2)=2049$ and $-24-2049=-2073=-3\cdot691$; $\sigma_{11}(3)=177148$ and $177148-252=176896=256\cdot691$.
\end{solution}

\begin{solution}{ex:eisensteinhecke}
\idea Hecke operators act on $q$-expansions by a rule that only involves the multiplicative structure of the indices; divisor sums are already multiplicative, which makes $E_k$ an eigenform.
\writeup Let $s(n)=\sigma_{k-1}(n)$, with $s(n/p)=0$ if $p\nmid n$. We claim $s(pn)+p^{k-1}s(n/p)=(1+p^{k-1})s(n)$. Write $n=p^rm$ with $p\nmid m$. The divisors of $p^{r+1}m$ are those of $p^rm$ and the $p^{r+1}e$ with $e\mid m$, so
$s(pn)=s(n)+p^{(r+1)(k-1)}s(m)$; similarly, if $r\ge1$, $p^{k-1}s(n/p)=p^{k-1}\bigl(s(n)-p^{r(k-1)}s(m)\bigr)$, and if $r=0$ both $p^{k-1}s(n/p)$ and $p^{k-1}(s(n)-p^{r(k-1)}s(m))$ are $0$. Adding gives the claim. For $n=0$
the constant terms give $(1+p^{k-1})a_0$. So $T_pE_k=(1+p^{k-1})E_k$.

For commutation, write $T_pf=\sum b_nq^n$ with $b_n=a_{pn}+p^{k-1}a_{n/p}$. The coefficient of $q^n$ in $T_\ell T_pf$ is
\[
\begin{aligned}b_{\ell n}+\ell^{k-1}b_{n/\ell}&=a_{p\ell n}+p^{k-1}a_{\ell n/p}\eqbr+\ell^{k-1}a_{pn/\ell}+(p\ell)^{k-1}a_{n/(p\ell)},\end{aligned}
\]
with the convention $a_x=0$ when $x$ is not an integer, which is consistent because $p\ne\ell$: for instance $\frac{pn}\ell$ is an integer exactly when $\ell\mid n$, that is, exactly when $b_{n/\ell}$ is defined.
This expression is symmetric in $p$ and $\ell$, so $T_\ell T_p=T_pT_\ell$.
\end{solution}

\begin{solution}{ex:heckebound}
\idea The weight is exactly what makes $(\operatorname{Im}\tau)^{k/2}|f|$ invariant. An invariant function is controlled by its values on $F$, where a cusp form decays exponentially. Then the Fourier coefficient is an
integral along a horizontal line, and the best line is at height $\frac1n$.
\writeup (a) $\phi(\gamma\tau)=\frac{(\operatorname{Im}\tau)^6}{|c\tau+d|^{12}}|c\tau+d|^{12}|\Delta(\tau)|=\phi(\tau)$. By Theorem~\ref{thm:fundamental}, $\sup_{\mathbb H}\phi=\sup_F\phi$. On $F$, $|q|\le e^{-\pi\sqrt3}$ and $\Delta=q\,u(q)$ with
$u$ continuous on the closed disc $|q|\le e^{-\pi\sqrt3}$, so $\phi(\tau)\le\max|u|\cdot y^6e^{-2\pi y}$ with $y=\operatorname{Im}\tau$, which is bounded. So $\phi\le C$ on $\mathbb H$.

(b) For fixed $y>0$, the series
\[
\Delta(x+iy)=\sum_m\tau(m)e^{-2\pi my}e^{2\pi imx}
\]
converges uniformly in $x$, so integrating term by term over $[0,1]$,
\[
\tau(n)e^{-2\pi ny}=\int_0^1\Delta(x+iy)e^{-2\pi inx}\dd x.
\]
Hence $|\tau(n)|\le e^{2\pi ny}\max_x|\Delta(x+iy)|$, which is at most $e^{2\pi ny}Cy^{-6}$. With $y=\frac1n$: $|\tau(n)|\le e^{2\pi}Cn^6$.
\end{solution}

\begin{solution}{pb:singer}
\idea Multiplication by $g$ permutes the points of the projective plane $\mathbb P(\mathbb F_{q^3})$ cyclically. Two lines meet in
one point, and a line is never mapped to itself by a non-trivial power of $g$: that is the whole proof.

\writeup (a) $g^i$ and $g^{i+N}$ differ by the factor $g^N$, which has order $q-1$ and so generates $\mathbb F_q^*$; hence
$g^i\in U\iff g^{i+N}\in U$ and $D$ is well defined. The map $i\mapsto\mathbb F_qg^i$ is a bijection from $\Z/N\Z$ to the
$\frac{q^3-1}{q-1}=N$ points (lines through $0$), and $U$ contains $\frac{q^2-1}{q-1}=q+1$ of them: $|D|=q+1$.

Now fix $d\not\equiv0\pmod N$. The pairs $(i,j)\in D^2$ with $i-j=d$ correspond to the $j$ with $g^j\in U$ and $g^{j+d}\in U$, that is
to the points of $U\cap g^{-d}U$. Let $c=g^d$; then $c\notin\mathbb F_q$, because $\mathbb F_q^*=\langle g^N\rangle$. So
$\mathbb F_q(c)$ is a subfield of $\mathbb F_{q^3}$ strictly bigger than $\mathbb F_q$, and since its degree divides $3$ it is
$\mathbb F_{q^3}$. If $cU=U$, then $U$ would be stable under $\mathbb F_q[c]=\mathbb F_{q^3}$, a sub-$\mathbb F_{q^3}$-space of dimension
$2$ over $\mathbb F_q$, impossible. So $g^{-d}U$ and $U$ are distinct $2$-dimensional subspaces of a $3$-dimensional space; their
intersection has dimension $2+2-3=1$, a single point. So $d$ is a difference in exactly one way.

(b) Two points $x\ne y$ lie on $D+t$ when $x-t,y-t\in D$; their difference $x-y$ is a difference of $D$ in exactly one way, which
determines $x-t$, hence $t$. Two translates $D+t$ and $D+t'$, $t\ne t'$, meet at the points $x$ with $x-t$ and $x-t'$ in $D$; the
difference $t'-t$ is represented once, so there is one such $x$. (The translates are exactly the images $g^{-t}U$, the lines of the
plane.)

(c) Modulo $7$, the differences of $\{0,1,3\}$ are $\pm1,\pm2,\pm3$. Modulo $13$, those of $\{0,1,3,9\}$ are
\[
\pm1,\ \pm3,\ \pm9,\ \pm2,\ \pm8,\ \pm6,
\]
which are $\pm1,\dots,\pm6$ since $9\equiv-4$ and $8\equiv-5$. Modulo $21$, the differences of $\{0,1,4,14,16\}$ are
\[
\pm1,\ \pm4,\ \pm14,\ \pm16,\ \pm3,\ \pm13,\ \pm15,\ \pm10,\ \pm12,\ \pm2,
\]
and with $14\equiv-7$, $16\equiv-5$, $13\equiv-8$, $15\equiv-6$, $12\equiv-9$ these are $\pm1,\dots,\pm10$.

\tools Volume~2, Chapter~9 (subfields of $\mathbb F_{q^3}$, cyclicity of $\mathbb F_q^*$); Volume~0, Chapter~7 (dimension of an
intersection).
\leads Every known projective plane of order $n$ has $n$ a prime power, and every known cyclic one comes from Singer's construction;
the prime power conjecture for cyclic planes has been checked up to orders in the billions. A numerical calculation rebuilds Singer sets for
$q=2,3,5,7$ from an explicit $\mathbb F_{q^3}$.
\end{solution}

\begin{solution}{pb:consecsqfree}
\idea Squarefreeness is a condition at every prime. For finitely many primes it is periodic and its density is a finite
product; for the large primes, being divisible by $p^2$ is so rare that the total effect vanishes.

\writeup (a) Fix $y\ge2$, let $Q=\prod_{p\le y}p^2$, and let $S_y$ be the set of $n$ such that neither $n$ nor $n+1$ is divisible by
$p^2$ for any $p\le y$. Membership in $S_y$ depends only on $n\bmod Q$, and by the Chinese remainder theorem the number of good
residues modulo $Q$ is $\prod_{p\le y}(p^2-2)$: modulo $p^2$ we must avoid $n\equiv0$ and $n\equiv-1$. So
$\#(S_y\cap[1,N])=N\prod_{p\le y}(1-\frac2{p^2})+O(Q)$.

Let $G$ be the set of $n$ with $n$ and $n+1$ squarefree. Then $G\subset S_y$, and an $n\le N$ in $S_y\setminus G$ has $p^2\mid n$ or
$p^2\mid n+1$ for some prime $p>y$ with $p^2\le N+1$. Their number is at most
\[
\sum_{y<p\le\sqrt{N+1}}2\Bigl(\frac N{p^2}+1\Bigr)\le\frac{2N}y+2\sqrt{N+1},
\]
using $\sum_{p>y}\frac1{p^2}\le\sum_{m>y}\frac1{m(m-1)}\le\frac1y$. Dividing by $N$ and letting $N\to\infty$,
\[
\prod_{p\le y}\Bigl(1-\frac2{p^2}\Bigr)-\frac2y\le\liminf\frac{\#(G\cap[1,N])}N\le\limsup\frac{\#(G\cap[1,N])}N\le\prod_{p\le y}\Bigl(1-\frac2{p^2}\Bigr).
\]
Now let $y\to\infty$; the product converges since $\sum\frac2{p^2}<\infty$.

(b) The same argument, with three forbidden residues $0,-1,-2$ modulo each $p^2$, gives $\prod_p(1-\frac3{p^2})$; for $p=2$ the
three residues $0,3,2$ modulo $4$ are distinct, so the factor is $\frac14$, and the product is not $0$. Among four consecutive
integers one is divisible by $4$.

(c) Given $L$, choose distinct primes $p_1,\dots,p_L$ and, by the Chinese remainder theorem, $n$ with $n\equiv-j\pmod{p_j^2}$ for
$j=1,\dots,L$. Then $p_j^2\mid n+j$, and none of $n+1,\dots,n+L$ is squarefree.

\tools Volume~2, Chapter~2 (Chinese remainders); Volume~4, Exercise~10.2 (the density $\frac6{\pi^2}$ of squarefree numbers).
\leads By (c) the first run of $L$ non-squarefree integers starts below $\prod_{j\le L}p_j^2$, and in practice much
earlier: one finds numerically that a run of $4$ first starts at $242$ ($242=2\cdot11^2$, $243=3^5$, $244=4\cdot61$, $245=5\cdot7^2$).
How long the longest run below $x$ is remains open; a heuristic count (each of the $L$ integers needs its own prime square,
of size about $L$) suggests about $\frac{\log x}{2\log\log x}$.
\end{solution}

\begin{solution}{pb:pairwise}
\idea At each prime the $k$ integers behave like independent coins of bias $\frac1p$, and ``pairwise coprime at $p$'' means at
most one coin shows heads. The probability is the product of these local probabilities, by the same truncation as in
Problem~\ref{pb:consecsqfree}.

\writeup (a) For a prime $p$, the probability that at most one of $k$ independent residues modulo $p$ is $0$ is
$(1-\frac1p)^k+k\frac1p(1-\frac1p)^{k-1}=(1-\frac1p)^{k-1}(1+\frac{k-1}p)$. Fix $y$ and let $Q=\prod_{p\le y}p$. The condition ``no
prime $p\le y$ divides two of the $a_i$'' depends only on the $a_i$ modulo $Q$, and by Chinese remainders its density in
$\{1,\dots,N\}^k$ is $c_y=\prod_{p\le y}(1-\frac1p)^{k-1}(1+\frac{k-1}p)$, up to $O(Q/N)$. The $k$-tuples satisfying it but not
pairwise coprime have two entries divisible by a common prime $p>y$; there are at most
$\binom k2\sum_{p>y}N^{k-2}\lfloor N/p\rfloor^2\le\binom k2N^k/y$ of them. Let $N\to\infty$, then $y\to\infty$, as before.

(b) The local factor in (a) for $k=3$ is $(1-\frac1p)^2(1+\frac2p)=1-\frac3{p^2}+\frac2{p^3}$, whereas $(1-\frac1{p^2})^3=1-\frac3{p^2}+\frac3{p^4}-\frac1{p^6}$.
The difference is $\frac2{p^3}-\frac3{p^4}+\frac1{p^6}>0$. The three events are positively correlated at every prime: if $p$
divides $a$ and $b$, the events $p\mid\gcd(b,c)$ and $p\mid\gcd(a,c)$ happen together, so the bad events overlap and the good
event is more likely than independence predicts.

\tools Volume~2, Chapter~2; Volume~3, Chapter~1 (inclusion--exclusion, independence); Volume~5, Exercise~2.3 (the gcd of $k$ integers).
\leads For $k=4$ the probability is $0.1149\ldots$ As $k\to\infty$ it tends to $0$ very fast, because each small prime
$p$ is then very likely to divide at least two of the $k$ integers. The same local reasoning gives the probability that
the $\binom k2$ differences of $k$ random integers are all coprime to $m$.
\end{solution}

\begin{solution}{pb:twosqfree}
\idea The expected number of representations $n=a+(n-a)$ with both parts squarefree is about $0.32\,n$. To turn this into a
proof one needs an \emph{explicit} error, so the sieve uses only the primes up to $13$, where inclusion--exclusion has few terms,
and a crude tail bound for the others.

\writeup Let $A(n)$ be the number of $a\in\{1,\dots,n-1\}$ with $a$ and $n-a$ squarefree, $P=2\cdot3\cdot5\cdot7\cdot11\cdot13$, and
$B(n)$ the number of $a$ in the same range such that $p^2\nmid a$ and $p^2\nmid n-a$ for every $p\mid P$. By inclusion--exclusion,
\[
B(n)=\sum_{d\mid P}\sum_{e\mid P}\mu(d)\mu(e)\,\#\{1\le a\le n-1:\ d^2\mid a,\ e^2\mid n-a\}.
\]
The two congruences $a\equiv0\pmod{d^2}$, $a\equiv n\pmod{e^2}$ are compatible if and only if $g^2\mid n$, where $g=\gcd(d,e)$, and
then they define one class modulo $L=\operatorname{lcm}(d^2,e^2)$, which contains $\frac{n-1}L+\theta$ elements of $[1,n-1]$, with
$|\theta|<1$. So $B(n)=(n-1)\Pi(n)+E$ with $|E|\le4^6=4096$ (the number of pairs $(d,e)$), and $\Pi(n)$ the multiplicative sum
\[
\Pi(n)=\sum_{d,e\mid P}\frac{\mu(d)\mu(e)[g^2\mid n]}L=\prod_{p\mid P}\Bigl(1-\frac2{p^2}+\frac{[p^2\mid n]}{p^2}\Bigr)\ge\prod_{p\le13}\Bigl(1-\frac2{p^2}\Bigr)=0.33350\ldots
\]
(the sum factors over primes because the summand is a product of local factors, each coming from $d,e\in\{1,p\}$).

An $a$ counted by $B(n)$ but not by $A(n)$ has $p^2\mid a$ or $p^2\mid n-a$ for some prime $17\le p\le\sqrt n$; there are at most
$2\sum_{p\ge17}\frac n{p^2}<0.0331\,n$ such $a$, since $\sum_{p\ge17}p^{-2}=0.016546\ldots$ (this is the prime zeta value
$\sum_pp^{-2}=0.452247\ldots$ minus its first six terms). Hence
\[
A(n)\ge0.3335(n-1)-4096-0.0331\,n>0.30\,n-4097,
\]
which is positive for $n\ge13657$. For $2\le n<14000$ one finds numerically a representation directly (for instance $2=1+1$, $3=1+2$),
and in fact the least $a$ that works is always very small.

\tools Volume~1, Chapter~4 (Möbius function, multiplicativity); Volume~2, Chapter~2 (Chinese remainders); Problem~\ref{pb:consecsqfree}.
\leads Estermann (1931) showed that the number of representations is $\sim c(n)\,n$ with an explicit singular series
$c(n)=\prod_p(1-\frac2{p^2}+\frac{[p^2\mid n]}{p^2})$, the exact analogue of the Hardy--Littlewood formula for Goldbach's problem,
with squarefree numbers in place of primes. For primes the same conjectural formula is far beyond reach, because primes are
not defined by local conditions alone.
\end{solution}

\begin{solution}{pb:dsquared}
\idea Euler products turn the identity into one about power series in a single prime, and the asymptotics follow from the known
average of $\tau_4$ by a convolution with a sparse function.

\writeup (a) Both sides are Euler products of multiplicative functions, absolutely convergent for $s>1$. With $x=p^{-s}$, the
factor at $p$ on the left is $\sum_k(k+1)^2x^k=\frac{1+x}{(1-x)^3}$ (differentiate $\sum x^{k+1}$ twice), and on the right it is
$\frac{1-x^2}{(1-x)^4}=\frac{1+x}{(1-x)^3}$.

(b) By (a), $d^2=\tau_4*g$, where $\tau_4=\mathbf1^{*4}$ has Dirichlet series $\zeta^4$ and $g(m^2)=\mu(m)$, $g=0$ off the squares (its
series is $\frac1{\zeta(2s)}$). So $\sum_{n\le x}d(n)^2=\sum_{m\le\sqrt x}\mu(m)T(x/m^2)$ with $T(y)=\sum_{n\le y}\tau_4(n)$. By Volume~5,
Exercise~2.7, $T(y)=\frac16y\log^3y+O(y\log^2y)$ for $y\ge3$, and trivially $T(y)=O(y)$ for $1\le y<3$. Hence
\[
\sum_{n\le x}d(n)^2=\frac x6\sum_{m\le\sqrt x}\frac{\mu(m)}{m^2}\Bigl(\log\frac x{m^2}\Bigr)^3+O\Bigl(\sum_{m\le\sqrt x}\frac x{m^2}\log^2x\Bigr).
\]
The error is $O(x\log^2x)$. Expanding $(\log x-2\log m)^3=\log^3x+O(\log^2x\,(1+\log m)^3)$ and using $\sum_m\frac{(1+\log m)^3}{m^2}<\infty$,
the main term is $\frac x6\log^3x\sum_{m\le\sqrt x}\frac{\mu(m)}{m^2}+O(x\log^2x)$. Finally
$\sum_{m\le\sqrt x}\frac{\mu(m)}{m^2}=\frac1{\zeta(2)}+O(x^{-1/2})=\frac6{\pi^2}+O(x^{-1/2})$, and $\frac16\cdot\frac6{\pi^2}=\frac1{\pi^2}$.

\tools Volume~5, Exercise~2.7 (average of $\tau_k$); Volume~1, Chapter~4 (Möbius, multiplicative functions); Volume~4, Chapter~10
($\zeta(2)=\frac{\pi^2}6$).
\leads Ramanujan stated (a) in 1916, together with similar identities such as $\sum\sigma_a(n)\sigma_b(n)n^{-s}=\frac{\zeta(s)\zeta(s-a)\zeta(s-b)\zeta(s-a-b)}{\zeta(2s-a-b)}$.
The true asymptotic is $xP(\log x)$ with $P$ a cubic, and the convergence of the leading term is slow: one finds numerically
$\frac{\sum d(n)^2}{x\log^3x}\approx0.17$ at $x=10^5$, against the limit $0.101$.
\end{solution}

\begin{solution}{pb:energy}
\idea The energy counts collisions of sums. If there are few distinct sums, collisions are forced; Cauchy--Schwarz turns this into a
bound. For products of integers up to $N$, the solutions of $ab=cd$ can be listed explicitly.

\writeup (a) Let $r(s)$ be the number of pairs $(a,b)\in A^2$ with $a+b=s$. Then $\sum_sr(s)=|A|^2$ and $\sum_sr(s)^2=E_+(A)$. By
Cauchy--Schwarz over the $|A+A|$ values of $s$, $|A|^4\le|A+A|\,E_+(A)$. The multiplicative case is identical.

(b) Let $ab=cd$ with $a,b,c,d\in[1,N]$. Let $g=\gcd(a,c)$, $a=gu$, $c=gv$ with $u,v$ coprime. Then $ub=vd$, so $v\mid b$; write $b=lv$, and then
$d=lu$. Conversely every choice of coprime $u,v$ and integers $g,l\ge1$ gives a solution $(gu,lv,gv,lu)$, which lies in $[1,N]^4$ if and only
if $g,l\le\frac N{\max(u,v)}$; and $(u,v,g,l)$ is recovered from the solution as $g=\gcd(a,c)$. This proves the formula. For $M\ge2$ there
are $2\varphi(M)$ coprime pairs with $\max(u,v)=M$, and one pair, $(1,1)$, for $M=1$. Since $\lfloor N/M\rfloor^2=\frac{N^2}{M^2}+O(\frac NM)$,
\[
\begin{aligned}
E_\times&=N^2+2N^2\sum_{2\le M\le N}\frac{\varphi(M)}{M^2}\\
&\qquad+O\Bigl(N\sum_{M\le N}\frac{\varphi(M)}M\Bigr)\\
&=2N^2\sum_{M\le N}\frac{\varphi(M)}{M^2}+O(N^2).
\end{aligned}
\]
By partial summation from $\Phi(x)=\sum_{M\le x}\varphi(M)$, which is $\frac3{\pi^2}x^2+O(x\log x)$ (Volume~5, Theorem~2.6),
\[
\begin{aligned}
\sum_{M\le N}\frac{\varphi(M)}{M^2}&=\frac{\Phi(N)}{N^2}+2\int_1^N\frac{\Phi(x)}{x^3}dx\\
&=\frac6{\pi^2}\log N+O(1),
\end{aligned}
\]
so $E_\times=\frac{12}{\pi^2}N^2\log N+O(N^2)$.

(c) By (a) and (b), $|A\cdot A|\ge\frac{N^4}{E_\times}=\frac{\pi^2}{12}\frac{N^2}{\log N}(1+o(1))$.

\tools Volume~0, Theorem~4.4 (Cauchy--Schwarz); Volume~5, Theorem~2.6 (the mean of $\varphi$); Volume~1, Chapter~1 (gcd).
\leads Volume~5 shows that the table has $o(N^2)$ entries; the truth, found by Ford in 2008, is of order
$\frac{N^2}{(\log N)^{\delta}(\log\log N)^{3/2}}$ with $\delta=1-\frac{1+\log\log2}{\log2}=0.086\ldots$, so the energy bound is off only by a power of
$\log N$. Direct computation gives both.
\end{solution}

\begin{solution}{pb:kleitman}
\idea Split the cube into two halves according to the last coordinate. Each half of an increasing family is increasing, and the
upper half contains the lower one; Chebyshev's inequality for two numbers glues the halves.

\writeup (a) Induction on $n$; for $n=0$ the families are $\varnothing$ or $\{\varnothing\}$ and the claim is clear. For $n\ge1$ and an increasing $\mathcal A$, let
$\mathcal A_0=\{S\subseteq[n-1]:S\in\mathcal A\}$ and $\mathcal A_1=\{S\subseteq[n-1]:S\cup\{n\}\in\mathcal A\}$. Both are increasing families on $[n-1]$, and $\mathcal A_0\subseteq\mathcal A_1$. Do
the same for $\mathcal B$. Then, by induction,
\[
\begin{aligned}
|\mathcal A\cap\mathcal B|&=|\mathcal A_0\cap\mathcal B_0|+|\mathcal A_1\cap\mathcal B_1|\\
&\ge\frac{|\mathcal A_0||\mathcal B_0|+|\mathcal A_1||\mathcal B_1|}{2^{n-1}}\\
&\ge\frac{(|\mathcal A_0|+|\mathcal A_1|)(|\mathcal B_0|+|\mathcal B_1|)}{2^n}\\
&=\frac{|\mathcal A||\mathcal B|}{2^n},
\end{aligned}
\]
where the last inequality is $2(xy+XY)-(x+X)(y+Y)=(X-x)(Y-y)\ge0$ with $x=|\mathcal A_0|\le X=|\mathcal A_1|$ and $y\le Y$ likewise. For decreasing
families apply this to the increasing families $\{[n]\setminus S:S\in\mathcal A\}$.

(b) If $\mathcal F$ is intersecting, so is its up-closure $\{T:T\supseteq S\text{ for some }S\in\mathcal F\}$, which contains $\mathcal F$; so we may assume the
families $\mathcal F_1,\dots,\mathcal F_k$ increasing. An intersecting family contains at most one of $S$ and $[n]\setminus S$, so $|\mathcal F_j|\le2^{n-1}$, and its complement
$\mathcal G_j$ in $2^{[n]}$ is decreasing with $|\mathcal G_j|\ge2^{n-1}$. By (a) applied repeatedly (an intersection of decreasing families is decreasing),
$|\mathcal G_1\cap\dots\cap\mathcal G_k|\ge2^n\prod_j\frac{|\mathcal G_j|}{2^n}\ge2^{n-k}$. The union of the $\mathcal F_j$ is the complement of this intersection. The families
$\mathcal F_j=\{S:j\in S\}$, $j\le k$, are intersecting, and their union misses exactly the $2^{n-k}$ sets disjoint from $[k]$.

\tools Volume~0, Chapter~8 (counting); Problem~\ref{pb:chebyshevsum} (the two-term case).
\leads (a) is the simplest case of the FKG inequality (Fortuin, Kasteleyn and Ginibre, 1971), central in statistical physics: in a
random graph, increasing properties such as ``connected'' and ``contains a triangle'' are positively correlated. Kleitman's (b)
was the answer to a question of Erdős.
\end{solution}

\begin{solution}{pb:chebyshevsum}
\idea Two similarly ordered sequences go up together, so every product $(a_i-a_j)(b_i-b_j)$ is non-negative; summing them is exactly the
inequality.

\writeup (a) $\sum_{i,j}(a_i-a_j)(b_i-b_j)\ge0$ since each term is a product of two numbers of the same sign. Expanding,
$2n\sum_ia_ib_i-2\sum_ia_i\sum_jb_j\ge0$.

(b) Let $Y$ be an independent copy of $X$. Then $(f(X)-f(Y))(g(X)-g(Y))\ge0$, and taking expectations,
$2\mathbb E[fg(X)]-2\mathbb E[f(X)]\mathbb E[g(X)]\ge0$.

(c) $M_k=\mathbb E[R^k]$ for a random variable $R$ with values in $[0,1]$: the limiting law of $\varphi(n)/n$ (Problem~\ref{pb:phiovern}). The functions
$x^j$ and $x^k$ are non-decreasing on $[0,1]$, so by (b) $M_{j+k}\ge M_jM_k$. Numerically $M_2=0.4282\ge M_1^2=0.3696$.

\tools Volume~3, Chapter~5 (expectation, independence); Problem~\ref{pb:phiovern}.
\leads (b) is the one-dimensional FKG inequality; Problem~\ref{pb:kleitman} is its version on the cube. Chebyshev used (a) in
1882 for integrals; with $f=g$ it is just $\operatorname{Var}f(X)\ge0$.
\end{solution}

\begin{solution}{pb:primitive}
\idea Attach to each $a$ the integers whose ``beginning'', in the factorisation sorted by primes, is exactly $a$. Primitivity makes these
sets disjoint, so their densities add up to at most $1$, and by Mertens the density attached to $a$ is about $\frac1{a\log a}$ or more.

\writeup (a) For $N=am\in M_a$, all primes of $m$ are $\ge P(a)$ while all primes of $a$ are $\le P(a)$. So $a$ is obtained from the prime
factorisation of $N$, listed in increasing order with multiplicity, by keeping an initial segment (all primes $<P(a)$ with their full
multiplicity, and some copies of $P(a)$). Two initial segments of the same list are comparable for divisibility. So if $N\in M_a\cap M_{a'}$
then $a\mid a'$ or $a'\mid a$; for a primitive set, $a=a'$.

(b) $N\in M_a$ if and only if $a\mid N$ and $N/a$ has no prime factor $<P(a)$. The second condition is periodic modulo $Q=\prod_{p<P(a)}p$ and holds
for a proportion $\prod_{p<P(a)}(1-\frac1p)$ of residues. So $M_a$ has density $\frac1a\prod_{p<P(a)}(1-\frac1p)$.

(c) By (a), for any finite subset $A'$ of $A$ the union of the $M_a$, $a\in A'$, has density $\sum_{a\in A'}d(M_a)\le1$ (densities of finitely many
disjoint periodic sets add). By Mertens's theorem (Volume~5, Theorem~1.6), there is $c>0$ with $\prod_{p<y}(1-\frac1p)\ge\frac c{\log y}$ for all $y\ge2$
(the asymptotic formula handles large $y$, and finitely many small $y$ are positive). As $P(a)\le a$,
\[
\sum_{a\in A'}\frac c{a\log a}\le\sum_{a\in A'}\frac1a\prod_{p<P(a)}\Bigl(1-\frac1p\Bigr)\le1,
\]
so $\sum_{a\in A}\frac1{a\log a}\le\frac1c$.

\tools Volume~5, Theorem~1.6 (Mertens); Volume~2, Chapter~2 (periodic sets and Chinese remainders); Volume~1, Chapter~2 (unique
factorisation).
\leads For the primes, $\sum_p\frac1{p\log p}=1.6366\ldots$, and Erdős conjectured in 1988 that no primitive set does better. Jared Duker
Lichtman proved it in 2022. Without the $\log a$, a primitive subset of $[1,x]$ can have $\sum\frac1a$ as large as $c\frac{\log x}{\sqrt{\log\log x}}$
(Besicovitch, Behrend).
\end{solution}

\begin{solution}{pb:divisorbound}
\idea $d(n)/n^\alpha$ is multiplicative, so its maximum is a product of local maxima, one for each prime; only the primes below
$2^{1/\alpha}$ can do better than $e=0$, and each local factor is unimodal in $e$.

\writeup (a) For $n=\prod p^{e_p}$, $\frac{d(n)}{n^\alpha}=\prod_pf_p(e_p)$ with $f_p(e)=\frac{e+1}{p^{e\alpha}}$, $f_p(0)=1$. For $p^\alpha\ge2$, $f_p(e)\le\frac{e+1}{2^e}\le1$,
so the maximum over $e$ is $1$; for each of the finitely many $p<2^{1/\alpha}$ the maximum over $e$ exists since $f_p(e)\to0$. Choosing each $e_p$ at
a maximiser shows that the product of local maxima is attained, and no $n$ exceeds it.

(b) The smallest maximiser $n_\alpha$ takes, for each $p$, the smallest maximiser of $f_p$. Now
\[
\frac{f_p(e+1)}{f_p(e)}=\frac{e+2}{e+1}\,p^{-\alpha},
\]
and $\frac{e+2}{e+1}$ decreases with $e$. So $f_p$ increases strictly as long as $\frac{e+2}{e+1}>p^\alpha$ and then never increases again: its
smallest maximiser is the least $e$ with $\frac{e+2}{e+1}\le p^\alpha$. When $p$ grows, $p^\alpha$ grows, the condition becomes easier and
this least $e$ can only decrease: $e_2\ge e_3\ge e_5\ge\cdots$. Finally $e_p\ge1$ if and only if the condition fails at $e=0$, that is $2>p^\alpha$.

(c) $\alpha=\frac13$: the primes with $p^{1/3}<2$ are those below $8$. The thresholds $\frac{e+2}{e+1}$ are $2,\frac32,\frac43,\frac54,\dots$, and
$2^{1/3}=1.260$, $3^{1/3}=1.442$, $5^{1/3}=1.710$, $7^{1/3}=1.913$. So $e_2=3$ (first threshold below $1.260$ is $\frac54$), $e_3=2$, $e_5=e_7=1$, and
$n_{1/3}=2^3\cdot3^2\cdot5\cdot7=2520$, with $d=48$ and $\frac{48}{2520^{1/3}}=3.527\ldots<3.53$.
For $\alpha=\frac14$ the primes with $p<16$ count, with $2^{1/4}=1.189$, $3^{1/4}=1.316$, $5^{1/4}=1.495$, $7^{1/4}=1.627$ and
$11^{1/4},13^{1/4}\in\left]1.5,2\right[$: so $e_2=5$ (the threshold $\frac76=1.167$), $e_3=3$ (the threshold $\frac54$), $e_5=2$ (the
threshold $\frac43$), $e_7=e_{11}=e_{13}=1$ (the threshold $\frac32$), and $n_{1/4}=2^5\cdot3^3\cdot5^2\cdot7\cdot11\cdot13=21\,621\,600$, with
$d(n_{1/4})=6\cdot4\cdot3\cdot2\cdot2\cdot2=576$ and $\frac{576}{n_{1/4}^{1/4}}=8.446\ldots$

\tools Volume~1, Chapter~4 (multiplicative functions).
\leads Wigert (1907) found the true order of the maximum: $\log d(n)\le(\log2+o(1))\frac{\log n}{\log\log n}$, with equality infinitely often.
Ramanujan studied the record holders, the \emph{highly composite numbers}, such as $12$, $2520$ and $5040$, in a long paper of 1915.
\end{solution}

\begin{solution}{pb:logtwo}
\idea A polynomial with integer coefficients, divided by $1+x$, integrates to a rational number plus a multiple of $\log2$, and
the denominators of the rationals are controlled by $d_n$. Rodrigues's formula $P_n=\frac1{n!}D^n(x^n(1-x)^n)$ makes the integral small:
the $n$ derivatives can be thrown onto $\frac1{1+x}$.

\writeup (a) $x^n(1-x)^n=\sum_k(-1)^k\binom nkx^{n+k}$, and $\frac1{n!}D^nx^{n+k}=\binom{n+k}nx^k$, so $P_n(x)=\sum_k(-1)^k\binom nk\binom{n+k}kx^k$ has integer
coefficients. The polynomial $x^n(1-x)^n$ has zeros of order $n$ at $0$ and $1$, so its derivatives of order $<n$ vanish there and $n$
integrations by parts give $\int_0^1P_nf=\frac{(-1)^n}{n!}\int_0^1x^n(1-x)^nf^{(n)}$ for every smooth $f$. With $f(x)=\frac1{1+x}$,
$f^{(n)}(x)=\frac{(-1)^nn!}{(1+x)^{n+1}}$, which gives the second formula.

(b) Dividing $x^k$ by $1+x$,
\[
\frac{x^k}{1+x}=\sum_{j=1}^k(-1)^{j-1}x^{k-j}+\frac{(-1)^k}{1+x},
\]
so
\[
\int_0^1\frac{x^k}{1+x}\dd x=(-1)^k\log2+\sum_{j=1}^k\frac{(-1)^{j-1}}{k-j+1},
\]
and $d_k$, hence $d_n$ for $k\le n$, times the last sum is an integer. Summing against the integer coefficients of $P_n$ gives
$d_nI_n=a_n+b_n\log2$ with $b_n=d_nP_n(-1)$ and $a_n\in\Z$.

(c) Let $g(x)=\frac{x(1-x)}{1+x}$. Then $g'(x)=\frac{1-2x-x^2}{(1+x)^2}$ vanishes on $[0,1]$ only at $x=\sqrt2-1$, where $g=(\sqrt2-1)^2=3-2\sqrt2$, its
maximum. So $0<I_n=\int_0^1\frac{g(x)^n}{1+x}\dd x\le(3-2\sqrt2)^n\log2$. Next, $d_n=\prod_{p\le n}p^{\lfloor\log n/\log p\rfloor}$; the primes $p>\sqrt n$ appear
to the first power, and the others contribute at most $n$ each, so by Volume~1, Theorem~2.4(i), $d_n\le4^nn^{\sqrt n}$. Suppose
$\log2=\frac pq$. Then $q(a_n+b_n\log2)=qa_n+pb_n$ is a non-zero integer (as $I_n>0$), and
\[
1\le|qd_nI_n|\le q\,n^{\sqrt n}\bigl(4(3-2\sqrt2)\bigr)^n\log2.
\]
Since $4(3-2\sqrt2)=0.686\ldots<1$ and $n^{\sqrt n}=e^{\sqrt n\log n}$ grows more slowly than any exponential, the right side tends to $0$:
a contradiction.

\tools{Volume~1, Theorem~2.4 (Chebyshev's bound); Volume~0, Chapter~6 (integration by parts); Volume~4, Chapter~4.}
\leads The same integrals give an \emph{irrationality measure}: $|\log2-\frac pq|>q^{-\mu}$ for large $q$. With $d_n=e^{n(1+o(1))}$ and
$|P_n(-1)|=(3+2\sqrt2)^{n(1+o(1))}$, the argument gives $\mu=1+\frac{1+\log(3+2\sqrt2)}{\log(3+2\sqrt2)-1}=4.62\ldots$; refinements of the integrals (Rukhadze, Hata, Marcovecchio) brought $\mu$ down to $3.57\ldots$ (R.~Marcovecchio,
2009). For comparison, almost every real number has $\mu=2$.
\end{solution}

\begin{solution}{pb:doubleint}
\idea Expand $\frac1{1-xy}$ as a geometric series; each term is a product of two one-variable integrals, and the sums telescope.

\writeup (a) The terms of $\sum_k(xy)^kx^ry^s$ are non-negative, so we may integrate term by term (Volume~4, Chapter~4):
\[
\int_0^1\!\!\int_0^1\frac{x^ry^s}{1-xy}\dd x\dd y=\sum_{k\ge0}\frac1{(r+k+1)(s+k+1)}.
\]
For $r>s$,
\[
\frac1{(r+k+1)(s+k+1)}=\frac1{r-s}\Bigl(\frac1{s+k+1}-\frac1{r+k+1}\Bigr),
\]
and the sum telescopes to $\frac1{r-s}\sum_{j=s+1}^r\frac1j$. For
$r=s$ the sum is $\sum_{j>r}\frac1{j^2}=\zeta(2)-\sum_{j\le r}\frac1{j^2}$.

(b) We have $-\log(xy)=-\log x-\log y$, and integration by parts gives
\[
\int_0^1x^a(-\log x)\dd x=\frac1{(a+1)^2}.
\]
So the term $(xy)^kx^ry^s$ contributes $\frac1{A^2B}+\frac1{AB^2}$, where $A=r+k+1$ and $B=s+k+1$. For $r>s$ this equals
$\frac1{r-s}\bigl(\frac1{B^2}-\frac1{A^2}\bigr)$, because $\frac1{B^2}-\frac1{A^2}=\frac{(A-B)(A+B)}{A^2B^2}$ and $A-B=r-s$. The sum telescopes again. For $r=s$ the term is
$\frac2{(r+k+1)^3}$, and the sum is $2\sum_{j>r}\frac1{j^3}$.

(c) $\frac1{1+xy}=\sum_k(-xy)^k$, and the partial sums lie between $0$ and $1$, so term-by-term integration is legitimate (dominated
convergence, Volume~4, Chapter~4). The result is $\sum_{k\ge0}\frac{(-1)^k}{(k+1)^2}=\zeta(2)-2\cdot\frac{\zeta(2)}4=\frac{\pi^2}{12}$, by Volume~4, Theorem~3.7.

\tools{Volume~4, Chapter~4 (term-by-term integration) and Theorem~3.7 ($\zeta(2)=\frac{\pi^2}6$).}
\leads The denominators are the point: for $r,s\le n$ the rational parts in (a) become integers when multiplied by $d_n^2$, and those
in (b) when multiplied by $d_n^3$. This is what Problems~\ref{pb:beukerstwo} and~\ref{pb:beukersthree} use.
\end{solution}

\begin{solution}{pb:beukerstwo}
\idea Problem~\ref{pb:doubleint} shows that any polynomial in $x$ and $y$ of degree at most $n$ in each, integrated against $\frac1{1-xy}$,
gives a rational number with denominator $d_n^2$ plus a rational multiple of $\zeta(2)$. Choosing the polynomial $(1-y)^nP_n(x)$ makes
the integral small, by the same trick as for $\log2$.

\writeup (a) Expand $(1-y)^nP_n(x)=\sum_{r,s\le n}p_rq_sx^ry^s$ with integers $p_r,q_s$. By Problem~\ref{pb:doubleint}(a), the term with $r\ne s$
contributes $\frac{p_rq_s}{|r-s|}\sum_{j=\min+1}^{\max}\frac1j$; here $|r-s|$ divides $d_n$ and $d_n\sum\frac1j\in\Z$, so $d_n^2$ times it is an
integer. The term with $r=s$ contributes $p_rq_r\bigl(\zeta(2)-\sum_{j\le r}\frac1{j^2}\bigr)$, and $d_n^2\sum_{j\le r}\frac1{j^2}\in\Z$. So
$d_n^2J_n=a_n+b_n\zeta(2)$ with $b_n=d_n^2\sum_rp_rq_r$ and $a_n\in\Z$.

(b) For fixed $y\in[0,1[$, the function $f(x)=\frac1{1-xy}$ is smooth on $[0,1]$ with $f^{(n)}(x)=\frac{n!\,y^n}{(1-xy)^{n+1}}$. As in
Problem~\ref{pb:logtwo}(a),
\[
\int_0^1\frac{P_n(x)}{1-xy}\dd x=\frac{(-1)^n}{n!}\int_0^1x^n(1-x)^n\frac{n!\,y^n}{(1-xy)^{n+1}}\dd x.
\]
Multiply by $(1-y)^n$ and integrate over $y\in[0,1[$.

(c) Let $t=\sqrt{xy}$. By AM--GM, $x+y\ge2t$, so $(1-x)(1-y)=1-(x+y)+xy\le(1-t)^2$. Hence
\[
\frac{x(1-x)y(1-y)}{1-xy}\le\frac{t^2(1-t)^2}{1-t^2}=\frac{t^2(1-t)}{1+t}=h(t).
\]
On $]0,1[$, $(\log h)'=\frac2t-\frac1{1-t}-\frac1{1+t}$ vanishes when $2(1-t^2)=2t$, that is $t^2+t-1=0$, $t=\tau=\frac{\sqrt5-1}2$, and $h$
is maximal there. Since $\tau^2=1-\tau$ and $\tau(1+\tau)=1$, $h(\tau)=\tau^2\cdot\tau^2\cdot\tau=\tau^5$. Therefore, by (b),
\[
0<|J_n|\le\tau^{5n}\int_0^1\!\!\int_0^1\frac{\mathrm dx\dd y}{1-xy}=\zeta(2)\tau^{5n}.
\]

(d) By the prime number theorem (Volume~5, Theorem~5.5), $\log d_n=\psi(n)\sim n$, so $d_n\le3^n$ for $n$ large. If $\zeta(2)=\frac pq$, then
$q d_n^2J_n=qa_n+pb_n$ is a non-zero integer with $|qd_n^2J_n|\le q\zeta(2)(9\tau^5)^n$, and $9\tau^5=0.81\ldots<1$: a contradiction for large $n$.

\tools{Problems~\ref{pb:logtwo} and~\ref{pb:doubleint}; Volume~0, Theorem~4.3 (AM--GM); Volume~5, Theorem~5.5 (prime number theorem).}
\leads The numbers $b_n$ are, up to sign and the factor $d_n^2$, the integers $\sum_k\binom nk^2\binom{n+k}k=1,3,19,147,1251,\dots$ that Apéry also
used for $\zeta(2)$; they satisfy a three-term recurrence, checked numerically. Volume~4 proved the irrationality of
$\pi^2$ by Niven's method, which says nothing about $\zeta(3)$; Beukers's method does.
\end{solution}

\begin{solution}{pb:beukersthree}
\idea Same plan as for $\zeta(2)$, one dimension up. The kernel $\frac{-\log(xy)}{1-xy}$ produces $\zeta(3)$ with denominator $d_n^3$
(Problem~\ref{pb:doubleint}(b)). Writing the logarithm as an integral in a third variable gives room for a change of variables that
lets us integrate by parts in both $x$ and $y$.

\writeup (a) As in Problem~\ref{pb:beukerstwo}(a), using Problem~\ref{pb:doubleint}(b): for $r\ne s$, $|r-s|$ divides $d_n$ and $d_n^2\sum\frac1{j^2}\in\Z$; for
$r=s$, $d_n^3\sum_{j\le r}\frac1{j^3}\in\Z$.

(b) $\int_0^1\frac{\mathrm dz}{1-(1-xy)z}=\Bigl[\frac{-\log(1-(1-xy)z)}{1-xy}\Bigr]_0^1=\frac{-\log(xy)}{1-xy}$.

(c) So $K_n=\int\!\!\int\!\!\int\frac{P_n(x)P_n(y)}{1-(1-xy)z}\dd x\dd y\dd z$ (the integrand is positive up to the bounded factor $P_n(x)P_n(y)$, and
$\int\!\!\int\frac{-\log(xy)}{1-xy}=2\zeta(3)$ is finite, so the order of integration may be changed). For fixed $y,z$, $f(x)=\frac1{1-z+xyz}$ has
$f^{(n)}=\frac{(-1)^nn!\,(yz)^n}{(1-(1-xy)z)^{n+1}}$, and $n$ integrations by parts in $x$ give
\[
K_n=\int\!\!\int\!\!\int\frac{(xyz)^n(1-x)^nP_n(y)}{(1-(1-xy)z)^{n+1}}\dd x\dd y\dd z.
\]
Now fix $x,y$ and put $w=\frac{1-z}{1-(1-xy)z}$, a decreasing bijection of $[0,1]$. Then $z=\frac{1-w}{1-(1-xy)w}$,
$1-(1-xy)z=\frac{xy}{1-(1-xy)w}$ and $\mathrm dz=-\frac{xy\dd w}{(1-(1-xy)w)^2}$, so
\[
\frac{(xyz)^n\dd z}{(1-(1-xy)z)^{n+1}}=-\frac{(1-w)^n\dd w}{1-(1-xy)w}
\]
(the powers of $xy$ and of $1-(1-xy)w$ cancel). Hence
\[
K_n=\int\!\!\int\!\!\int\frac{(1-x)^n(1-w)^nP_n(y)}{1-(1-xy)w}\dd x\dd y\dd w.
\] For fixed $x,w$,
$g(y)=\frac1{1-w+xyw}$ has $g^{(n)}=\frac{(-1)^nn!\,(xw)^n}{(1-(1-xy)w)^{n+1}}$, and $n$ integrations by parts in $y$ give the formula.

(d) Fix $x,y\in]0,1]$ and let $\rho=\sqrt{xy}$, $s=1-xy=1-\rho^2$. The function $w\mapsto\frac{w(1-w)}{1-sw}$ has derivative of the sign of
$(1-2w)(1-sw)+sw(1-w)=1-2w+sw^2$, which vanishes at $w=\frac{1-\rho}{s}=\frac1{1+\rho}$; there $1-w=\frac\rho{1+\rho}$, $1-sw=\rho$, and the value
is $\frac1{(1+\rho)^2}$, the maximum. As in Problem~\ref{pb:beukerstwo}(c), $x(1-x)y(1-y)\le\rho^2(1-\rho)^2$, so the whole expression is at most
\[
\Bigl(\frac{\rho(1-\rho)}{1+\rho}\Bigr)^2\le(3-2\sqrt2)^2=(\sqrt2-1)^4,
\]
by the maximum of $\frac{t(1-t)}{1+t}$ found in Problem~\ref{pb:logtwo}(c). By (c) with $n=0$, the integral of $\frac1{1-(1-xy)w}$ over the cube is
$K_0=\int\!\!\int\frac{-\log(xy)}{1-xy}=2\zeta(3)$ (Problem~\ref{pb:doubleint}(b)). So $0<K_n\le2\zeta(3)(\sqrt2-1)^{4n}$.

(e) For large $n$, $d_n\le3^n$, and $27(\sqrt2-1)^4=0.79\ldots<1$. If $\zeta(3)=\frac pq$, the non-zero integers $qd_n^3K_n$ would tend to $0$.

\tools{Problems~\ref{pb:logtwo}, \ref{pb:doubleint} and~\ref{pb:beukerstwo}; Volume~5, Theorem~5.5; Volume~4, Chapter~4 (changing the order of integration).}
\leads The coefficient $b_n$ is $2d_n^3$ times the Apéry number $\sum_k\binom nk^2\binom{n+k}k^2=1,5,73,1445,\dots$, which satisfies
$n^3u_n=(34n^3-51n^2+27n-5)u_{n-1}-(n-1)^3u_{n-2}$; Apéry's original proof ran on this recurrence. The method gives an irrationality
measure $13.41\ldots$ for $\zeta(3)$, later improved to $5.51\ldots$ by Rhin and Viola (2001). For $\zeta(5)$ the analogous integrals do
not decay fast enough against $d_n^5$, and Rivoal's and Zudilin's theorems use cleverer integrals of the same family.
\end{solution}

\begin{solution}{pb:fracint}
\idea After $x=\frac1t$, the integral runs over $[1,\infty[$ with weight $\frac1{t^2}$, and on each $[k,k+1[$ the fractional part is a
polynomial. The sums that come out are partial sums of the harmonic series and of $\log k$, so $\gamma$ and Stirling's constant appear.

\writeup (a) $\int_0^1\{\frac1x\}\dd x=\int_1^\infty\frac{\{t\}}{t^2}\dd t$, and
\[
\int_k^{k+1}\frac{t-k}{t^2}\dd t=\log\frac{k+1}k-\frac1{k+1}.
\]
Summing for $k\le K$ gives $\log(K+1)-(H_{K+1}-1)\to1-\gamma$ (Volume~4, Corollary~11.6).

(b) Similarly
\[
\int_k^{k+1}\frac{(t-k)^2}{t^2}\dd t=1-2k\log\frac{k+1}k+\frac k{k+1}.
\]
Now $\sum_{k\le K}\frac k{k+1}=K+1-H_{K+1}$ and, by summation by parts,
$\sum_{k\le K}k\log\frac{k+1}k=K\log(K+1)-\log K!$. So the integral over $[1,K+1]$ is
\[
2K+1-H_{K+1}-2K\log(K+1)+2\log K!.
\]
By Stirling (Volume~4, Theorem~5.8), $2\log K!=2K\log K-2K+\log(2\pi K)+o(1)$, while $2K\log(K+1)=2K\log K+2+o(1)$ and $H_{K+1}=\log K+\gamma+o(1)$.
Everything cancels except $1-2-\gamma+\log(2\pi)$: the limit is $\log(2\pi)-\gamma-1=0.2606\ldots$

(c) On $]\frac1{k+1},\frac1k]$, $\lfloor\frac1x\rfloor=k$, so the integral is $\sum_k\frac1k\bigl(\frac1k-\frac1{k+1}\bigr)=\zeta(2)-\sum_k\frac1{k(k+1)}=\zeta(2)-1$.

\tools{Volume~4, Corollary~11.6 (harmonic numbers), Theorem~5.8 (Stirling).}
\leads $\int_0^1\{\frac1x\}x^{s-1}\dd x=\frac1{s-1}-\frac{\zeta(s)}s$ for $\operatorname{Re}s>1$, and the left side is holomorphic for $\operatorname{Re}s>0$; this is how the
fractional part continues $\zeta$ to $\operatorname{Re}s>0$ (Volume~4, Chapter~10). Part (a) is the value at $s=1$, where the pole of $\zeta$ cancels.
\end{solution}

\begin{solution}{pb:primefrac}
\idea The values $\{\frac xp\}$ are $\varphi(\frac px)$ with $\varphi(t)=\{\frac1t\}$. By the prime number theorem the points $\frac px$ are spread over
$]0,1]$ with density $\frac x{\log x}\dd t$ on every fixed interval $[a,b]\subset]0,1]$, because $\log(xt)\sim\log x$ there. So the average of
$\varphi(\frac px)$ should be $\int_0^1\varphi=1-\gamma$. The work is in the approximation.

\writeup Let $\varepsilon>0$ and $K\ge2$. On $[\frac1K,1]$, $\varphi$ is bounded with finitely many discontinuities, hence Riemann integrable: there
are step functions $s\le\varphi\le S$ on $[\frac1K,1]$, constant on the intervals $]t_{j-1},t_j]$ of a subdivision $\frac1K=t_0<\cdots<t_m=1$, with
$\int(S-s)<\varepsilon$. By the prime number theorem (Volume~5, Theorem~5.5), for each fixed $t\in]0,1]$, $\pi(xt)=\frac{xt}{\log(xt)}(1+o(1))=\frac{xt}{\log x}(1+o(1))$.
Hence, for each $j$, the number of primes in $]xt_{j-1},xt_j]$ is $\frac x{\log x}(t_j-t_{j-1})+o(\frac x{\log x})$, and
\[
\sum_{x/K<p\le x}S\Bigl(\frac px\Bigr)=\frac x{\log x}\int_{1/K}^1S+o\Bigl(\frac x{\log x}\Bigr),
\]
and the same for $s$. The sum of $\varphi(\frac px)$ is squeezed between these two. The primes $p\le\frac xK$ contribute between $0$ and
$\pi(\frac xK)\le\frac{2x}{K\log x}$ for large $x$. Finally $\int_0^{1/K}\varphi\le\frac1K$. Dividing by $\pi(x)\sim\frac x{\log x}$,
\[
\limsup_{x\to\infty}\Bigl|\frac1{\pi(x)}\sum_{p\le x}\Bigl\{\frac xp\Bigr\}-\int_0^1\varphi\Bigr|\le\varepsilon+\frac3K,
\]
and $\varepsilon$, $K$ are arbitrary. By Problem~\ref{pb:fracint}(a), $\int_0^1\varphi=1-\gamma$.

\tools{Volume~5, Theorem~5.5 (prime number theorem); Volume~4, Chapter~4 (Riemann integrability, step functions); Problem~\ref{pb:fracint}.}
\leads De la Vallée Poussin proved this in 1898, two years after his proof of the prime number theorem, as an example of a ``mean
value'' over primes. The rate of convergence is tied to the error term in the prime number theorem, and therefore to the zeros of
$\zeta$: the same remark as for the divisor problem, with primes instead of integers.
\end{solution}

\begin{solution}{pb:leastpf}
\idea Weighting the primes by $p$ makes the largest primes count most, so the sum behaves like $x\cdot\pi(x)$ corrected by an integral. The least prime factor of a composite is tiny,
so the composites do not matter.

\writeup (a) By Abel summation (Volume~5, Lemma~1.1) with $a_p=p$, or directly, $\sum_{p\le x}p=x\pi(x)-\int_2^x\pi(t)\dd t$. By the prime number theorem $x\pi(x)\sim\frac{x^2}{\log x}$.
For the integral, fix $\varepsilon>0$; for $t\ge t_0$, $\pi(t)\le(1+\varepsilon)\frac t{\log t}$, and $\frac t{\log t}\le\frac t{\log\sqrt x}=\frac{2t}{\log x}$ for $t\ge\sqrt x$. So
\[
\int_2^x\pi(t)\dd t\le O(t_0^2)+\int_2^{\sqrt x}t\dd t+(1+\varepsilon)\int_{\sqrt x}^x\frac t{\log t}\dd t.
\]
On $[x^{1-\varepsilon},x]$, $\frac1{\log t}\le\frac1{(1-\varepsilon)\log x}$, and on $[\sqrt x,x^{1-\varepsilon}]$ the integrand is at most $\frac{2t}{\log x}$, with
$\int_{\sqrt x}^{x^{1-\varepsilon}}t\dd t\le\frac12x^{2-2\varepsilon}$. Hence $\int_2^x\pi(t)\dd t\le\frac{1+\varepsilon}{1-\varepsilon}\frac{x^2}{2\log x}+o(\frac{x^2}{\log x})$, and the same argument
with $(1-\varepsilon)$ and $\frac1{\log t}\ge\frac1{\log x}$ gives the matching lower bound. So $\int_2^x\pi(t)\dd t\sim\frac{x^2}{2\log x}$ and $\sum_{p\le x}p\sim\frac{x^2}{\log x}-\frac{x^2}{2\log x}=\frac{x^2}{2\log x}$.

(b) If $n$ is prime, $P^-(n)=n$; if $n$ is composite, $P^-(n)^2\le n$, so $P^-(n)\le\sqrt x$. Hence $0\le\sum_{2\le n\le x}P^-(n)-\sum_{p\le x}p\le x\cdot\sqrt x=o\bigl(\frac{x^2}{\log x}\bigr)$.

(c) The sum of the first $n$ primes is $\sum_{p\le p_n}p\sim\frac{p_n^2}{2\log p_n}$ by (a). Since $p_n\sim n\log n$ (Volume~5, Exercise~5.4), $\log p_n\sim\log n$ and the sum is
$\sim\frac{n^2\log^2n}{2\log n}=\frac12n^2\log n$.

(d) Heuristically $\mathrm d\pi(t)\approx\frac{\mathrm dt}{\log t}$, so $\sum_{p\le x}p\approx\int_2^x\frac t{\log t}\dd t$. With $s=t^2$, $\mathrm ds=2t\dd t$ and $\log t=\frac12\log s$, so
$\frac{t\dd t}{\log t}=\frac{\mathrm ds}{\log s}$ and the integral is $\operatorname{Li}(x^2)$ up to a constant. Precisely, write $\pi=\operatorname{Li}+E$ with $E(t)=O(te^{-c\sqrt{\log t}})$ (Volume~5,
Corollary~7.9). Then
\[
\sum_{p\le x}p=\int_{2^-}^xt\dd\pi(t)=\int_2^x\frac t{\log t}\dd t+\Bigl[tE(t)\Bigr]_2^x-\int_2^xE(t)\dd t.
\]
The first integral is $\operatorname{Li}(x^2)-\operatorname{Li}(4)$. The boundary term is $O(x^2e^{-c\sqrt{\log x}})$, and so is the last integral, since $te^{-c\sqrt{\log t}}$ is increasing for
large $t$ and $\int_2^x\le x\cdot\max$. So $\sum_{p\le x}p=\operatorname{Li}(x^2)+O(x^2e^{-c\sqrt{\log x}})$. The expansion $\operatorname{Li}(y)=\frac y{\log y}(1+\frac1{\log y}+\dots)$ at
$y=x^2$ gives $\frac{x^2}{2\log x}(1+\frac1{2\log x}+\dots)$; at $x=10^7$ the correction $\frac1{2\log x}$ is $3.1\%$, which explains the ratio $1.0326$.

\tools{Volume~5, Lemma~1.1 (Abel summation), Theorem~5.5 (prime number theorem), Corollary~7.9 (its error term), Exercise~5.4 ($p_n\sim n\log n$).}
\leads Weighting by $p$ moves the action to the top of the range, which is why the crude form $\frac{x^2}{2\log x}$ is so poor: it inherits the relative error $\frac1{\log x}$
of $\frac x{\log x}$. The same mechanism drives the next problem, where the weight is the largest prime factor.
\end{solution}

\begin{solution}{pb:largestpf}
\idea The largest prime factor is large exactly when $n$ is a prime times a small cofactor. Counting by the cofactor $k=\frac np$ turns the sum into $\sum_k\sum_{p\le x/k}p$,
and each inner sum is $\frac{x^2}{2k^2\log x}$ to first order: the $\frac{\pi^2}{12}$ is $\frac12\zeta(2)$.

\writeup (a) Every $n\ge2$ is uniquely $n=pm$ with $p=P^+(n)$ and $m\le\frac xp$ having no prime factor larger than $p$; conversely every such pair gives an $n\le x$ with $P^+(n)=p$.
So $\sum_{2\le n\le x}P^+(n)=\sum_{p\le x}p\,\Psi(\frac xp,p)$.

(b) For $p\le\sqrt x$, $p\,\Psi(\frac xp,p)\le p\cdot\frac xp=x$, and there are at most $\sqrt x$ such primes: total $O(x^{3/2})$. For $p>\sqrt x$, every $m\le\frac xp<\sqrt x<p$
has all its prime factors below $p$, so $\Psi(\frac xp,p)=\lfloor\frac xp\rfloor$.

(c) Let $T(y)=\sum_{p\le y}p$. Since $\lfloor\frac xp\rfloor=\#\{k\ge1:k\le\frac xp\}=\#\{k\ge1:p\le\frac xk\}$,
\[
\sum_{\sqrt x<p\le x}p\Bigl\lfloor\frac xp\Bigr\rfloor=\sum_{k\ge1}\ \sum_{\sqrt x<p\le x/k}p=\sum_{1\le k<\sqrt x}\bigl(T(\tfrac xk)-T(\sqrt x)\bigr),
\]
because the inner sum is empty when $\frac xk\le\sqrt x$. The terms $T(\sqrt x)\le x$ contribute at most $x^{3/2}$. For the rest, by Problem~\ref{pb:leastpf}(a), for each fixed $k$,
$\frac{\log x}{x^2}T(\frac xk)\to\frac1{2k^2}$. For a bound, $T(y)\le y\pi(y)\le C\frac{y^2}{\log y}$ (Chebyshev), and for $k<\sqrt x$, $\log\frac xk\ge\frac12\log x$, so
$\frac{\log x}{x^2}T(\frac xk)\le\frac{2C}{k^2}$. By dominated convergence for series,
\[
\frac{\log x}{x^2}\sum_{1\le k<\sqrt x}T\Bigl(\frac xk\Bigr)\to\sum_{k\ge1}\frac1{2k^2}=\frac{\pi^2}{12}.
\]
With (a) and (b), $\sum_{2\le n\le x}P^+(n)=\frac{\pi^2}{12}\frac{x^2}{\log x}(1+o(1))+O(x^{3/2})$, which is the claim.

(d) Write $n=P^+(n)\,m$. Then $A(n)-P^+(n)=A(m)$ is a sum of $\Omega(m)\le\log_2n$ primes, each at most $P^+(m)$. If $m>1$, $P^+(m)$ is a prime factor of $n$ other than the
largest one (counted with multiplicity), so $P^+(m)^2\le P^+(m)P^+(n)\le n$. Hence $0\le A(n)-P^+(n)\le\sqrt n\log_2n$, and summing,
$0\le\sum_{n\le x}A(n)-\sum_{n\le x}P^+(n)\le x^{3/2}\log_2x=o\bigl(\frac{x^2}{\log x}\bigr)$.

(e) In (c), $T(\frac xk)\approx\frac{x^2}{2k^2\log(x/k)}$, and $\log\frac xk<\log x$: each term is larger than its limit, by a factor $1+\frac{\log k}{\log x}+\dots$ The
corrections are of relative size $\frac1{\log x}$, about $6\%$ at $x=10^7$, so the convergence is logarithmically slow.

\tools{Volume~5, Theorem~5.5 and Proposition~1.2; Problem~\ref{pb:leastpf}; $\zeta(2)=\frac{\pi^2}6$ (Volume~4).}
\leads Alladi and Erd\H os (1977) found this asymptotic formula, in the equivalent form ``$A(n)$ has average order $\frac{\pi^2}6\frac n{\log n}$'', and extended it to the $k$-th largest prime
factor. The average of $\frac{\log P^+(n)}{\log n}$, a different question, leads to the Golomb--Dickman constant and to Dickman's function of Section~3.
\end{solution}

\begin{solution}{pb:factorialdiv}
\idea A prime $p>\sqrt n$ divides $n!$ exactly $\lfloor\frac np\rfloor$ times, and there are about $\frac n{k\log n}-\frac n{(k+1)\log n}$ primes with $\lfloor\frac np\rfloor=k$. Each contributes
$\log(1+k)$; summing by parts gives $\sum_k\frac1k\log(1+\frac1k)$.

\writeup (a) By Legendre's formula $n!=\prod_{p\le n}p^{v_p(n!)}$ with $v_p(n!)=\sum_j\lfloor\frac n{p^j}\rfloor$, and a number $\prod p^{a_p}$ has $\prod(1+a_p)$ divisors.

(b) For $p\le\sqrt n$, $\log(1+v_p(n!))\le\log(1+n)$, and there are at most $\sqrt n$ such primes: $O(\sqrt n\log n)$, which is $o(\frac n{\log n})$. (Even $O(\sqrt n)$ with
Chebyshev's bound, but we do not need it.) For $p>\sqrt n$, $p^2>n$, so $v_p(n!)=\lfloor\frac np\rfloor$.

(c) Telescoping, $\log(1+m)=\sum_{k=1}^m\log(1+\frac1k)$. So
\[
\sum_{\sqrt n<p\le n}\log\Bigl(1+\Bigl\lfloor\frac np\Bigr\rfloor\Bigr)=\sum_{k\ge1}\log\Bigl(1+\frac1k\Bigr)\#\Bigl\{\sqrt n<p\le\frac nk\Bigr\},
\]
which equals
\[
\sum_{1\le k<\sqrt n}\log\Bigl(1+\frac1k\Bigr)\Bigl(\pi\Bigl(\frac nk\Bigr)-\pi(\sqrt n)\Bigr).
\]
The terms with $\pi(\sqrt n)$ add up to at most $\sqrt n\sum_{k<\sqrt n}\frac1k=O(\sqrt n\log n)$. For each fixed $k$, $\frac{\log n}n\pi(\frac nk)\to\frac1k$, and for $k<\sqrt n$,
$\frac{\log n}n\pi(\frac nk)\le\frac{\log n}n\cdot\frac{Cn}{k\log(n/k)}\le\frac{2C}k$. Since $\log(1+\frac1k)\le\frac1k$, the terms are dominated by $\frac{2C}{k^2}$, and
\[
\frac{\log n}n\log\tau(n!)\to\sum_{k\ge1}\frac1k\log\Bigl(1+\frac1k\Bigr)=C.
\]
Finally $e^C=\prod_k(1+\frac1k)^{1/k}=2\cdot(\frac32)^{1/2}(\frac43)^{1/3}\cdots$

(d) The convergence is slow for the same reason as in Problem~\ref{pb:largestpf}: $\pi(\frac nk)\approx\frac n{k\log(n/k)}(1+\frac1{\log(n/k)})$, so the true sum is larger than its limit
by terms of relative order $\frac1{\log n}$; the excess over $C$ is $0.33$ at $n=10^4$ and $0.16$ at $n=10^7$, about $\frac3{\log n}$ in both cases.

\tools{Legendre's formula (Volume~1); Volume~5, Theorem~5.5 and Proposition~1.2; dominated convergence for series (Volume~0).}
\leads Erd\H os, Graham, Ivi\'c and Pomerance (1996) gave a full asymptotic expansion of $\log\tau(n!)$ in powers of $\frac1{\log n}$, and studied how fast $\tau((n+k)!)$ outgrows $\tau(n!)$.
\end{solution}

\begin{solution}{pb:centralbinom}
\idea For a large prime $p$, whether $p$ divides $\binom{2n}n$ depends only on the fractional part of $\frac np$: it does when there is a carry in $n+n$ in base $p$. The primes for
which this happens fill a union of intervals, whose total "prime mass" is $\sum_k\bigl(\frac1{k+1/2}-\frac1{k+1}\bigr)=2\log2$.

\writeup (a) By Legendre, $v_p(\binom{2n}n)=\sum_j\bigl(\lfloor\frac{2n}{p^j}\rfloor-2\lfloor\frac n{p^j}\rfloor\bigr)$, and for $p^2>2n$ only $j=1$ survives. With $t=\frac np=\lfloor t\rfloor+\{t\}$,
$\lfloor2t\rfloor-2\lfloor t\rfloor=\lfloor2\{t\}\rfloor$, which is $1$ if $\{t\}\ge\frac12$ and $0$ otherwise.

(b) $\{\frac np\}\ge\frac12$ means $k+\frac12\le\frac np<k+1$ with $k=\lfloor\frac np\rfloor\ge0$, that is $\frac n{k+1}<p\le\frac n{k+1/2}$.

(c) Let $I_k=\bigl]\frac n{k+1},\frac n{k+1/2}\bigr]$, of length $\ell_k=n\bigl(\frac1{k+1/2}-\frac1{k+1}\bigr)=\frac n{(2k+1)(k+1)}$. The primes $\le\sqrt{2n}$ contribute at most $\sqrt{2n}$.
For each fixed $k$, the prime number theorem gives $\#(I_k\cap\mathcal P)=\pi(\frac n{k+1/2})-\pi(\frac n{k+1})\sim\frac{\ell_k}{\log n}$. For $k\le n^{1/4}$, $\ell_k\ge\frac n{2(k+1)^2}\ge\frac{\sqrt n}3$
for large $n$, so $\log\ell_k\ge0.4\log n$ and Brun--Titchmarsh (Volume~5, Theorem~8.4) gives $\#(I_k\cap\mathcal P)\le\frac{3\ell_k}{\log\ell_k}\le\frac{8\ell_k}{\log n}$. As $\frac{\ell_k}n\le\frac1{(k+1)^2}$
is summable, dominated convergence gives
\[
\frac{\log n}n\sum_{k\le n^{1/4}}\#(I_k\cap\mathcal P)\to\sum_{k\ge0}\frac1{(2k+1)(k+1)}=2\sum_{k\ge0}\Bigl(\frac1{2k+1}-\frac1{2k+2}\Bigr)=2\log2.
\]
The intervals with $k>n^{1/4}$ lie below $n^{3/4}$ and contain at most $\pi(n^{3/4})=o(\frac n{\log n})$ primes. Hence $\omega\bigl(\binom{2n}n\bigr)\sim2\log2\cdot\frac n{\log n}=\log4\cdot\frac n{\log n}$.

(d) For $n=5\cdot10^6$, $\log4\cdot\frac n{\log n}=449\,370$, and $453\,587$ is $1.0094$ times this. (The ratio is above $1$ for the reason met in the previous problems: the
intervals $I_k$ sit at $\frac nk$, where $\log$ is smaller than $\log n$.)

\tools{Legendre's formula and Kummer's carries (Volume~1); Volume~5, Theorems~5.5 and~8.4.}
\leads The primes of $]n,2n]$ ($k=0$) give $\frac n{\log n}$, and those of $]\frac n2,\frac{2n}3]$ ($k=1$) give a sixth of that. Kummer's theorem says more: $v_p(\binom{2n}n)$ is the number of carries in $n+n$ in base $p$,
which for small $p$ can be large; but small primes are few.
\end{solution}

\begin{solution}{pb:superprimes}
\idea Applying the prime number theorem twice: $\pi(\pi(x))$ loses one factor $\log x$ at each step.

\writeup (a) With $m=p_n\sim n\log n$, $p_m\sim m\log m$ and $\log m\sim\log n$, so $p_{p_n}\sim n\log^2n$. Next, $\pi(x)\to\infty$ and $\pi(y)\sim\frac y{\log y}$, so
$\pi(\pi(x))\sim\frac{\pi(x)}{\log\pi(x)}$; since $\log\pi(x)=\log x-\log\log x+o(1)\sim\log x$, $\pi(\pi(x))\sim\frac x{\log^2x}$.

(b) By (a), $\frac1{p_{p_n}}\le\frac2{n\log^2n}$ for large $n$, and $\sum\frac1{n\log^2n}$ converges (compare with $\int\frac{\mathrm dt}{t\log^2t}=-\frac1{\log t}$). The partial sum up to $10^8$ stops at
$n=397\,557$, and the tail is about $\sum_{n>N}\frac1{n\log^2n}\approx\frac1{\log N}=0.078$. (This overestimates a little, since $p_{p_n}$ is larger than $n\log^2n$; the sum
of the series is a little above $1$.)

(c) $\pi(10^8)=5\,761\,455$, and $\frac y{\log y}(1+\frac1{\log y})$ at $y=5\,761\,455$ gives $393\,900$, within $1\%$ of $397\,557$, against $35\%$ for $\frac x{\log^2x}=294\,700$.

\tools{Volume~5, Theorem~5.5 and Exercise~5.4.}
\leads Dressler and Parker (1975) showed with a computer-aided argument that every integer larger than $96$ is a sum of distinct primes of the form $p_{p_n}$: a thin set of primes,
of density $\frac1{\log^2x}$, is still a basis in this weak sense.
\end{solution}

\begin{solution}{pb:ramprimes}
\idea The function $f(x)=\pi(x)-\pi(\frac x2)$ is a random-looking staircase that drifts upwards like $\frac x{2\log x}$. The last time it is below $n$ is followed by an upward
step, and upward steps happen only at primes.

\writeup (a) $f$ is a right-continuous step function, increasing by $1$ at each prime $x=p$ and decreasing by $1$ at each $x=2p$; an up-step and a down-step never occur at the same
point, since a prime is never twice a prime. By the prime number theorem $f(x)\sim\frac x{2\log x}\to\infty$, so $S=\{x>0:f(x)<n\}$ is bounded;
it contains $]0,2[$, where $f=0$. Let $R_n=\sup S$. For $x>R_n$, $f(x)\ge n$, and by right continuity $f(R_n)\ge n$. There are points of $S$ arbitrarily close to $R_n$ from the
left, so $f(R_n^-)<n\le f(R_n)$: $f$ jumps up at $R_n$, so $R_n$ is a prime, the jump is $1$, and $f(R_n)=n$, $f(R_n^-)=n-1$. Clearly $R_n$ is the least real with the required property.

(b) $R_n\to\infty$ since $f$ is bounded on bounded sets. From $f(R_n)=n$, $n\sim\frac{R_n}{2\log R_n}$. Also $2n=\pi(p_{2n})\sim\frac{p_{2n}}{\log p_{2n}}$. Hence
$\frac{R_n}{\log R_n}\sim\frac{p_{2n}}{\log p_{2n}}$. Taking logarithms, $\log R_n-\log\log R_n=\log p_{2n}-\log\log p_{2n}+o(1)$, so $\log R_n\sim\log p_{2n}$, and then
$R_n\sim p_{2n}$.

(c) With $L=\log x$, $\pi(x)\approx\frac xL(1+\frac1L)$ and $\pi(\frac x2)\approx\frac x{2(L-\log2)}(1+\frac1L)\approx\frac x{2L}(1+\frac{1+\log2}L)$, so
$f(x)\approx\frac x{2L}(1+\frac{1-\log2}L)$. Equating $f(R_n)=n$ with $\pi(p_{2n})=2n$, that is $\frac{p_{2n}}{2L}(1+\frac1L)=n$, and ignoring the difference between $\log R_n$ and $\log p_{2n}$
in second-order terms,
\[
\frac{R_n}{p_{2n}}\approx\frac{\log R_n}{\log p_{2n}}\Bigl(1+\frac{\log2}{\log p_{2n}}\Bigr)\approx1+\frac{\log2}{\log p_{2n}}+\frac{\log(R_n/p_{2n})}{\log p_{2n}}.
\]
At $n=10^4$, $\log p_{2n}=12.32$ and $1+\frac{\log2}{12.32}=1.056$; including the last term, $1.056+\frac{0.074}{12.32}=1.062$. The true ratio is $\frac{242\,057}{224\,737}=1.077$. The
approximation explains the bulk of the excess, and also why it decays only like $\frac1{\log n}$.

\tools{Volume~5, Theorem~5.5; Bertrand's postulate (Volume~1).}
\leads Sondow (2009) also proved $p_{2n}<R_n$ for $n>1$ and conjectured $R_n<p_{3n}$, proved by Laishram (2010). Ramanujan's original paper (1919) found $R_1,\dots,R_5$ as the thresholds $2,11,17,29,41$.
\end{solution}

\begin{solution}{pb:pminusone}
\idea For a fixed set of small primes $q$, the condition on $p-1$ is a congruence condition on $p$, and the prime number theorem in progressions gives its density exactly. The
large $q$ are a nuisance, and Brun--Titchmarsh shows that they affect only a small proportion of primes.

\writeup (a) For $q\le x^{1/4}$, $d=q^2\le\sqrt x$, and~\eqref{eq:bt13} gives $\pi(x;q^2,1)\le\frac{5x}{q(q-1)\log x}$. Since $\sum_{q>y}\frac1{q(q-1)}\le\sum_{m>y}\frac1{m(m-1)}\le\frac2y$, these $q$
account for at most $\frac{10x}{y\log x}$ primes. For $q>x^{1/4}$, the number of $p\le x$ with $q^2\mid p-1$ is at most $\frac x{q^2}$, and
$\sum_{q>x^{1/4}}\frac x{q^2}\le x\sum_{m>x^{1/4}}\frac1{m(m-1)}\le2x^{3/4}$.

(b) Fix $y$ and let $M=\prod_{q\le y}q^2$. For $p>y$, whether $q^2\mid p-1$ for some $q\le y$ depends only on $p\bmod M$. By the Chinese remainder theorem, the classes
$a\in(\Z/M\Z)^\times$ with $q^2\nmid a-1$ for all $q\le y$ number $\prod_{q\le y}(\varphi(q^2)-1)$. By the prime number theorem in progressions (Volume~5, Theorem~5.8), each class
receives $\sim\frac{\pi(x)}{\varphi(M)}$ primes, so the primes $p\le x$ with $q^2\nmid p-1$ for all $q\le y$ number $\sim A_y\pi(x)$, where
\[
A_y=\prod_{q\le y}\frac{\varphi(q^2)-1}{\varphi(q^2)}=\prod_{q\le y}\Bigl(1-\frac1{q(q-1)}\Bigr).
\]
The set of $p$ with $p-1$ squarefree is contained in this one, and differs from it by at most the primes counted in (a). As $\pi(x)\sim\frac x{\log x}$, the proportion of
primes $p\le x$ with $p-1$ squarefree has $\limsup\le A_y$ and $\liminf\ge A_y-\frac{10}y$. Let $y\to\infty$: $A_y\to A$, and the proportion tends to $A$.

(c) Using $\frac{\varphi(m)}m=\sum_{d\mid m}\frac{\mu(d)}d$, split according to $d\le y$, $y<d\le\sqrt x$, $d>\sqrt x$. First,
\[
\sum_{p\le x}\sum_{\substack{d\mid p-1\\d\le y}}\frac{\mu(d)}d=\sum_{d\le y}\frac{\mu(d)}d\pi(x;d,1)=\pi(x)\Bigl(\sum_{d\le y}\frac{\mu(d)}{d\varphi(d)}+o(1)\Bigr)
\]
by Volume~5, Theorem~5.8, $y$ being fixed. Second, by~\eqref{eq:bt13},
$\sum_{y<d\le\sqrt x}\frac1d\pi(x;d,1)\le\frac{5x}{\log x}\sum_{d>y}\frac1{d\varphi(d)}$, and the series converges since $\varphi(d)\ge\sqrt{d/2}$ (for odd $p$, $\varphi(p^k)=p^{k-1}(p-1)\ge p^{k/2}$, and
$\varphi(2^k)\ge2^{k/2}/\sqrt2$). Third, if $d\mid p-1$ with $d>\sqrt x$, then $e=\frac{p-1}d<\sqrt x$ and $\frac1d=\frac e{p-1}$; writing $p-1=em$,
\[
\sum_{p\le x}\sum_{\substack{d\mid p-1\\d>\sqrt x}}\frac1d\le\sum_{e<\sqrt x}e\sum_{m\le x/e}\frac1{em}\le\sqrt x(1+\log x)=o(\pi(x)).
\]
Hence $\frac1{\pi(x)}\sum_{p\le x}\frac{\varphi(p-1)}{p-1}$ lies within $6\sum_{d>y}\frac1{d\varphi(d)}+o(1)$ of $\sum_{d\le y}\frac{\mu(d)}{d\varphi(d)}$. Letting $y\to\infty$, the
limit is $\sum_d\frac{\mu(d)}{d\varphi(d)}=\prod_q\bigl(1-\frac1{q\varphi(q)}\bigr)=A$.

(d) $\mathbb F_p^\times$ is cyclic of order $p-1$, so it has $\varphi(p-1)$ generators: the proportion of primitive roots among the non-zero residues is $\frac{\varphi(p-1)}{p-1}$. Averaging over $p\le x$,
(c) says that a uniformly random prime $p\le x$ and a uniformly random non-zero residue modulo $p$ give a primitive root with probability tending to $A$. Artin's conjecture asks the same for a
\emph{fixed} residue such as $2$, and the data agree to four decimals.

\tools{Volume~5, Theorem~5.8 and Exercise~8.4 (Brun--Titchmarsh); Euler products (Volume~4); cyclicity of $\mathbb F_p^\times$ (Volume~1).}
\leads The coincidence of the two limits is no accident: $p-1$ squarefree and ``a random residue is a primitive root'' are both conditions prime by prime, $q^2\nmid p-1$ with probability $1-\frac1{q(q-1)}$
and ``not a $q$-th power when $q\mid p-1$'' with probability $1-\frac1{q(q-1)}$. Stephens (1969) proved (c) with an error term, and Hooley (1967) proved Artin's conjecture assuming the generalised
Riemann hypothesis.
\end{solution}

\begin{solution}{pb:titchmarsh}
\idea Pairing each divisor $d$ of $m$ with $\frac md$ shows that $\tau(m)$ is twice the number of small divisors; counting the primes $p\equiv1\pmod d$ for $d\le\sqrt x$ with Brun--Titchmarsh
then costs only a constant factor.

\writeup (a) $\frac d{\varphi(d)}=\prod_{p\mid d}\frac p{p-1}=\prod_{p\mid d}\bigl(1+\frac1{p-1}\bigr)=\sum_{e\mid d}\frac{\mu(e)^2}{\varphi(e)}$ (Volume~5, Exercise~2.2). Hence, writing $d=em$,
\[
\sum_{d\le y}\frac1{\varphi(d)}=\sum_{e\le y}\frac{\mu(e)^2}{e\varphi(e)}\sum_{m\le y/e}\frac1m\le(1+\log y)\sum_e\frac1{e\varphi(e)}=c(1+\log y),
\]
the last series converging as in Problem~\ref{pb:pminusone}(c).

(b) The divisors of $m$ come in pairs $(d,\frac md)$ with $d\le\sqrt m$ or $\frac md\le\sqrt m$, so $\tau(m)\le2\#\{d\mid m:d\le\sqrt m\}$. With $m=p-1<x$,
\[
\sum_{p\le x}\tau(p-1)\le2\sum_{d\le\sqrt x}\pi(x;d,1)\le\frac{10x}{\log x}\sum_{d\le\sqrt x}\frac1{\varphi(d)}\le\frac{10x}{\log x}c\Bigl(1+\frac12\log x\Bigr)=O(x)
\]
by~\eqref{eq:bt13} and (a).

(c) Now $\frac m{\varphi(m)}=\sum_{d\mid m}\frac{\mu(d)^2}{\varphi(d)}$. The argument of Problem~\ref{pb:pminusone}(c) applies. For $d\le y$, the main term is $\pi(x)\sum_{d\le y}\frac{\mu(d)^2}{\varphi(d)^2}$.
For $y<d\le\sqrt x$, the contribution is at most $\frac{5x}{\log x}\sum_{d>y}\frac1{\varphi(d)^2}$; the series converges because $\frac d{\varphi(d)}\le2^{\omega(d)}\le\tau(d)\le C d^{1/4}$
(Volume~5, Exercise~2.6), so $\varphi(d)\ge\frac1Cd^{3/4}$. For $d>\sqrt x$, put $e=\frac{p-1}d<\sqrt x$; since $\varphi(de)\le e\varphi(d)$, $\frac1{\varphi(d)}\le\frac e{\varphi(p-1)}\le\frac{Ce}{(p-1)^{3/4}}$, and
\[
\sum_{p\le x}\sum_{\substack{d\mid p-1\\d>\sqrt x}}\frac1{\varphi(d)}\le C\sum_{e<\sqrt x}e\sum_{m\le x/e}\frac1{(em)^{3/4}}\le4C\sum_{e<\sqrt x}e^{1/4}\Bigl(\frac xe\Bigr)^{1/4}\le4Cx^{3/4}.
\]
Letting $x\to\infty$ and then $y\to\infty$, the mean tends to $\sum_d\frac{\mu(d)^2}{\varphi(d)^2}=\prod_q\bigl(1+\frac1{(q-1)^2}\bigr)=B$.

(d) Pretend that $\tau(p-1)=2\#\{d\mid p-1:d\le\sqrt x\}$ and that $\pi(x;d,1)\approx\frac{\pi(x)}{\varphi(d)}$ for all $d\le\sqrt x$. Then
\[
\sum_{p\le x}\tau(p-1)\approx2\pi(x)\sum_{d\le\sqrt x}\frac1{\varphi(d)}\approx2\frac x{\log x}\cdot Z\log\sqrt x=Zx.
\]
The data, $1.916$, $1.927$, $1.932$, creep up towards $Z=1.9436$.

\tools{Volume~5, Exercises~2.2, 2.6 and~8.4, Theorem~5.8.}
\leads Two different constants appear: on average over primes, $\frac{p-1}{\varphi(p-1)}$ is $B=2.83$, larger than the mean $Z=1.94$ of $\frac n{\varphi(n)}$ over all integers, because $p-1$ is
always even and more often divisible by small primes. The asymptotic formula $\sum_{p\le x}\tau(p-1)\sim Zx$ is the \emph{Titchmarsh divisor problem}: Titchmarsh (1930) proved it under the
generalised Riemann hypothesis, Linnik (1961) unconditionally with his dispersion method, and Halberstam (1967) gave a short proof from the Bombieri--Vinogradov theorem, the result that
makes the heuristic of (d) rigorous on average over $d\le\sqrt x/\log^Bx$.
\end{solution}

\begin{solution}{pb:menon}
\idea The number of roots is a gcd, and averaging a gcd over a group of units is a job for Burnside's lemma. The answer $\tau(n)$ is the number of irreducible factors of $x^n-1$ over $\Q$,
and each factor contributes one root on average.

\writeup (a) $\mathbb F_p^\times$ is cyclic of order $p-1$, and in a cyclic group of order $N$ the equation $x^n=1$ has exactly $\gcd(n,N)$ solutions. ($0$ is never a root.)

(b) Two elements $x,x'$ of $\Z/n\Z$ are in the same orbit if and only if $\gcd(x,n)=\gcd(x',n)$. Indeed multiplication by a unit preserves $\gcd(x,n)$; conversely if
$\gcd(x,n)=\gcd(x',n)=d$, then $x=dv$, $x'=dv'$ with $v,v'$ units modulo $\frac nd$, and since $(\Z/n\Z)^\times\to(\Z/\frac nd\Z)^\times$ is onto, some unit $a$ modulo $n$
reduces to $v'v^{-1}$ modulo $\frac nd$, and then $ax\equiv x'\pmod n$. So the orbits correspond to the divisors $d$ of $n$, and there are $\tau(n)$ of them. The fixed points of $a$
are the $x$ with $(a-1)x\equiv0\pmod n$, and there are $\gcd(a-1,n)$ of them. Burnside's lemma (Volume~2) gives
$\tau(n)=\frac1{\varphi(n)}\sum_{a\in(\Z/n\Z)^\times}\gcd(a-1,n)$.

(c) For $p\nmid n$, $\gcd(n,p-1)=\gcd(n,a-1)$ where $a\equiv p\pmod n$. Ignoring the finitely many $p\mid n$,
\[
\frac1{\pi(x)}\sum_{p\le x}\gcd(n,p-1)=\sum_{a\in(\Z/n\Z)^\times}\gcd(a-1,n)\frac{\pi(x;n,a)}{\pi(x)}+o(1),
\]
which tends to
\[
\frac1{\varphi(n)}\sum_a\gcd(a-1,n)=\tau(n)
\]
by Volume~5, Theorem~5.8, and (b).

(d) For $p\nmid n$, $x^n-1$ is separable modulo $p$ (its derivative $nx^{n-1}$ has no common root with it), and its roots are the elements of $\mathbb F_p^\times$ of order dividing $n$.
The roots of $\Phi_d$ are those of order exactly $d$: there are $\varphi(d)$ of them if $d\mid p-1$, none otherwise. By Volume~5, Theorem~5.8, a proportion $\frac1{\varphi(d)}$ of primes satisfy
$d\mid p-1$, so $\Phi_d$ has on average $\varphi(d)\cdot\frac1{\varphi(d)}=1$ root. Summing over $d\mid n$ gives $\tau(n)$ again, and the identity
$\gcd(n,p-1)=\sum_{d\mid\gcd(n,p-1)}\varphi(d)$ is the bridge between (c) and (d).

\tools{Burnside's lemma (Volume~2); cyclotomic polynomials and cyclicity of $\mathbb F_p^\times$ (Volumes~1 and~2); Volume~5, Theorem~5.8.}
\leads This is a special case of a theorem that goes back to Kronecker and Frobenius: for any $f\in\Z[x]$, the average number of roots of $f$ modulo $p$ is the number of irreducible factors of
$f$ over $\Q$. The general proof needs the Chebotarev density theorem (Stevenhagen and Lenstra tell the story). Menon's identity has since been generalised in dozens of ways.
\end{solution}

\begin{solution}{pb:omegapoly}
\idea Exchange the order of summation: $\sum_n\omega(n^2+1)$ counts pairs $(n,p)$ with $p\mid n^2+1$, and for each $p$ there are $\rho(p)$ residue classes of $n$. The primes with
$\rho(p)=2$ and $\rho(p)=0$ balance out to an average of one class per prime, as for $\omega(n)$ itself. For the variance, the classes modulo $p$ and $q$ are independent by the
Chinese remainder theorem, exactly as in Tur\'an's proof for $\omega(n)$.

\writeup (a) $n^2\equiv-1\pmod2$ means $n$ odd: one class. For odd $p$, $-1$ is a square modulo $p$ if and only if $p\equiv1\pmod4$ (Euler's criterion, Volume~1), and then it has
two square roots. So $\rho(p)=2$ or $0$, that is $1+\chi(p)$.

(b) If $n\le x$ and $n^2+1$ had two prime factors larger than $x$, then $n^2+1\ge(\lfloor x\rfloor+1)^2>\lfloor x\rfloor^2+1\ge n^2+1$. So
\[
\sum_{n\le x}\omega(n^2+1)=\sum_{p\le x}\#\{n\le x:p\mid n^2+1\}+O(x).
\]
For $p\le x$ the $n\le x$ with $p\mid n^2+1$ fill $\rho(p)$ classes modulo $p$, so there are $\rho(p)(\frac xp+O(1))$ of them. The errors add up to $O(\pi(x))$, and the claim follows.

(c) $\sum_{p\le x}\frac{\rho(p)}p=\sum_{p\le x}\frac1p+\sum_{p\le x}\frac{\chi(p)}p$ (the term $p=2$ fits since $\chi(2)=0$). The first sum is $\log\log x+O(1)$ (Volume~5, Theorem~1.5),
and the second converges, by Volume~5, Corollary~4.5 applied to both classes modulo~$4$. So $\sum_{n\le x}\omega(n^2+1)=x\log\log x+O(x)$.

(d) Let $y=x^{1/4}$ and $\omega_y(m)$ the number of prime factors $p\le y$ of $m$. For $x\ge3$, $n^2+1\le x^2+1<y^9$, so $0\le\omega(n^2+1)-\omega_y(n^2+1)\le8$. Let
$S=\sum_{p\le y}\frac{\rho(p)}p=\log\log x+O(1)$. Then
\[
\sum_{n\le x}\omega_y(n^2+1)^2=\sum_{p,q\le y}\#\{n\le x:p\mid n^2+1,\ q\mid n^2+1\}.
\]
For $p\ne q$, the condition is $pq\mid n^2+1$, which by the Chinese remainder theorem defines $\rho(p)\rho(q)$ classes modulo $pq$, so the count is $\rho(p)\rho(q)(\frac x{pq}+O(1))$;
for $p=q$ it is $\rho(p)(\frac xp+O(1))$. Hence $\sum_n\omega_y^2\le xS^2+xS+O(\pi(y)^2)$, and similarly $\sum_n\omega_y=xS+O(\pi(y))$. Expanding the square,
\[
\sum_{n\le x}\bigl(\omega_y(n^2+1)-S\bigr)^2\le xS^2+xS-2S(xS-O(\pi(y)))+xS^2+O(\pi(y)^2)=xS+O(\sqrt x).
\]
Finally $\omega(n^2+1)-\log\log x=(\omega_y-S)+O(1)$, and $(a+b)^2\le2a^2+2b^2$ gives $\sum_{n\le x}(\omega(n^2+1)-\log\log x)^2=O(x\log\log x)$. By Chebyshev's inequality, for every
$\varepsilon>0$ the $n\le x$ with $|\omega(n^2+1)-\log\log x|>\varepsilon\log\log x$ number $O\bigl(\frac x{\varepsilon^2\log\log x}\bigr)=o(x)$; and for $\sqrt x<n\le x$,
$|\log\log n-\log\log x|\le\log2$. So $\omega(n^2+1)$ has normal order $\log\log n$.

(e) Yes: (c) says that the difference is bounded, and it is, at about $0.6$. Its limit is a constant built from the Meissel--Mertens constant, $\sum_p\frac{\chi(p)}p$, and the
density of the $n$ for which $n^2+1$ has a prime factor larger than $n$.

\tools{Euler's criterion (Volume~1); Volume~5, Theorem~1.5, Corollary~4.5, and the method of Tur\'an's variance (Chapter~9); Chebyshev's inequality (Volume~3).}
\leads The same argument works for any irreducible polynomial $f$, with $\rho_f(p)$ averaging to $1$ by the theorem of Kronecker and Frobenius met in Problem~\ref{pb:menon}; and the
Erd\H os--Kac theorem holds for $\omega(f(n))$ too. The variance here is only bounded above; one finds numerically it smaller than $\log\log x$, partly because the large prime factors
have been thrown away.
\end{solution}

\begin{solution}{pb:rankin}
\idea Counting the smooth numbers up to $x$ is hard; summing $(\frac xn)^\sigma$ over \emph{all} smooth numbers is easy, because it factors as an Euler product. Choosing $\sigma$ slightly less
than $1$ balances the gain $x^{\sigma-1}$ against the growth of the product.

\writeup (a) For a $y$-smooth $n\le x$, $(\frac xn)^\sigma\ge1$. So $\Psi(x,y)\le\sum_{n\ y\text{-smooth}}(\frac xn)^\sigma=x^\sigma\prod_{p\le y}(1-p^{-\sigma})^{-1}$, the sum over all
$y$-smooth $n$ being an Euler product with finitely many factors.

(b) Since $\sigma\ge\frac23$, $\sum_p\sum_{k\ge2}\frac{p^{-k\sigma}}k\le\sum_p\frac{p^{-4/3}}{1-2^{-2/3}}$ is bounded by a constant. Next, with $t=\frac{\log p}{\log y}\in[0,1]$,
$p^{-\sigma}=\frac1pe^{\eta t}\le\frac1p\bigl(1+(e^\eta-1)\frac{\log p}{\log y}\bigr)$, so by Mertens's first and second theorems (Volume~5, Theorems~1.3 and~1.5)
\[
\sum_{p\le y}p^{-\sigma}\le\log\log y+c_1+(e^\eta-1)\frac{\log y+c_2}{\log y}\le\log\log y+c_1+\Bigl(1+\frac{c_2}{\log2}\Bigr)e^\eta.
\]
Exponentiating, $\prod_{p\le y}(1-p^{-\sigma})^{-1}\le C\log y\,e^{Ce^\eta}$, and $x^\sigma=xe^{-\eta\log x/\log y}=xe^{-\eta u}$.

(c) With $\eta=\log u$ (allowed when $u>1$ and $\log u\le\frac13\log y$), $e^{-\eta u+Ce^\eta}=e^{-u\log u+Cu}$. For $y=e^{\sqrt{\log x}}$, $u=\log y=\sqrt{\log x}$, and the condition
$\frac12\log\log x\le\frac13\sqrt{\log x}$ holds for large $x$. Then $\frac{\Psi(x,y)}x\le C\sqrt{\log x}\,e^{-\frac12\sqrt{\log x}\log\log x+C\sqrt{\log x}}=e^{-(\frac12+o(1))\sqrt{\log x}\log\log x}$.

(d) The best bounds miss by factors $\frac{9.95\cdot10^7}{8.80\cdot10^6}=11.3$ and $\frac{1.26\cdot10^{11}}{6.47\cdot10^9}=19.4$. The loss is the price of (a): the smooth $n>x$ are counted with weights
$(\frac xn)^\sigma$ that decay slowly, and a whole range of $n<x$ is counted with weights larger than $1$. The product $\prod_{p\le y}(1-\frac1p)^{-1}\sim e^\gamma\log y$ already shows the
factor $\log y$ that the method cannot remove.

\tools{Volume~5, Theorems~1.3 and~1.5 (Mertens); Euler products (Volume~4).}
\leads Rankin (1938) used this trick to find large gaps between primes. Refined with a saddle point, the same idea gives $\Psi(x,y)$ asymptotically in a huge range (Hildebrand and Tenenbaum):
the loss in (d) is a factor of order $\sqrt u\,\log y$, which the saddle-point method recovers.
\end{solution}

\begin{solution}{pb:renyi}
\idea $z^{\Omega(n)-\omega(n)}$ is a multiplicative function that is $1$ on primes and differs from $1$ only at prime powers, so it is a convolution $\mathbf1*h$ with $h$ supported on
squarefull numbers, and Wintner's theorem gives its mean as an Euler product. The pole of the factor at $p=2$ decides the tail of the law.

\writeup (a) $\Omega-\omega$ is additive, $\Omega(p^k)-\omega(p^k)=k-1$, so $f_z$ is multiplicative with $f_z(p^k)=z^{k-1}$. Then $h_z=f_z*\mu$ is multiplicative with
$h_z(p)=f_z(p)-1=0$ and $h_z(p^k)=z^{k-1}-z^{k-2}=z^{k-2}(z-1)$ for $k\ge2$. For $|z|<2$,
\[
\sum_d\frac{|h_z(d)|}d=\prod_p\Bigl(1+|z-1|\sum_{k\ge2}\frac{|z|^{k-2}}{p^k}\Bigr)=\prod_p\Bigl(1+\frac{|z-1|}{p(p-|z|)}\Bigr)<\infty,
\]
since $p-|z|>0$ for all $p\ge2$ and $\sum_p\frac1{p^2}$ converges.

(b) By Wintner's theorem the mean value is $\sum_d\frac{h_z(d)}d=\prod_p\bigl(1+\frac{z-1}{p(p-z)}\bigr)$. And
$(1-\frac1p)(1+\frac1{p-z})=\frac{(p-1)(p-z+1)}{p(p-z)}=\frac{p^2-pz+z-1}{p(p-z)}=1+\frac{z-1}{p(p-z)}$.

(c) Let $d_k(x)=\frac1x\#\{n\le x:\Omega(n)-\omega(n)=k\}$ and $P_x(z)=\sum_kd_k(x)z^k=\frac1x\sum_{n\le x}z^{\Omega(n)-\omega(n)}$, a polynomial. By Cauchy's formula on the unit circle,
$d_k(x)=\frac1{2\pi}\int_0^{2\pi}P_x(e^{i\theta})e^{-ik\theta}\dd\theta$. By (b) $P_x(e^{i\theta})\to R(e^{i\theta})$ for every $\theta$, and $|P_x|\le\frac{\lfloor x\rfloor}x\le1$ there, so by dominated
convergence $d_k(x)\to\frac1{2\pi}\int R(e^{i\theta})e^{-ik\theta}\dd\theta=[z^k]R$, $R$ being holomorphic on $|z|<2$ (Volume~4, Theorem~9.10). So $d_k$ exists and $\sum d_kz^k=R(z)$.
Then $d_0=R(0)=\prod_p(1-\frac1{p^2})=\frac1{\zeta(2)}=\frac6{\pi^2}$: the squarefree numbers, as it should be. The mean is $R'(1)=\sum_p\frac{\mathrm d}{\mathrm dz}\log\bigl(1+\frac{z-1}{p(p-z)}\bigr)\big|_{z=1}=\sum_p\frac1{p(p-1)}$,
since $R(1)=1$ and the derivative of $\frac{z-1}{p(p-z)}$ at $1$ is $\frac1{p(p-1)}$.

(d) Write $R(z)=\frac12\bigl(1+\frac1{2-z}\bigr)R_3(z)$ with $R_3(z)=\prod_{p\ge3}(1-\frac1p)(1+\frac1{p-z})$, holomorphic on $|z|<3$. Then
$R(z)-\frac{R_3(2)}{2(2-z)}$ is holomorphic on $|z|<3$, so by Cauchy's estimates (Volume~4, Corollary~9.12) its coefficients are $O(2.5^{-k})$, while $\frac{R_3(2)}{2(2-z)}=\frac{R_3(2)}4\sum_k\frac{z^k}{2^k}$. Hence
$2^kd_k\to\frac{R_3(2)}4$, and
\[
R_3(2)=\prod_{p\ge3}\frac{p-1}p\cdot\frac{p-1}{p-2}=\prod_{p\ge3}\frac{(p-1)^2}{p(p-2)}=\prod_{p\ge3}\Bigl(1-\frac1{(p-1)^2}\Bigr)^{-1}=\frac1{C_2}.
\]

(e) The agreement is to five or six decimals: the convergence of $\frac1x\sum_{n\le x}f(n)$ to the mean is fast here, with error $O(x^{-1/2+\varepsilon})$ since $h$ lives on squarefull numbers.

\tools{Volume~5, Theorem~2.2 (Wintner); Volume~4, Theorem~9.10 and Corollary~9.12 (Cauchy's formula and estimates).}
\leads R\'enyi's law is a law of repeated prime factors, governed by the small primes, unlike the Erd\H os--Kac law of $\omega(n)$, governed by all of them. The twin prime constant enters
through the factor $\frac{p-1}{p-2}$, the same factor that makes pairs $p,p+2$ more common than independent coins would (Problem~\ref{pb:digitpairs}).
\end{solution}

\begin{solution}{pb:twosquares}
\idea Squaring $F$ produces $\zeta(s)L(s,\chi_4)$, so $F$ behaves like $(s-1)^{-1/2}$: a ``half pole''. Rankin's trick turns this into the size $\sqrt{\log y}$ of the logarithmic
sum $S(y)$, and two elementary arguments, one through prime powers and one through primes $p\equiv1\pmod4$, transfer it to $B(x)$.

\writeup (a) By Volume~1, Theorem~7.3, $b$ is multiplicative, with $b(2^k)=1$, $b(p^k)=1$ for $p\equiv1\pmod4$, and $b(p^k)=[k\text{ even}]$ for $p\equiv3\pmod4$. So $F(s)$ has Euler
factors $(1-2^{-s})^{-1}$, $(1-p^{-s})^{-1}$ and $(1-p^{-2s})^{-1}$. The factors of $\zeta(s)L(s,\chi_4)$ are $(1-2^{-s})^{-1}$, $(1-p^{-s})^{-2}$ and $(1-p^{-s})^{-1}(1+p^{-s})^{-1}=(1-p^{-2s})^{-1}$.
Comparing squares gives the identity. As $\sigma\to1^+$, $\zeta(\sigma)\sim\frac1{\sigma-1}$, $L(1,\chi_4)=\frac\pi4$, $\frac1{1-2^{-1}}=2$, and the last product converges at $s=1$, so
$F(\sigma)^2\sim\frac{C^2}{\sigma-1}$.

(b) For $n\le y$, $n^{-\delta}\ge e^{-1}$, so $\frac1n\le e\,n^{-1-\delta}$ and $S(y)\le eF(1+\delta)$. For $n>y$, $n^{-1-\delta}\le y^{-\delta/2}n^{-1-\delta/2}=e^{-1/2}n^{-1-\delta/2}$, so
$S(y)\ge\sum_{n\le y}b(n)n^{-1-\delta}\ge F(1+\delta)-e^{-1/2}F(1+\frac\delta2)$. By (a), $F(1+\delta)\sim C\sqrt{\log y}$ and $F(1+\frac\delta2)\sim C\sqrt{2\log y}$, so
$C\bigl(1-\sqrt2e^{-1/2}+o(1)\bigr)\sqrt{\log y}\le S(y)\le C(e+o(1))\sqrt{\log y}$, and $1-\sqrt2e^{-1/2}=0.142\ldots>0$.

(c) For $p\le\sqrt y$ there are at most $\frac{\log y}{\log p}$ values of $k$, so $\sum_{p^k\le y}k\log p\le\sum_{p\le\sqrt y}\frac{\log^2y}{\log p}+\sum_{p\le y}\log p=O(\sqrt y\log y)+O(y)$
by Chebyshev (Volume~5, Proposition~1.2). If $b(n)=1$ and $p^k\parallel n$, then $m=\frac n{p^k}$ has $b(m)=1$. Writing $\log n=\sum_{p^k\parallel n}k\log p$ and counting $n$ by $m$,
\[
\sum_{n\le x}b(n)\log n\le\sum_{m\le x}b(m)\sum_{p^k\le x/m}k\log p\le c\sum_{m\le x}b(m)\frac xm=c\,x\,S(x).
\]
With $L=\log x$, $\log\frac xL\,\bigl(B(x)-B(\frac xL)\bigr)\le\sum_{n\le x}b(n)\log n=O(x\sqrt{\log x})$ by (b), so $B(x)\le\frac x{\log x}+O\bigl(\frac{x\sqrt{\log x}}{\log x}\bigr)=O\bigl(\frac x{\sqrt{\log x}}\bigr)$.

(d) If $b(m)=1$ and $p\equiv1\pmod4$, then $b(mp)=1$, and conversely $b(n)=1$, $p\mid n$ give $b(\frac np)=1$. So pairs $(m,p)$ with $mp\le x$ correspond to pairs $(n,p)$ with $p\mid n$, and
\[
\sum_{n\le x}b(n)\log n\ge\sum_{n\le x}b(n)\sum_{p\mid n,\,p\equiv1\,(4)}\log p=\sum_{m\le x}b(m)\,\theta\bigl(\tfrac xm;4,1\bigr)\ge\sum_{m\le x/\log^3x}b(m)\,\theta\bigl(\tfrac xm;4,1\bigr).
\]
For $m\le\frac x{\log^3x}$, $\frac xm\ge\log^3x\to\infty$, and $\theta(y;4,1)\sim\frac y2$ (Volume~5, Theorem~5.8), so the right side is at least $(\frac12-o(1))\,x\,S(\frac x{\log^3x})$, which
is $\ge c\,x\sqrt{\log x}$ by (b). Since $\sum_{n\le x}b(n)\log n\le B(x)\log x$, $B(x)\ge\frac{c\,x}{\sqrt{\log x}}$.

(e) By partial summation, for $\sigma>1$, $F(\sigma)=\sigma\int_1^\infty B(x)x^{-\sigma-1}\dd x$. If $B(x)=(K+o(1))\frac x{\sqrt{\log x}}$, the part $x\le X$ is bounded for each fixed $X$, and
with $x=e^u$,
\[
\int_X^\infty\frac{x^{-\sigma}}{\sqrt{\log x}}\dd x=\int_{\log X}^\infty\frac{e^{-(\sigma-1)u}}{\sqrt u}\dd u=\frac{\Gamma(\frac12)+o(1)}{\sqrt{\sigma-1}}\qquad(\sigma\to1^+).
\]
So $F(\sigma)\sim K\sqrt\pi(\sigma-1)^{-1/2}$, and (a) forces $K=\frac C{\sqrt\pi}$. Numerically $\prod_{p\equiv3\,(4)}(1-p^{-2})^{-1}=1.168076\ldots$, $C=1.35455\ldots$ and $K=0.764224\ldots$
the numerical value $0.7922$ is $1.037K$.

\tools{Volume~1, Theorem~7.3; Volume~5, Proposition~1.2 and Theorem~5.8; $L(1,\chi_4)=\frac\pi4$ (Leibniz); Rankin's trick (Problem~\ref{pb:rankin}); $\Gamma(\frac12)=\sqrt\pi$.}

\leads Landau proved in 1908 that $B(x)\sim\frac{Kx}{\sqrt{\log x}}$ indeed, with complex analysis: $F(s)$ is $\sqrt{\zeta(s)}$ times a function holomorphic near $s=1$, and he integrated
around the branch cut. Volume~7 proves the logarithmic version $\sum_{n\le x}\frac{b(n)}n\sim2K\sqrt{\log x}$ with Karamata's Tauberian theorem. The slow convergence seen in the
data comes from a secondary term: $B(x)=\frac{Kx}{\sqrt{\log x}}\bigl(1+\frac{c}{\log x}+\cdots\bigr)$, and at $x=10^8$ the correction $\frac c{\log x}$ is still $3.7\%$.
\end{solution}

\begin{solution}{pb:ramanujanE2}
\idea $E_2$ is almost a modular form: its transformation law has an extra term $\frac{6\tau}{\pi i}$, inherited from the $\sqrt{-i\tau}$ in Dedekind's law. Differentiating a modular
form also spoils modularity by a term of the same kind, and the two defects cancel in $DE_2-\frac{E_2^2}{12}$. Since the spaces $M_4$, $M_6$, $M_8$ are one-dimensional, each corrected
derivative is a multiple of a known form.

\writeup (a) $\log\eta=\frac{2\pi i\tau}{24}+\sum\log(1-q^n)$, so $D\log\eta=\frac1{24}-\sum_n\frac{nq^n}{1-q^n}=\frac1{24}-\sum_N\sigma(N)q^N=\frac{E_2}{24}$.

(b) Apply $D$ to $\log\eta(-\frac1\tau)=\log\eta(\tau)+\frac12\log(-i\tau)$. By the chain rule the left side gives $\frac1{\tau^2}\frac{E_2}{24}(-\frac1\tau)$, the right side
$\frac{E_2(\tau)}{24}+\frac1{2\pi i}\frac1{2\tau}$. Multiplying by $24\tau^2$: $E_2(-\frac1\tau)=\tau^2E_2(\tau)+\frac{12\tau}{2\pi i}=\tau^2E_2(\tau)+\frac{6\tau}{\pi i}$.

(c) Write $c=\frac{12}{2\pi i}$, so that $E_2(-\frac1\tau)=\tau^2E_2+c\tau$. Applying $D$ and the chain rule, $(DE_2)(-\frac1\tau)=\tau^4DE_2+\frac{2\tau^3E_2+c\tau^2}{2\pi i}$, and
$E_2(-\frac1\tau)^2=\tau^4E_2^2+2c\tau^3E_2+c^2\tau^2$. In $f(-\frac1\tau)$ the $\tau^3E_2$ terms are $\frac2{2\pi i}-\frac{2c}{12}=0$ and the $\tau^2$ terms are $c\bigl(\frac1{2\pi i}-\frac c{12}\bigr)=0$,
so $f(-\frac1\tau)=\tau^4f(\tau)$. Also $f(\tau+1)=f(\tau)$, and $f=-\frac1{12}+O(q)$ is holomorphic at $i\infty$. So $f\in M_4=\C E_4$ (Volume~5, Theorem~12.9), and $f=-\frac{E_4}{12}$.
In the same way, for a form $g$ of weight $k$, $(Dg)(-\frac1\tau)=\tau^{k+2}Dg+\frac{k\tau^{k+1}g}{2\pi i}$, so $Dg-\frac k{12}E_2g$ has weight $k+2$. For $g=E_4$ it lies in $M_6=\C E_6$
with constant term $-\frac13$; for $g=E_6$ it lies in $M_8=\C E_4^2$ with constant term $-\frac12$. These are Ramanujan's equations.

(d) The coefficient of $q^n$ in $12DE_2$ is $-288n\sigma(n)$; in $E_2^2$ it is
\[
-48\sigma(n)+576\sum_{m=1}^{n-1}\sigma(m)\sigma(n-m);
\]
in $E_4$ it is $240\sigma_3(n)$. Solving,
$576\sum\sigma(m)\sigma(n-m)=240\sigma_3(n)+(48-288n)\sigma(n)$, which is the formula after division by $48$.

(e) In $3DE_4=E_2E_4-E_6$, the coefficients of $q^n$ give $720n\sigma_3(n)=240\sigma_3(n)-24\sigma(n)-5760\sum\sigma(m)\sigma_3(n-m)+504\sigma_5(n)$; divide by $24$.

(f) With $\sigma(p)=p+1$ and $\sigma_3(p)=p^3+1$, (d) gives $\frac{5p^3+5+(1-6p)(p+1)}{12}=\frac{5p^3-6p^2-5p+6}{12}=\frac{(p-1)(5p^2-p-6)}{12}$. For $p=7$: $\sigma(1),\dots,\sigma(6)=1,3,4,7,6,12$ and
$1\cdot12+3\cdot6+4\cdot7+7\cdot4+6\cdot3+12\cdot1=116=\frac{6\cdot232}{12}$.

\tools{Volume~5, Theorem~12.7 (Dedekind's $\eta(-\frac1\tau)=\sqrt{-i\tau}\,\eta(\tau)$) and Theorem~12.9 (dimensions of $M_k$, so $E_8=E_4^2$).}

\leads Ramanujan published these equations in 1916, in the paper where he introduced $\tau(n)$. They say that the ring $\C[E_2,E_4,E_6]$ is closed under $D$; it is the ring of
\emph{quasimodular forms}. Every convolution sum $\sum\sigma_a(m)\sigma_b(n-m)$ with $a$ and $b$ odd can be evaluated this way, except that from weight $12$ on the cusp form $\Delta$
enters and its coefficients $\tau(n)$ appear in the answer.
\end{solution}

\begin{solution}{pb:landauzeros}
\idea Integrate $x^s\frac{\zeta'}\zeta(s)$ around a tall rectangle. The residues are the $x^\rho$. On the right edge the Dirichlet series of $\frac{\zeta'}\zeta$ is integrated term by term,
and only the term $n=x$ does not oscillate: it grows like $T$. Everything else is smaller, so the zeros, summed with the weights $x^\rho$, ``detect'' whether $x$ is a prime power.

\writeup Fix $x>1$ and $T\ge3$, and choose $T'\in[T,T+1]$ as in Volume~5, Lemma~7.3, so that $|\frac{\zeta'}\zeta(\sigma+iT')|\le C\log^2T$ for $-1\le\sigma\le3$. No zero has $0<\gamma\le1$.

(a) On $\operatorname{Re}s=2$ the series $\frac{\zeta'}\zeta(s)=-\sum\Lambda(n)n^{-s}$ converges absolutely, so
\[
\frac1{2\pi}\int_1^Tx^{2+it}\frac{\zeta'}\zeta(2+it)\dd t=-\frac1{2\pi}\sum_n\Lambda(n)\Bigl(\frac xn\Bigr)^2\int_1^T\Bigl(\frac xn\Bigr)^{it}\dd t.
\]
The term $n=x$ (if $x$ is an integer) gives $-\frac{T-1}{2\pi}\Lambda(x)$. For $n\ne x$ the integral is at most $\frac2{|\log(x/n)|}$, and $\sum_{n\ne x}\Lambda(n)\frac{x^2}{n^2|\log(x/n)|}$
converges, since $|\log\frac xn|$ is bounded below for $n\ne x$ and $\sum\frac{\Lambda(n)}{n^2}<\infty$. The total is $-\frac T{2\pi}\Lambda(x)+O_x(1)$.

(b) Take the logarithmic derivative of $\xi(s)=\xi(1-s)$ with $\xi(s)=\frac12s(s-1)\pi^{-s/2}\Gamma(\frac s2)\zeta(s)$:
\[
\frac1s+\frac1{s-1}-\frac{\log\pi}2+\frac12\frac{\Gamma'}\Gamma\Bigl(\frac s2\Bigr)+\frac{\zeta'}\zeta(s)=\frac1s+\frac1{s-1}+\frac{\log\pi}2-\frac12\frac{\Gamma'}\Gamma\Bigl(\frac{1-s}2\Bigr)-\frac{\zeta'}\zeta(1-s),
\]
which is the stated formula. On $s=-1+it$, $\frac{\zeta'}\zeta(1-s)=\frac{\zeta'}\zeta(2-it)$ is again an absolutely convergent series, and integrating $x^{-1+it}\Lambda(n)n^{-2+it}$ gives
$O(\frac1{\log(xn)})$, summable in $n$. The constant $\log\pi$ gives $O(\frac1{\log x})$. The function $\varphi(t)=-\frac12\frac{\Gamma'}\Gamma(\frac{-1+it}2)-\frac12\frac{\Gamma'}\Gamma(\frac{2-it}2)$ satisfies
$\varphi(t)=O(\log t)$ and $\varphi'(t)=O(\frac1t)$ for $t\ge1$, by Stirling's formula for $\frac{\Gamma'}\Gamma$ (Volume~5, Theorem~6.1) and the reflection formula to move
$\frac{-1+it}2$ to the right half-plane. Integrating by parts,
$\int_1^{T}x^{it}\varphi(t)\dd t=\bigl[\frac{x^{it}}{i\log x}\varphi\bigr]_1^{T}-\int_1^{T}\frac{x^{it}}{i\log x}\varphi'(t)\dd t=O_x(\log T)$.

(c) The rectangle with vertices $2+i$, $2+iT'$, $-1+iT'$, $-1+i$ contains the zeros with $1<\gamma<T'$ and no other pole of $x^s\frac{\zeta'}\zeta$; the residue at $\rho$ is
$x^\rho$ times its multiplicity. So $\sum_{0<\gamma<T'}x^\rho$ is $\frac1{2\pi i}$ times the integral around the boundary. The right edge gives $-\frac{T'}{2\pi}\Lambda(x)+O_x(1)$ by (a), the
left edge $O_x(\log T)$ by (b), the bottom edge $O_x(1)$, and the top edge at most $\frac1{2\pi}\cdot3x^2\cdot C\log^2T$. Finally $T'-T\le1$, and there are $O(\log T)$ zeros with
$T<\gamma\le T'$ (Volume~5, Corollary~6.10), each with $|x^\rho|\le x$. This gives Landau's formula.

(d) Under the Riemann hypothesis $x^\rho=\sqrt x\,e^{i\gamma\log x}$. Take real parts in (c) and divide by $\sqrt x$.

(e) By Weyl's criterion (Volume~4, Theorem~6.13) it is enough to show that $\frac1N\sum_{n\le N}e^{2\pi ik\alpha\gamma_n}\to0$ for every integer $k\ne0$; replacing $k$ by $-k$ conjugates the sum,
so let $k\alpha>0$, and $y=e^{2\pi k\alpha}>1$. Then $e^{2\pi ik\alpha\gamma}=y^{i\gamma}=y^{-1/2}y^\rho$, so by (c)
$\sum_{0<\gamma\le T}e^{2\pi ik\alpha\gamma}=-\frac T{2\pi}\frac{\Lambda(y)}{\sqrt y}+O_y(\log^2T)=O(T)$, while the number of terms is $N(T)\sim\frac T{2\pi}\log T$ (Volume~5, Theorem~6.9).
The ratio tends to $0$, also along $T=\gamma_N$.

\tools{Volume~5, Proposition~6.2, Theorem~6.1, Lemma~7.3, Theorem~6.9 and Corollary~6.10; Weyl's criterion (Volume~4, Theorem~6.13).}

\leads Landau proved his formula in 1911; Gonek (1993) made it uniform in $x$, which is what one needs to use it as a tool. Part (e) is Rademacher's theorem (1956); Hlawka (1975)
removed the Riemann hypothesis. The spikes of Figure~\ref{fig:landau} are the formula with $2000$ zeros: at a prime power $p^k$ the value is within $3\%$ of $\log p$, and between the
spikes it is a small ripple.
\end{solution}

\begin{solution}{pb:hearzeros}
\idea The explicit formula is the residue theorem applied to $H\frac{\zeta'}\zeta$, where $H$ is a smooth bump on the critical line that is still entire. The zeros come out as
residues, the primes from the Dirichlet series on the right edge, and the functional equation folds the left edge back onto the right. Choosing $H$ concentrated near height
$T$ turns the formula into a zero detector built from primes.

\writeup (a) $h(w)=e^{-(w-T)^2/2s^2}+e^{-(w+T)^2/2s^2}$ is entire and even in $w$, so $H$ is entire and $H(1-z)=h(i(z-\frac12))=h(-i(z-\frac12))=H(z)$; and $H(\frac12+it)=h(t)$.
For $z=\sigma+it$, $w=t-i(\sigma-\frac12)$, and $|e^{-(w\mp T)^2/2s^2}|=e^{-((t\mp T)^2-(\sigma-\frac12)^2)/2s^2}$, a Gaussian in $t$ uniformly for bounded $\sigma$. The value of $g$ is the
Fourier transform of a Gaussian: $\frac1{2\pi}\int e^{-(t-T)^2/2s^2}e^{-iut}\dd t=\frac s{\sqrt{2\pi}}e^{-iuT}e^{-s^2u^2/2}$, plus the same with $-T$.

(b) Integrate $H\frac{\zeta'}\zeta$ around the rectangle $-1\le\sigma\le2$, $|t|\le T'$, with good heights $T'\to\infty$ (Volume~5, Lemma~7.3): by (a) the horizontal sides tend to $0$,
and the residues are $H(\rho)$ at the zeros and $-H(1)$ at the pole of $\zeta$. So
$\sum_\rho H(\rho)-H(1)=\frac1{2\pi i}\int_{(2)}H\frac{\zeta'}\zeta-\frac1{2\pi i}\int_{(-1)}H\frac{\zeta'}\zeta$, both lines upward.

\emph{Right edge.} Expanding $\frac{\zeta'}\zeta$ and moving each term to the line $\frac12$ (the integrand $Hn^{-z}$ is entire and decays),
\[
\frac1{2\pi i}\int_{(2)}H(z)n^{-z}\dd z=\frac1{2\pi}\int h(t)n^{-1/2-it}\dd t=n^{-1/2}g(\log n).
\]
The right edge is $-\sum\Lambda(n)n^{-1/2}g(\log n)$.

\emph{Left edge.} Insert the formula of Problem~\ref{pb:landauzeros}(b). The term $-\frac{\zeta'}\zeta(1-z)$, after $z\mapsto1-z$ and $H(1-z)=H(z)$, gives the right edge
once more. The constant $\log\pi$ gives $-\log\pi\cdot\frac1{2\pi}\int h=-g(0)\log\pi$. The remaining term is $\frac1{2\pi i}\int_{(-1)}H(z)\cdot\frac12\bigl(\frac{\Gamma'}\Gamma(\frac z2)+\frac{\Gamma'}\Gamma(\frac{1-z}2)\bigr)\dd z$.
Moving it to the line $\frac12$ crosses one pole, at $z=0$, where $\frac12\frac{\Gamma'}\Gamma(\frac z2)\sim-\frac1z$; so it equals $H(0)$ plus
$\frac1{2\pi}\int h(t)\operatorname{Re}\frac{\Gamma'}\Gamma(\frac14+\frac{it}2)\dd t$, because $\frac{1-z}2=\overline{\frac z2}$ on that line. Collecting terms gives the formula.
It converges absolutely because $H(\rho)$ decays faster than any power of $|\gamma|$ while there are $O(T\log T)$ zeros up to height $T$.

(c) Under the Riemann hypothesis $H(\rho)=h(\gamma)$, and $\sum_\rho h(\gamma)=2\sum_{\gamma>0}h(\gamma)=2Z(T)+O(e^{-cT})$, since $e^{-(\gamma+T)^2/2s^2}\le e^{-\gamma^2/2s^2}e^{-14T/s^2}$.
Also $|H(0)|=|H(1)|=|h(\frac i2)|\le2e^{(1/4-T^2)/2s^2}$. In the integral, $h$ is concentrated near $\pm T$, where
$\operatorname{Re}\frac{\Gamma'}\Gamma(\frac14+\frac{it}2)=\log\frac{|t|}2+O(t^{-2})$ (Stirling). Expanding $\log\frac t2=\log\frac T2+\frac{t-T}T+O(\frac{(t-T)^2}{T^2})$ against the Gaussian,
the linear term integrates to $0$, so the integral is $\frac{2s}{\sqrt{2\pi}}\log\frac T2+O(\frac1{T^2})$. With $g(0)=\frac{2s}{\sqrt{2\pi}}$ the formula becomes
\[
2Z(T)=\frac{2s}{\sqrt{2\pi}}\log\frac T{2\pi}-2\sum_n\frac{\Lambda(n)}{\sqrt n}g(\log n)+O\Bigl(\frac1T\Bigr),
\]
and dividing by $2$ gives the claim.

(d) A zero $\gamma$ contributes a bump of height $1$ at $T=\gamma$, and the other zeros add little, so the peaks of Figure~\ref{fig:hearzeros} sit at $14.13$, $21.02$, $25.01$, \dots{}
The closest pair below $50$, $48.01$ and $49.77$, is $1.76=2.9s$ apart and still resolved. A smaller $s$ makes narrower bumps and separates closer zeros, but $g(\log n)$ decays only
like $e^{-s^2\log^2n/2}$, so the primes up to about $e^{c/s}$ are needed: halving the width squares the range of primes. This is the uncertainty principle: sharp in $t$ means wide
in $\log n$. A direct computation verifies that the right side agrees with $Z(T)$ computed from the zeros to within $0.003$ on $[5,50]$.

\tools{Problem~\ref{pb:landauzeros}(b); Volume~5, Lemma~7.3 and Theorem~6.1; Fourier transform of a Gaussian (Volume~4).}

\leads Guinand (1948) and Weil (1952) stated the explicit formula for general test functions; Weil showed that the Riemann hypothesis is equivalent to the positivity of a
quadratic form built from it. Zeros are located in practice with the Riemann--Siegel formula, which Siegel found in 1932 in Riemann's unpublished notes; the computation
of (d) is a different route, which uses nothing but the primes.
\end{solution}

\begin{solution}{pb:cramersign}
\idea In Cramér's model, $\pi_{\mathcal C}(x)-\operatorname{li}(x)$ is a sum of independent centred coin tosses, a random walk with no preference of sign. The real primes have a preference,
and its source is visible in (c): $\pi$ counts primes, while the explicit formula naturally counts prime powers with weights $\frac1k$. The squares of primes, which $\pi$ omits,
push $\pi(x)$ below $\operatorname{li}(x)$ by about $\frac{\sqrt x}{\log x}$, which is also the size of the honest fluctuations.

\writeup (a) $\mathbb E\,\pi_{\mathcal C}(x)=\sum_{3\le n\le x}\frac1{\log n}$ differs from $\int_3^x\frac{\dd t}{\log t}$ by at most $\frac1{\log3}$ (a decreasing function), so it is
$\operatorname{li}(x)+O(1)$. $\operatorname{Var}\pi_{\mathcal C}(x)=\sum\frac1{\log n}\bigl(1-\frac1{\log n}\bigr)\sim\frac x{\log x}$. The summands are bounded and the variance tends to infinity, so
Lindeberg's condition holds and $\frac{\pi_{\mathcal C}(x)-\mathbb E\pi_{\mathcal C}(x)}{\sqrt{\operatorname{Var}}}\to\mathcal N(0,1)$ (Volume~3, Theorem~6.10). As
$\frac{\operatorname{li}(x)-\mathbb E\pi_{\mathcal C}(x)}{\sqrt{\operatorname{Var}}}\to0$, $P(\pi_{\mathcal C}(x)>\operatorname{li}(x))\to P(Z>0)=\frac12$.

(b) Let $V_j=\operatorname{Var}\pi_{\mathcal C}(n_j)$ and let $D_j$ be the centred increment of $\pi_{\mathcal C}$ over $]n_{j-1},n_j]$. The $D_j$ are independent, $\operatorname{Var}D_j\sim V_j$ (as
$n_{j-1}=\sqrt{n_j}$), and by the central limit theorem $P(D_j>2\sqrt{V_j})\to P(Z>2)>0$. So $\sum_jP(D_j>2\sqrt{V_j})=\infty$, and by the second Borel--Cantelli lemma (Volume~3,
Theorem~4.11) $D_j>2\sqrt{V_j}$ for infinitely many $j$, almost surely; likewise $D_j<-2\sqrt{V_j}$. The centred value $Y_{j-1}$ at $n_{j-1}$ satisfies
$P(|Y_{j-1}|>\sqrt{V_j})\le\frac{V_{j-1}}{V_j}=O(n_j^{-1/2})$ by Chebyshev, a summable sequence, so by the first lemma $|Y_{j-1}|\le\sqrt{V_j}$ for all large $j$. Then
$\pi_{\mathcal C}(n_j)-\operatorname{li}(n_j)=Y_{j-1}+D_j+O(1)$ is $>\sqrt{V_j}-O(1)>0$ infinitely often and $<0$ infinitely often.

(c) We have
\[
\Pi(x)=\pi(x)+\frac12\pi(\sqrt x)+\sum_{3\le k\le\log_2x}\frac1k\pi(x^{1/k}),
\]
and the last sum is at most $x^{1/3}+x^{1/4}\log_2x=O(x^{1/3})$. Next, $\Pi(x)=\sum_{n\le x}\frac{\Lambda(n)}{\log n}$,
so by Stieltjes integration by parts, $\Pi(x)=\frac{\psi(x)}{\log x}+\int_2^x\frac{\psi(t)}{t\log^2t}\dd t$. In the same way
$\operatorname{li}(x)-\operatorname{li}(2)=\int_2^x\frac{\dd t}{\log t}=\frac x{\log x}-\frac2{\log2}+\int_2^x\frac{t}{t\log^2t}\dd t$. Subtracting gives the identity.

(d) By (c), $\pi(x)-\operatorname{li}(x)=\bigl(\Pi(x)-\operatorname{li}(x)\bigr)-\frac12\pi(\sqrt x)+O(x^{1/3})$, and $\frac12\pi(\sqrt x)\approx\frac{\sqrt x}{\log x}$. The term $\Pi-\operatorname{li}$ is dominated by
$\frac{\psi(x)-x}{\log x}$, which oscillates around $0$ (its $66\,112$ sign changes below $10^8$); the shift by $-\frac{\sqrt x}{\log x}$, about one unit on the scale of the data, is larger
than the oscillation in this range, so $\pi<\operatorname{li}$ throughout. Under the Riemann hypothesis, the explicit formula gives $\psi(x)-x=-\sqrt x\sum_\rho\frac{x^{i\gamma}}\rho+O(\log^2x)$,
so $\frac{\psi(x)-x}{\log x}$ has size $\frac{\sqrt x}{\log x}$ times an oscillating sum whose terms decay like $\frac1{|\gamma|}$. The integral in (c) is smaller: integrating $t^{\rho-1}$ against
$\frac1{\log^2t}$ divides each term by about $|\rho|\log x$, leaving $O(\frac{\sqrt x}{\log^2x})$. So both $\Pi-\operatorname{li}$ and $\pi-\operatorname{li}$ have fluctuations of size $\frac{\sqrt x}{\log x}$,
which is $543$ at $10^8$, while the coin model fluctuates by $\sqrt{x/\log x}=2330$: the real primes are far more regular than independent coins, because their errors are all
governed by the same zeros.

\tools{Volume~3, Theorems~4.11 and~6.10; Volume~5, Definition~10.1 and Theorem~7.5; Stieltjes integration by parts (Volume~5, Lemma~1.1).}

\leads The bias does not last forever: Littlewood proved in 1914 that $\pi(x)-\operatorname{li}(x)$ changes sign infinitely often, because $\sum_\rho\frac{x^{i\gamma}}\rho$ occasionally beats
the shift $-1$. But the first sign change is beyond $10^{19}$ (Büthe, 2018, verified $\pi(x)<\operatorname{li}(x)$ for $2\le x\le10^{19}$). Rubinstein and Sarnak (1994) showed, under the
Riemann hypothesis and a linear independence hypothesis on the $\gamma$, that the set of $x$ with $\pi(x)>\operatorname{li}(x)$ has logarithmic density $0.00000026\ldots$: Chebyshev's bias in
its sharpest form.
\end{solution}

\begin{solution}{pb:selberg}
\idea $\Lambda=\mu*\log$ sees one prime at a time; $\Lambda_2=\mu*\log^2$ sees primes and products of two primes, with total weight $2\log x$ on average. The point is that
$\sum\Lambda_2$ can be computed from Möbius sums alone, with an error $O(x)$, because the $\log^2$ is smooth enough for the error terms to cancel.

\writeup (a) $f\mapsto f\log$ is a derivation of the convolution ring: $(f*g)\log=(f\log)*g+f*(g\log)$, since $\log(dm)=\log d+\log m$. From $\log=1*\Lambda$ we get
$\log^2=\log*\Lambda+1*(\Lambda\log)$; convolving with $\mu$ and using $\mu*\log=\Lambda$, $\mu*1=\delta$ gives $\Lambda_2=\Lambda*\Lambda+\Lambda\log$. Both terms are $\ge0$, and both vanish when
$n$ has three distinct prime factors, since $\Lambda(d)\Lambda(m)\ne0$ forces $d,m$ to be prime powers.

(b) From $\sum_{m\le y}\frac1m=\log y+\gamma+O(\frac1y)$,
\[
\sum_{d\le x}\frac{\mu(d)}d\log\frac xd=\sum_{d\le x}\frac{\mu(d)}d\sum_{m\le x/d}\frac1m-\gamma\sum_{d\le x}\frac{\mu(d)}d+O(1),
\]
which is $\sum_{n\le x}\frac1n\sum_{d\mid n}\mu(d)+O(1)=O(1)$, using
$|\sum_{d\le x}\frac{\mu(d)}d|\le1$ (Volume~5, Exercise~1.6). For the second sum, the hyperbola method (Volume~5, Lemma~2.8) gives
$\sum_{m\le y}\frac{\tau(m)}m=\frac12\log^2y+2\gamma\log y+c+O(y^{-1/2}\log y)$. So $\log^2\frac xd=2\sum_{m\le x/d}\frac{\tau(m)}m-4\gamma\log\frac xd-2c+O\bigl((\frac dx)^{1/2}\log\frac xd\bigr)$. Multiply
by $\frac{\mu(d)}d$ and sum: since $\mu*\tau=1$, the first part is $2\sum_{n\le x}\frac1n=2\log x+O(1)$; the others are $O(1)$ by the first sum and Exercise~1.6, and the error
terms add up to $O\bigl(\sum_{d\le x}\frac1d(\frac dx)^{1/2}\log\frac xd\bigr)=O(1)$.

(c) $\sum_{n\le x}\Lambda_2(n)=\sum_{d\le x}\mu(d)\sum_{m\le x/d}\log^2m$, and $\sum_{m\le y}\log^2m=y(\log^2y-2\log y+2)+O(\log^2y+1)$ by comparison with the integral. So
\[
\sum_{n\le x}\Lambda_2(n)=x\sum_{d\le x}\frac{\mu(d)}d\Bigl(\log^2\frac xd-2\log\frac xd+2\Bigr)+O\Bigl(\sum_{d\le x}\bigl(\log^2\tfrac xd+1\bigr)\Bigr),
\]
which is $x\bigl(2\log x+O(1)\bigr)+O(x)$ by (b), Exercise~1.6, and $\sum_{d\le x}\log^2\frac xd=O(x)$.

(d) By (a),
\[
\sum_{n\le x}\Lambda_2(n)=\sum_{n\le x}\Lambda(n)\log n+\sum_{dm\le x}\Lambda(d)\Lambda(m).
\]
Prime powers $p^k$, $k\ge2$, contribute $O(\sqrt x\log^2x)$ to the first sum and at most
\[
\sum_{p^k,\,k\ge2}\log p\,\psi\Bigl(\frac x{p^k}\Bigr)=O\Bigl(x\sum_{p,\,k\ge2}\frac{\log p}{p^k}\Bigr)=O(x)
\]
to the second. And $\sum_{p\le x}\log^2p=\theta(x)\log x-\int_2^x\frac{\theta(t)}t\dd t=\theta(x)\log x+O(x)$.

(e) Write the formula as $\theta(x)\log x+\sum_{p\le x}\log p\,\theta(\frac xp)=2x\log x+O(x)$; by Chebyshev, $A<\infty$. Let $\varepsilon>0$; there is $X$ with $\theta(y)\le(A+\varepsilon)y$ for $y\ge X$,
and $\theta(\frac xp)=O_\varepsilon(1)$ for $p>\frac xX$. By Mertens (Volume~5, Theorem~1.3), $\sum_{p\le x}\log p\,\theta(\frac xp)\le(A+\varepsilon)x\log x+O_\varepsilon(x)$. Along a sequence with
$\theta(x)\le(a+\varepsilon)x$, the formula gives $2\le a+A+2\varepsilon+o(1)$. Symmetrically, $\theta(y)\ge(a-\varepsilon)y$ for $y\ge X$ gives
$\sum_p\log p\,\theta(\frac xp)\ge(a-\varepsilon)x\sum_{p\le x/X}\frac{\log p}p=(a-\varepsilon)x\log x-O_\varepsilon(x)$, and along a sequence with $\theta(x)\ge(A-\varepsilon)x$, $2\ge A+a-2\varepsilon-o(1)$.
So $a+A=2$.

(f) An error term $O(x)$ in a main term $2x\log x$ is a relative error $O(\frac1{\log x})$. The data say the error is $-3.67x+o(x)$, so the ratio is $1-\frac{3.67}{2\log x}$, which is
$0.886$ at $10^7$ and would reach $0.99$ only near $x=10^{80}$. There is no contradiction: the formula is precise in absolute terms, and $\log x$ grows slowly.

\tools{Volume~5, Exercise~1.6, Lemma~2.8, Proposition~1.2 and Theorem~1.3; the derivation $f\mapsto f\log$ (Volume~1, Dirichlet convolution).}

\leads Selberg's formula is the first step of the elementary proof. If $\theta(x)=x+R(x)$, it says roughly that $R(x)\log x\approx-\sum_p\log p\,R(\frac xp)$: a large error at $x$ forces
errors of the opposite sign at $\frac xp$. Selberg and Erd\H os turned this averaging into $R(x)=o(x)$, Erd\H os by showing that every interval $]x,(1+\varepsilon)x]$ has many primes.
Part (e) is the identity Selberg found in April 1948, before the full proof.
\end{solution}

\begin{solution}{pb:schur}
\idea The Fourier matrix is almost an involution: $F^2$ is $p$ times a reflection, so its eigenvalues are $\pm\sqrt p$ and $\pm i\sqrt p$ and its trace, the Gauss sum, is determined by
four multiplicities. Three equations come from $\operatorname{tr}F$, $\operatorname{tr}F^2$ and $|G|^2=p$; the fourth comes from the determinant, which Vandermonde computes exactly,
sign included. That fixes the sign of $G$, which Gauss needed four years to find.

\writeup (a) $\sum_{n=0}^{p-1}\zeta^{n^2}=1+\sum_{m\ne0}(1+\chi(m))\zeta^m=\sum_{m\ne0}\chi(m)\zeta^m$, as $\sum_{m\ne0}\zeta^m=-1$. By Volume~5, Theorem~3.5, $|G|^2=p$, and
$\overline G=\sum\chi(m)\zeta^{-m}=\chi(-1)G$, so $G^2=\chi(-1)G\overline G=\chi(-1)p$. And $\operatorname{tr}F=\sum_j\zeta^{j^2}=G$.

(b) $(F^2)_{jk}=\sum_l\zeta^{l(j+k)}=p\,[j+k\equiv0]$, so $F^2=pR$, $R^2=I$ and $F^4=p^2I$. So $F$ is diagonalisable with eigenvalues among the four fourth roots of $p^2$. Then
$p=a+b+c+d$; $\operatorname{tr}F^2=p\operatorname{tr}R=p$ (only $j=0$ is fixed by $j\mapsto-j$) gives $p(a+b)-p(c+d)=p$; and $\operatorname{tr}F=\sqrt p(a-b+ic-id)=G$. Finally
$|G|^2=p$ gives $(a-b)^2+(c-d)^2=1$.

(c) $F=(x_k^j)$ with $x_k=\zeta^k$, so $\det F=\prod_{j<k}(\zeta^k-\zeta^j)$, and $\zeta^k-\zeta^j=e(\frac{j+k}{2p})\cdot2i\sin\frac{\pi(k-j)}p$, with $\sin\frac{\pi(k-j)}p>0$. The
product of the phases is $e\bigl(\frac1{2p}\sum_{j<k}(j+k)\bigr)$, and $\sum_{j<k}(j+k)=(p-1)\sum_jj=\frac{p(p-1)^2}2$, so the phase is $e(\frac{(p-1)^2}4)=1$. Hence
$\det F=i^{p(p-1)/2}\Delta$ with $\Delta=2^{p(p-1)/2}\prod_{j<k}\sin\frac{\pi(k-j)}p>0$.

(d) The product of the eigenvalues is
\[
\sqrt p^{\,p}(-1)^bi^c(-i)^d=\sqrt p^{\,p}(-1)^bi^{c-d},
\]
so $(-1)^bi^{c-d}=i^{p(p-1)/2}$. From (b), $a+b=\frac{p+1}2$ and $c+d=\frac{p-1}2$.
If $p=4m+1$, then $a+b=2m+1$ is odd, so $a\ne b$; hence $|a-b|=1$, $c=d$, and $b\in\{m,m+1\}$. As $i^{p(p-1)/2}=i^{2m(4m+1)}=(-1)^m$, $b=m$, $a-b=1$ and $G=\sqrt p$. If $p=4m+3$,
then $c+d=2m+1$ is odd, so $|c-d|=1$, $a=b=m+1$, and $(-1)^{m+1}i^{c-d}=i^{(4m+3)(2m+1)}=i^{2m+3}=(-1)^{m+1}i$; so $c-d=1$ and $G=i\sqrt p$.

(e) Since $\chi(-1)=-1$, pairing $a$ with $p-a$ gives $T=\sum_{a<p/2}\bigl(a\chi(a)-(p-a)\chi(a)\bigr)=2U-pS$. Splitting by parity, the even $2a$ give $2\chi(2)U$ and the odd
$p-2a$ give $\sum_{a<p/2}(p-2a)\chi(p-2a)=-\chi(2)(pS-2U)$, so $T=\chi(2)(4U-pS)$. Eliminating $U$: $T=\chi(2)(2T+pS)$, so $(\chi(2)-2)T=pS$. Dirichlet's formula with
$\tau(\chi)=G=i\sqrt p$ reads $L(1,\chi)=-\frac{i\pi}{p\cdot i\sqrt p}T=-\frac{\pi T}{p^{3/2}}$. Hence $S=\frac{(2-\chi(2))\sqrt p}\pi L(1,\chi)$. Since $L(1,\chi)>0$ for a real character
(Volume~5, Theorem~3.8 and Exercise~3.4), $S>0$; as $S$ is an integer, $S\ge1$ and $L(1,\chi)\ge\frac\pi{(2-\chi(2))\sqrt p}\ge\frac\pi{3\sqrt p}$.

(f) $\chi(2)=1$ for $p\equiv7\pmod8$ ($7$, $23$, $10\,007$) and $\chi(2)=-1$ for $p\equiv3\pmod8$ ($163$, $1019$). For instance $\frac{\sqrt7}\pi1.18741=1.0000$, $3\frac{\sqrt{163}}\pi0.246069=3.0000$ and
$\frac{\sqrt{10007}}\pi2.41818=77.00$. For $p=163$, $S=3$ and $2-\chi(2)=3$, so $L(1,\chi)=\frac\pi{\sqrt{163}}$ exactly.

\tools{Volume~5, Theorems~3.5, 3.8 and~3.10, Exercise~3.4; Vandermonde determinants (Volume~2); eigenvalues and traces (Volume~2).}

\leads Gauss conjectured the sign in 1801 and proved it in 1805, after four years of trying; Schur's proof (1921) is the one given here. Dirichlet's class number formula says that
$h(-p)=\frac{\sqrt p}\pi L(1,\chi)$ for $p\equiv3\pmod4$, $p>3$, so (e) reads $h(-p)=\frac S{2-\chi(2)}$: the class number counts how many more residues than non-residues lie in the first
half. The values above are $h=1,3,1,13,77$; $h(-163)=1$ is the reason for Euler's polynomial (Problem~\ref{pb:euler41}). No proof of $S>0$ is known that avoids analysis.
\end{solution}

\begin{solution}{pb:euler41}
\idea A prime factor $p<q$ of $n^2+n+q$ is a value of a quadratic form of discriminant $1-4q$ which cannot be equivalent to $x^2+xy+qy^2$, since the latter does not represent small
numbers. So composite values and extra classes of forms are the same thing. For the density of primes, each prime $p$ divides $f(n)$ for $1+\chi(p)$ residues $n$, and $f$ is lucky:
$\chi(p)=-1$ for every $p<41$.

\writeup (a) $x\mapsto x+my$ turns $[a,b,c]$ into $[a,b+2ma,c']$, and $(x,y)\mapsto(-y,x)$ turns it into $[c,-b,a]$. Choose $m$ with $-a<b+2ma\le a$; if then $a>c$, swap. Each swap
lowers $a$, a positive integer, so the process stops at a form with $|b|\le a\le c$; if $b=-a$ translate once more, and if $a=c$ and $b<0$ swap. A reduced form has
$|D|=4ac-b^2\ge4a^2-a^2=3a^2$. A unimodular substitution maps coprime pairs bijectively to coprime pairs, which gives the last claim.

(b) Let $[a,b,c]$ be reduced of discriminant $D=1-4q$; $b$ is odd. Replacing $n$ by $-1-n$ does not change $f(n)$, so write $|b|=2n+1$ with $n\ge0$. Then
$ac=\frac{b^2-D}4=n^2+n+q=f(n)$, and $n<\frac a2\le\frac12\sqrt{(4q-1)/3}<\sqrt{q/3}\le q-2$ for $q\ge3$ (for $q=2$, $a=1$ is forced). If all $f(n)$, $0\le n\le q-2$, are prime, then
$a=1$, $|b|=1$, and $b=1$ by the convention: $[1,1,q]$ is the only reduced form. Conversely, suppose $f(n)$ is composite for some $0\le n\le q-2$. Since
$f(n)\le q^2-2q+2<q^2$, its least prime factor $p$ is $<q$. The form $[p,2n+1,\frac{f(n)}p]$ has discriminant $D$ and represents $p$ with $(x,y)=(1,0)$. But
$x^2+xy+qy^2=(x+\frac y2)^2+(q-\frac14)y^2$ is $\ge q$ when $y\ne0$, and $x^2$ with $\gcd(x,0)=1$ is $1$. So $[1,1,q]$ does not represent $p$ with coprime $(x,y)$, and by (a) the reduced form
equivalent to $[p,2n+1,\frac{f(n)}p]$ is another one.

(c) For $D=-163$: $a\le\sqrt{163/3}<8$, $b$ odd, $ac=\frac{b^2+163}4$, which is $41,43,47,53$ for $|b|=1,3,5,7$, all prime. So $a=1$, $|b|\le1$, and $[1,1,41]$ is the only reduced
form: $n^2+n+41$ is prime for $0\le n\le39$. The list $2,3,5,11,17,41$ is complete. If $1-4q=f^2D_0$ with $f>1$, then $f$ times the principal form of discriminant $D_0$ is a second
reduced form; so $D=1-4q$ must be a fundamental discriminant with one class of forms, and Heegner (1952), Baker (1966) and Stark (1967) proved that the only such $D<-4$ are
$-7,-8,-11,-19,-43,-67,-163$.

(d) For odd $p$, $4f(n)=(2n+1)^2+163$, and $2n+1$ runs over all residues modulo $p$ with $n$, so the number of roots is the number of $y$ with $y^2\equiv-163$, that is $1+\chi(p)$. And
$f(n)=n(n+1)+41$ is odd. For a random integer near $f(n)$, the chance of being prime is $\frac1{\log f(n)}$, and the chance of avoiding the prime factor $p$ is $1-\frac1p$; for $f(n)$ it is
$1-\frac{w(p)}p$ with $w(p)$ the number of roots. Correcting prime by prime gives $C=\prod_p\frac{1-w(p)/p}{1-1/p}$: the factor at $2$ is $\frac1{1/2}=2$, and at odd $p$ it is
$\frac{p-1-\chi(p)}{p-1}$. For the faster product, $\prod_{p>2}(1-\frac{\chi(p)}p)=\frac{1}{L(1,\chi)}\cdot\frac1{1-\chi(2)/2}=\frac{\sqrt{163}}\pi\cdot\frac23$, since $\chi(2)=-1$ ($-163\equiv5\pmod8$) and
$L(1,\chi)=\frac\pi{\sqrt{163}}$ (Problem~\ref{pb:schur}(f); $\chi$ is the Legendre symbol modulo $163$, by reciprocity). The Euler product of $L(1,\chi)$ converges, in the order of the primes,
to $L(1,\chi)$, by the prime number theorem in progressions. The quotients $\frac{1-\chi(p)/(p-1)}{1-\chi(p)/p}=1+O(p^{-2})$ make the remaining product absolutely convergent.

(e) For $n^2+1$, the factor at $2$ is $\frac{1-1/2}{1-1/2}=1$: half the values are even. For $f(n)=n^2+n+41$, the factor at $2$ is $2$, and $\chi(p)=-1$ for every prime $p\le37$ (otherwise
$p$ would divide some $f(n_0)$ with $0\le n_0<p$, a prime larger than $p$), so each such $p$ contributes $1+\frac1{p-1}$ instead of $1\pm\frac1{p-1}$. The primes up to $37$ alone give
a ratio $5.07$, and the larger primes, where $\chi$ and $\chi_4$ take both signs, bring it to $\frac{6.6395}{1.3728}=4.84$; the counts give $\frac{261\,081}{54\,110}=4.83$.

\tools{Binary quadratic forms (Volume~1); Volume~5, Theorem~3.10 and Theorem~5.8; Problem~\ref{pb:schur}.}

\leads Rabinowitsch proved (b) in 1913. Hardy and Littlewood stated their conjecture in 1923; it is open for every quadratic polynomial, even for $n^2+1$, which is one of
Landau's four problems of 1912. Jacobson and Williams (2003) searched for quadratic polynomials whose constant $C$ beats Euler's.
\end{solution}

\begin{solution}{pb:hecke}
\idea Dirichlet sorted primes by residue classes with characters of $(\Z/q\Z)^*$; Hecke sorted Gaussian primes by angle with the characters $e^{4ik\theta}$ of the circle. The
proof of the prime number theorem goes through unchanged: the $L$-functions continue a little to the left of $\operatorname{Re}s=1$, do not vanish on the line by the $3$-$4$-$1$ trick,
and Newman's theorem finishes. Weyl's criterion turns ``$\vartheta_k=o(x)$ for all $k\ne0$'' into equidistribution.

\writeup (a) A unit is $i^j$ and $(i^j)^{4k}=1$, so $\lambda_k$ depends only on $(\alpha)$, and it is multiplicative because $\frac{\alpha\beta}{|\alpha\beta|}=\frac\alpha{|\alpha|}\frac\beta{|\beta|}$. Let $Q_\alpha=\alpha+[0,1[^2$.
For $|\alpha|\ge2$ and $z\in Q_\alpha$, the arguments of $z$ and $\alpha$ differ by $O(\frac1{|\alpha|})$, so $|\lambda_k(\alpha)-(\frac z{|z|})^{4k}|=O(\frac{|k|}{|\alpha|})$. Hence
\[
S_k(x)=\int_{\bigcup Q_\alpha}\Bigl(\frac z{|z|}\Bigr)^{4k}\dd A+O\Bigl(\sum_{0<|\alpha|\le\sqrt x}\frac{|k|}{|\alpha|}\Bigr)+O(1).
\]
The union of the squares differs from the disc of radius $\sqrt x$ by a set of area $O(\sqrt x)$, the integral over the disc is $\int_0^{\sqrt x}\int_0^{2\pi}e^{4ik\theta}r\dd\theta\dd r=0$,
and $\sum_{|\alpha|\le\sqrt x}\frac1{|\alpha|}=O(\sqrt x)$ by comparison with $\int\frac{\dd A}{|z|}$.

(b) $L_k(s)=\sum_na_k(n)n^{-s}$ with $\sum_{n\le x}a_k(n)=\frac14S_k(x)=O(\sqrt x)$, so the series converges and is holomorphic on $\operatorname{Re}s>\frac12$ (partial summation, Volume~5,
Lemma~1.1). Since $\Z[i]$ has unique factorisation (Volume~1, Theorem~10.3), $L_k(s)=\sum_I\lambda_k(I)N(I)^{-s}$ over non-zero ideals is an Euler product over prime ideals. These are
(Volume~1, Theorem~10.4) $(1+i)$ with $\lambda_k=(\frac{1+i}{\sqrt2})^{4k}=e^{i\pi k}=(-1)^k$; $(a+bi)$ and $(a-bi)$ for $p\equiv1\pmod4$, with $\lambda_k=e^{\pm4ik\varphi_p}$; and $(p)$, of norm
$p^2$ and $\lambda_k=1$, for $p\equiv3\pmod4$. For $k=0$, $L_0(s)=\frac14\sum r_2(n)n^{-s}=\sum(1*\chi_4)(n)n^{-s}=\zeta(s)L(s,\chi_4)$ (Volume~1, Theorem~10.5).

(c) For $\sigma>1$,
\[
\log\bigl|L_0(\sigma)^3L_k(\sigma+it)^4L_{2k}(\sigma+2it)\bigr|=\sum_{\mathfrak p,m}\frac{N\mathfrak p^{-m\sigma}}m(3+4\cos\theta+\cos2\theta)\ge0,
\]
with
$\theta=m(\arg\lambda_k(\mathfrak p)-t\log N\mathfrak p)$, because $\lambda_{2k}=\lambda_k^2$. If $L_k(1+it)=0$ with $k\ne0$, then as $\sigma\to1^+$, $L_0(\sigma)^3=O((\sigma-1)^{-3})$,
$L_k(\sigma+it)^4=O((\sigma-1)^4)$ and $L_{2k}(\sigma+2it)$ stays bounded (it is holomorphic there, as $2k\ne0$): the product tends to $0$, a contradiction.

(d) In $-\frac{L_k'}{L_k}(s)=\sum_{\mathfrak p,m}\lambda_k(\mathfrak p)^m\log N\mathfrak p\,N\mathfrak p^{-ms}$, the terms with $m=1$ and $N\mathfrak p=p\equiv1\pmod4$ give $\sum_p2\cos(4k\varphi_p)\log p\,p^{-s}$,
and the other terms form a series absolutely convergent on $\operatorname{Re}s>\frac12$. So $\Phi_k(s)=\int_1^\infty x^{-s}\dd\vartheta_k(x)$ is $-\frac{L_k'}{L_k}(s)$ plus a function holomorphic on
$\operatorname{Re}s>\frac12$. For $k\ne0$, $\Phi_k$ is holomorphic on $\operatorname{Re}s\ge1$ by (b) and (c); and $\Phi_0(s)-\frac1{s-1}$ is holomorphic there, since $\zeta(s)L(s,\chi_4)$ has a
simple pole at $1$ and no zero on the line (Volume~5, Theorems~3.8 and~5.2, and the same argument for $L(s,\chi_4)$). The functions $A_\pm=\vartheta_0\pm\vartheta_k$ are
non-decreasing (their increments are $2(1\pm\cos)\log p\ge0$) and $O(x)$, and $\int_1^\infty(A_\pm(x)-x)x^{-s-1}\dd x=\frac{\Phi_0(s)\pm\Phi_k(s)}s-\frac1{s-1}$ is holomorphic on
$\operatorname{Re}s\ge1$. By Newman's theorem $\int_1^\infty\frac{A_\pm(x)-x}{x^2}\dd x$ converges, so $A_\pm(x)\sim x$ by Lemma~5.4, and $\vartheta_k=\frac12(A_+-A_-)=o(x)$.

(e) By partial summation,
\[
\sum_{p\le x,\ p\equiv1\,(4)}\cos(k\psi_p)=o(\pi(x))\qquad(k\ne0),
\]
where $\psi_p=4\varphi_p\in\,]0,\pi[$. List the primes $p\equiv1\pmod4$ in order, and form the
sequence $\frac{\psi_{p}}{2\pi},-\frac{\psi_{p}}{2\pi},\dots$ modulo $1$, each prime giving two terms. Its Weyl sums are
\[
\frac1{2N}\sum\bigl(e^{ik\psi_p}+e^{-ik\psi_p}\bigr)=\frac1N\sum\cos(k\psi_p)\to0,
\]
so it is equidistributed. The positive copies lie in $]0,\frac12[$ and the negative ones in $]\frac12,1[$, so for $[\alpha,\beta]\subset\,]0,\pi[$ the proportion of $p$ with
$\psi_p\in[\alpha,\beta]$ tends to $2\cdot\frac{\beta-\alpha}{2\pi}=\frac{\beta-\alpha}\pi$. So $\psi_p$ is uniform on $]0,\pi[$, and $\varphi_p$ on $]0,\frac\pi4[$.

(f) $b<\frac a{10}$ means $\varphi_p<\arctan\frac1{10}$, which has predicted proportion $\frac4\pi\arctan\frac1{10}=0.12690$, against $0.12658$ observed.

\tools{Volume~1, Theorems~10.3--10.5; Volume~5, Lemma~1.1, Theorems~3.8, 5.2, 5.3 and Lemma~5.4; Weyl's criterion (Volume~4, Theorem~6.13).}

\leads Hecke introduced these characters in 1918--1920 for every imaginary quadratic field, and in general ``Grössencharaktere'' for all number fields. Fixing $b$ instead of the angle is
much harder: Friedlander and Iwaniec proved in 1998 that there are infinitely many primes $a^2+b^4$, and Heath-Brown (2001) did the same for $a^3+2b^3$. Primes $a^2+1$, that
is $b=1$, are still out of reach (Problem~\ref{pb:euler41}).
\end{solution}

\begin{solution}{pb:shifted}
\idea Each small prime $q$ divides $p-1$ with probability $\frac1{q-1}$, independently, as the prime number theorem in progressions says; so $\omega(p-1)$ behaves like a sum of independent
coins with $\sum\frac1{q-1}\approx\log\log x$, the Erd\H os--Kac picture for shifted primes. Upper bounds for large $q$ come cheaply from Brun--Titchmarsh; matching lower bounds
need Bombieri--Vinogradov.

\writeup (a) Let $Q=\prod_{q\le y}q$. For $p\nmid Q$, the set of $q\le y$ dividing $p-1$ is determined by $p$ modulo $Q$, and by the Chinese remainder theorem the classes modulo $Q$ are
tuples of classes modulo the $q$. Each class prime to $Q$ has proportion $\frac1{\varphi(Q)}$ of the primes (Volume~5, Theorem~5.8). So for $S\subset\{q\le y\}$ the proportion of primes with
$q\mid p-1$ exactly for $q\in S$ tends to $\prod_{q\in S}\frac1{q-1}\prod_{q\notin S}\frac{q-2}{q-1}$, the probability that $B_q=1$ exactly for $q\in S$.

(b) $\frac1{q-1}=\frac1q+\frac1{q(q-1)}$, so by Mertens (Volume~5, Theorem~1.6) $\sum_{q\le y}\frac1{q-1}=\log\log y+M+\sum_q\frac1{q(q-1)}+o(1)$, and
$0.261497+0.773157=1.034654$. The variance $\sum_{q\le y}\frac1{q-1}(1-\frac1{q-1})$ is $\log\log y+O(1)\to\infty$ and the $B_q$ are bounded, so Lindeberg's theorem (Volume~3,
Theorem~6.10) applies, and the difference between the mean and $\log\log y$ is $O(1)=o(\sqrt{\log\log y})$.

(c) $\sum_{p\le x}\omega(p-1)=\sum_q\pi(x;q,1)$. At most one $q>\sqrt x$ divides $p-1$, so these $q$ contribute at most $\pi(x)$. Fix $\varepsilon>0$. For $q\le x^\varepsilon$,
Brun--Titchmarsh gives $\pi(x;q,1)\le\frac{(2+o(1))x}{(q-1)(1-\varepsilon)\log x}$, uniformly; summing with (b), these $q$ give at most $\frac{2+o(1)}{1-\varepsilon}\pi(x)\log\log x$. For
$x^\varepsilon<q\le\sqrt x$, $\pi(x;q,1)\le\frac{(4+o(1))x}{(q-1)\log x}$ and $\sum_{x^\varepsilon<q\le\sqrt x}\frac1{q-1}=\log\frac1{2\varepsilon}+o(1)$, so they give $O_\varepsilon(\pi(x))$. Let $\varepsilon\to0$.

(d) With $Q=\frac{\sqrt x}{\log^Bx}$, Bombieri--Vinogradov gives $\sum_{q\le Q}\bigl|\pi(x;q,1)-\frac{\operatorname{li}(x)}{q-1}\bigr|=O(\frac x{\log^2x})$, so the $q\le Q$ contribute
$\operatorname{li}(x)\sum_{q\le Q}\frac1{q-1}+O(\frac x{\log^2x})=\pi(x)(\log\log x+O(1))$. The $q\in\,]Q,\sqrt x]$ contribute $O(\pi(x)\sum_{Q<q\le\sqrt x}\frac1q)=o(\pi(x))$ by Brun--Titchmarsh, and the
$q>\sqrt x$ at most $\pi(x)$.

(e) The model says that $q\mid p-1$ with probability $\frac1{q-1}$ for every $q<x$. For $q\le\sqrt x$ this is (d) on average, but for $q>\sqrt x$ it is a statement about primes $p$
with $\frac{p-1}q$ small and $q$ prime, a twin-prime-like problem; so the constant $1.035$ is a conjecture, which the data support: the deficits $0.228$, $0.180$, $0.143$, $0.119$,
$0.102$ shrink roughly like $\frac2{\log x}$. The variance is another matter. In the model the indicators of large $q$ are independent, so two or three large prime factors can
occur together; in reality at most one prime factor of $p-1$ exceeds $\sqrt x$, at most two exceed $x^{1/3}$, and so on. These constraints make the large-$q$ indicators negatively
correlated and remove a constant from the variance, which is of size $\log\log x+O(1)$ with a large negative $O(1)$. The same happens for $\omega(n)$: variance $1.21$ at $10^8$,
against $\log\log10^8=2.91$. A constant is a lot when $\log\log x$ is $2.9$.

\tools{Volume~5, Theorems~1.6 and~5.8, Exercise~8.4, the remark after Proposition~8.11; Lindeberg (Volume~3, Theorem~6.10).}

\leads Erd\H os proved in 1935, with Brun's sieve, that $\omega(p-1)$ has normal order $\log\log p$. The central limit theorem for $\omega(p-1)$, the analogue of Erd\H os--Kac
(Volume~5, Theorem~9.7), was proved later along the same lines, with sieve bounds for primes in progressions replacing the trivial count of multiples. The same method handles $\omega(p+a)$ for any fixed $a\ne0$.
\end{solution}

\begin{solution}{pb:romanoff}
\idea There are about $\frac x{\log x}\cdot\log_2x\asymp x$ pairs $(p,k)$ with $p+2^k\le x$, so $\mathcal R$ has positive density unless the sums collide too often. Collisions are pairs of
primes differing by $2^k-2^l$, which a sieve bounds by the twin prime count times $(\frac d{\varphi(d)})^2$. The whole difficulty is that $2^m-1$ might have unusually many small prime
factors for many $m$; it does not, because a prime $q$ divides $2^m-1$ only when $\operatorname{ord}_q(2)\mid m$, and orders are rarely small.

\writeup (a) For $p\le\frac x2$ and $2^k\le\frac x2$, $p+2^k\le x$: at least $\pi(\frac x2)\lfloor\log_2\frac x2\rfloor\ge c\,x$ pairs. By Cauchy--Schwarz,
$(\sum_{n\le x}r(n))^2\le\#\{n\le x:r(n)>0\}\cdot\sum_{n\le x}r(n)^2$.

(b) $\sum_{n\le x}r(n)^2$ counts quadruples with $p+2^k=p'+2^l\le x$. The diagonal $k=l$ gives $\sum r(n)\le\pi(x)\log_2x=O(x)$. For $k>l$, $p'=p+2^l(2^m-1)$ with $m=k-l$, and
$d=2^l(2^m-1)$ has $\frac d{\varphi(d)}=2\frac{2^m-1}{\varphi(2^m-1)}$, so the number of $p$ is at most $4c\,f(m)\frac x{\log^2x}$. Summing over $l\le\log_2x$ and $m\le\log_2x$, and
doubling for $k<l$, gives $O\bigl(\frac x{\log x}\sum_{m\le\log_2x}f(m)\bigr)=O(x)$ if $\sum_{m\le L}f(m)=O(L)$.

(c) $\frac n{\varphi(n)}=\prod_{q\mid n}\frac q{q-1}$, so $(\frac n{\varphi(n)})^2=\prod_{q\mid n}(1+g(q))=\sum_{e\mid n}\mu(e)^2g(e)$ with $1+g(q)=\frac{q^2}{(q-1)^2}$, that is
$g(q)=\frac{2q-1}{(q-1)^2}$; and $g(q)\le\frac5q$ is $3q^2-9q+5\ge0$, true for $q\ge3$. As $2^m-1$ is odd, and $e\mid2^m-1$ if and only if $\operatorname{ord}_e(2)\mid m$,
\[
\sum_{m\le L}f(m)=\sum_{e\text{ odd}}\mu(e)^2g(e)\,\#\{m\le L:\operatorname{ord}_e(2)\mid m\}\le L\sum_{e\text{ odd}}\frac{\mu(e)^2g(e)}{\operatorname{ord}_e(2)}.
\]

(d) If $\operatorname{ord}_e(2)=j\le J$, then $e\mid2^j-1\mid P_J$, and $P_J<\prod_{j\le J}2^j\le2^{J^2}$. So $A(J)=\sum_{\operatorname{ord}_e(2)\le J}\mu(e)^2g(e)\le\prod_{q\mid P_J}(1+\frac5q)$. $P_J$ has
fewer than $J^2$ prime factors, so the product is at most $\prod_{q\le p_{J^2}}(1+\frac5q)\le\exp(5\sum_{q\le p_{J^2}}\frac1q)=O((\log J)^5)$ by Mertens, as $p_{J^2}=O(J^3)$. By partial summation,
$\sum_e\frac{\mu(e)^2g(e)}{\operatorname{ord}_e(2)}=\sum_j\frac{A(j)-A(j-1)}j=\sum_j\frac{A(j)}{j(j+1)}<\infty$. With (a)--(c), $\#(\mathcal R\cap[1,x])\ge\frac{c^2x^2}{O(x)}$.

(e) The product is $2\cdot3\cdot5\cdot7\cdot13\cdot17\cdot241=11\,184\,810$, and $n\equiv7\,629\,217$ is odd. The orders of $2$ modulo $3,5,7,13,17,241$ are $2,4,3,12,8,24$, and $n\equiv1,2,1,11,8,121$ modulo these primes. So $3\mid n-2^k$ when $k\equiv0\pmod2$,
$5\mid n-2^k$ when $k\equiv1\pmod4$, $7\mid n-2^k$ when $k\equiv0\pmod3$, $13\mid n-2^k$ when $k\equiv7\pmod{12}$, $17\mid n-2^k$ when $k\equiv3\pmod8$, and $241\mid n-2^k$ when
$k\equiv23\pmod{24}$. These congruences cover all $k$: even $k$ by the first; $k\equiv1\pmod4$ by the second; $k\equiv3\pmod8$ by the fifth; and $k\equiv7\pmod8$ is $7$, $15$ or $23$ modulo
$24$, covered by the fourth, the third and the sixth. A direct computation verifies all $24$ residues of $k$. So $n-2^k$ is
divisible by one of the six primes, and is not prime unless it equals that prime; for each such prime this happens for at most one $k$ per $n$, and only for a sparse set of $n$
in the progression. So a positive proportion, $\frac1{11\,184\,810}$ up to sparse exceptions, of all integers are odd and not in $\mathcal R$.

(f) $n=p+2^k$ with $k\ge1$ is odd unless $p=2$, so $\mathcal R$ contains almost no even numbers, and its density is at most $\frac12$; the data are half the odd proportion. An odd $n$
has on average about $\sum_{k\le\log_2n}\frac{2}{\log n}\approx\frac2{\log2}=2.9$ representations, adjusted by local factors from $3$, $5$, $7$, \dots; with so few expected
representations, a noticeable fraction of odd numbers get none, as the $8\%$ in the data show. The covering progression of (e) is far too thin to explain them. The values of
$\frac1L\sum f(m)$ level off near $2.5$, as (c)--(d) predict.

\tools{Volume~5, Theorem~8.2 and Exercise~8.5 (Selberg's sieve), Theorem~1.3 (Mertens), Proposition~1.2; orders modulo $e$ (Volume~1).}

\leads Romanoff proved the theorem in 1934. Erd\H os found the covering of (e) in 1950, and asked whether covering systems with all moduli distinct can have arbitrarily large
least modulus; Hough proved in 2015 that they cannot. The exact density of $\mathcal R$ is unknown.
\end{solution}

\section*{Sources and further reading}

\begin{itemize}
\item S.~Ramanujan, \emph{On certain arithmetical functions}, Transactions of the Cambridge Philosophical Society 22 (1916), 159--184.
\item L.~J.~Mordell, \emph{On Mr Ramanujan's empirical expansions of modular functions}, Proceedings of the Cambridge Philosophical Society 19 (1917), 117--124.
\item E.~Hecke, \emph{Über Modulfunktionen und die Dirichletschen Reihen mit Eulerscher Produktentwicklung I, II}, Mathematische Annalen 114 (1937), 1--28 and 316--351.
\item C.~L.~Siegel, \emph{A simple proof of $\eta(-1/\tau)=\eta(\tau)\sqrt{\tau/i}$}, Mathematika 1 (1954), 4.
\item P.~Deligne, \emph{La conjecture de Weil I}, Publications Mathématiques de l'IHÉS 43 (1974), 273--307.
\item J.-P.~Serre, \emph{A Course in Arithmetic}, Springer GTM 7, 1973, Chapter~VII. The classic short introduction.
\item T.~M.~Apostol, \emph{Modular Functions and Dirichlet Series in Number Theory}, 2nd ed., Springer GTM 41, 1990, Chapters~1--3 and~6. Siegel's proof and Hecke operators.
\item D.~Zagier, \emph{Elliptic modular forms and their applications}, in \emph{The 1-2-3 of Modular Forms}, Springer, 2008, 1--103. A delightful survey.
\item F.~Diamond and J.~Shurman, \emph{A First Course in Modular Forms}, Springer GTM 228, 2005.
\item T.~Barnet-Lamb, D.~Geraghty, M.~Harris and R.~Taylor, \emph{A family of Calabi--Yau varieties and potential automorphy II}, Publications of the RIMS 47 (2011), 29--98. The Sato--Tate conjecture.
\end{itemize}
